% Matrices et applications linéaires d'un plan vectoriel — Mathématiques 1ère C — Cours signé OnBuch+
% Programme officiel MINESEC. Compiler avec : tectonic N17-matrices-applications-lineaires.tex
\newcommand{\DOCMATIERE}{Mathématiques}
\newcommand{\DOCNIVEAU}{1ère C}
\newcommand{\DOCMODULE}{Configurations et transformations élémentaires du plan}
\newcommand{\DOCLECON}{N17}
\newcommand{\DOCTITRE}{Matrices et applications linéaires d'un plan vectoriel}
% OnBuch+ — préambule « cours signés » ENRICHI (compilé avec tectonic / XeLaTeX)
% Système visuel « Pop » d'OnBuch+ : fond crème (jamais blanc), bordures encre
% épaisses, ombres « dures » décalées, titres Archivo Black en capitales.
% Version MATHÉMATIQUES : graphes (pgfplots/tikz), boîtes pédagogiques
% (définition, propriété, méthode, attention, le savais-tu, exercice résolu,
% exercice, TP, bilan), page de garde et signature OnBuch+ ancrée en bas.
%
% Macros attendues dans main.tex AVANT \input{preamble} :
%   \DOCMATIERE  \DOCNIVEAU  \DOCMODULE  \DOCLECON (numéro)  \DOCTITRE
\documentclass[11pt]{article}

\usepackage{fontspec}
\usepackage[french]{babel}
\usepackage[a4paper,margin=2cm,top=3.4cm,bottom=2.8cm,headheight=1.4cm,footskip=1.3cm]{geometry}
\usepackage{amsmath,amssymb,amsfonts}
\usepackage{mathtools} % \xmapsto, \coloneqq... souvent utilisés par les modèles
\usepackage{cancel} % simplifications barrées
% \abs{z} (module, valeur absolue) et \norm{u} (norme), utilisés par les modèles
\providecommand{\abs}{}\let\abs\relax\DeclarePairedDelimiter{\abs}{\lvert}{\rvert}
\providecommand{\norm}{}\let\norm\relax\DeclarePairedDelimiter{\norm}{\lVert}{\rVert}
\usepackage[table]{xcolor} % [table] : fournit aussi \rowcolors (alternance de lignes)
\usepackage{pagecolor}
\usepackage{tikz}
\usetikzlibrary{arrows.meta,positioning,calc,shapes.geometric,shapes.misc,shapes.arrows,decorations.pathmorphing,decorations.pathreplacing,decorations.markings,patterns,shadows,mindmap,trees,fit,backgrounds,matrix,intersections,angles,quotes}
\usepackage{pgfplots}
\usepackage{pgf-pie} % diagrammes circulaires \pie (statistiques)
\pgfplotsset{compat=1.18}
\usepgfplotslibrary{fillbetween,groupplots} % aires entre courbes (calcul intégral)
\usepackage{enumitem}
\usepackage{fancyhdr}
% [most] charge toutes les bibliothèques tcolorbox (listings, minted, raster,
% documentation...) et gonfle inutilement la mémoire TeX sur un document long
% (~300 boîtes) : on ne charge que ce qui est réellement utilisé.
\usepackage[skins,breakable]{tcolorbox}
\usepackage{graphicx}
\usepackage{colortbl}
\usepackage{booktabs}
\usepackage{tabularx}
\usepackage{diagbox} % case d'en-tête diagonale des tableaux à double entrée
% Colonnes extensibles centrée / gauche / droite (C, L, R), souvent utilisées par les modèles
\newcolumntype{C}{>{\centering\arraybackslash}X}
\newcolumntype{L}{>{\raggedright\arraybackslash}X}
\newcolumntype{R}{>{\raggedleft\arraybackslash}X}
\usepackage{array}
\usepackage{multicol}
\usepackage{siunitx}
% parse-numbers=false : accepte aussi \SI{2,5 \times 10^{-3}}{...}
\sisetup{locale=FR,output-decimal-marker={,},group-minimum-digits=4,parse-numbers=false}
\DeclareSIUnit\ohm{\ensuremath{\symup{\Omega}}} % U+2126 absent de la police
\usepackage{qrcode}
% \degree : commande fréquemment utilisée par les modèles pour un angle ou une
% température en dehors de \SI{}{} (non fournie par les paquets chargés).
\providecommand{\degree}{^{\circ}}
% \qed (fin de démonstration), utilisé par les modèles sans amsthm
\providecommand{\qed}{\hfill$\square$}
% \barre : double barre des valeurs interdites dans un tableau de variations (tkz-tab)
\providecommand{\barre}{\ensuremath{\|}}
% environnement proof (amsthm non chargé) utilisé par les modèles
\ifdefined\proof\else\newenvironment{proof}[1][Démonstration]{\par\noindent\textit{#1.}\ }{\qed\par}\fi
\usepackage{lastpage}
\usepackage{hyperref}
\hypersetup{hidelinks}

% ============================================================
% Palette « Pop » OnBuch+ — mêmes hex que la classe P (plus_ui.dart)
% ============================================================
\definecolor{poporange}{HTML}{F59321}
\definecolor{popdark}{HTML}{E07A0C}
\definecolor{popcream}{HTML}{FFF6E8}
\definecolor{popink}{HTML}{1C1714}
\definecolor{poppurple}{HTML}{7A5AE0}
\definecolor{popgreen}{HTML}{2E9E6A}
\definecolor{poppink}{HTML}{E0457B}
\definecolor{popblue}{HTML}{2F7DE1}
\definecolor{popgold}{HTML}{C98A12}
\definecolor{popmuted}{HTML}{8A7F76}
\definecolor{popline}{HTML}{EADFD2}
\definecolor{popgray}{HTML}{8A7F76}
\colorlet{popgrey}{popgray}
% Alias tolérants (le modèle nomme parfois les couleurs différemment)
\colorlet{popwhite}{white}
\colorlet{popblack}{popink}
\colorlet{popred}{poppink}
\colorlet{popyellow}{popgold}
% Teintes claires pour fonds de tableaux / zones
\colorlet{popblueL}{popblue!10!white}
\colorlet{poporangeL}{poporange!14!white}
\colorlet{popgreenL}{popgreen!12!white}
\colorlet{poppurpleL}{poppurple!12!white}
\colorlet{poppinkL}{poppink!12!white}
\colorlet{popgoldL}{popgold!14!white}

\pagecolor{popcream}

% ============================================================
% Typographie
% ============================================================
\defaultfontfeatures{Ligatures=TeX}
\setmainfont{PlusJakartaSans-Medium.ttf}[Path=../../../fonts/,BoldFont=PlusJakartaSans-Bold.ttf]
% Maths en sans-serif (Fira Math) avec chiffres et lettres droites de Plus
% Jakarta, pour que les formules se fondent dans le texte.
\usepackage[math-style=ISO,bold-style=ISO]{unicode-math}
\setmathfont{FiraMath-Regular.otf}[Path=../../../fonts/]
\setmathfont{PlusJakartaSans-Medium.ttf}[Path=../../../fonts/,range={up/num,up/{Latin,latin},\mathup}]
\usepackage{icomma} % virgule décimale sans espace en mode math
% Caractères mathématiques Unicode absents des polices de texte (Plus Jakarta,
% Archivo Black) : rendus par leur équivalent mathématique, en texte comme en
% formule — sinon ils s'affichent en carré vide (ex. « Arithmétique dans ℤ »).
\usepackage{newunicodechar}
\newunicodechar{ℝ}{\ensuremath{\mathbb{R}}}
\newunicodechar{ℤ}{\ensuremath{\mathbb{Z}}}
\newunicodechar{ℂ}{\ensuremath{\mathbb{C}}}
\newunicodechar{ℕ}{\ensuremath{\mathbb{N}}}
\newunicodechar{ℚ}{\ensuremath{\mathbb{Q}}}
\newunicodechar{𝔻}{\ensuremath{\mathbb{D}}}
\newunicodechar{α}{\ensuremath{\alpha}}
\newunicodechar{β}{\ensuremath{\beta}}
\newunicodechar{γ}{\ensuremath{\gamma}}
\newunicodechar{δ}{\ensuremath{\delta}}
\newunicodechar{ε}{\ensuremath{\varepsilon}}
\newunicodechar{θ}{\ensuremath{\theta}}
\newunicodechar{λ}{\ensuremath{\lambda}}
\newunicodechar{μ}{\ensuremath{\mu}}
\newunicodechar{σ}{\ensuremath{\sigma}}
\newunicodechar{τ}{\ensuremath{\tau}}
\newunicodechar{φ}{\ensuremath{\varphi}}
\newunicodechar{ψ}{\ensuremath{\psi}}
\newunicodechar{ω}{\ensuremath{\omega}}
\newunicodechar{Σ}{\ensuremath{\Sigma}}
% majuscules produites par \MakeUppercase dans les titres (α-aminés -> Α)
\newunicodechar{Α}{\ensuremath{\alpha}}
\newunicodechar{Β}{\ensuremath{\beta}}
\newunicodechar{Γ}{\ensuremath{\Gamma}}
\newunicodechar{Δ}{\ensuremath{\Delta}}
\newunicodechar{Ε}{\ensuremath{\varepsilon}}
\newunicodechar{Θ}{\ensuremath{\Theta}}
\newunicodechar{Λ}{\ensuremath{\Lambda}}
\newunicodechar{Μ}{\ensuremath{\mu}}
\newunicodechar{Τ}{\ensuremath{\tau}}
\newunicodechar{Φ}{\ensuremath{\Phi}}
\newunicodechar{Ψ}{\ensuremath{\Psi}}
\newunicodechar{Ω}{\ensuremath{\Omega}}
\newunicodechar{↦}{\ensuremath{\mapsto}}
\newunicodechar{∈}{\ensuremath{\in}}
\newunicodechar{∉}{\ensuremath{\notin}}
\newunicodechar{∧}{\ensuremath{\wedge}}
\newunicodechar{⇔}{\ensuremath{\Leftrightarrow}}
\newunicodechar{⟺}{\ensuremath{\Longleftrightarrow}}
\newunicodechar{⇒}{\ensuremath{\Rightarrow}}
\newunicodechar{⟹}{\ensuremath{\Longrightarrow}}
\newunicodechar{∘}{\ensuremath{\circ}}
\newunicodechar{∪}{\ensuremath{\cup}}
\newunicodechar{∩}{\ensuremath{\cap}}
\newunicodechar{≡}{\ensuremath{\equiv}}
\newunicodechar{⊂}{\ensuremath{\subset}}
\newunicodechar{⊥}{\ensuremath{\perp}}
\newunicodechar{∃}{\ensuremath{\exists}}
\newunicodechar{∀}{\ensuremath{\forall}}
\newunicodechar{∛}{\ensuremath{\sqrt[3]{\,}}}
\newunicodechar{∜}{\ensuremath{\sqrt[4]{\,}}}
\newunicodechar{⃗}{\ensuremath{{}^{\to}}}
\newunicodechar{ⁿ}{\ifmmode^{n}\else\textsuperscript{\ensuremath{n}}\fi}
\newunicodechar{ˣ}{\ifmmode^{x}\else\textsuperscript{\ensuremath{x}}\fi}
\newunicodechar{ᵇ}{\ifmmode^{b}\else\textsuperscript{\ensuremath{b}}\fi}
\newunicodechar{ʳ}{\ifmmode^{r}\else\textsuperscript{\ensuremath{r}}\fi}
\newunicodechar{ᵘ}{\ifmmode^{u}\else\textsuperscript{\ensuremath{u}}\fi}
\newunicodechar{ʸ}{\ifmmode^{y}\else\textsuperscript{\ensuremath{y}}\fi}
\newunicodechar{ᵃ}{\ifmmode^{a}\else\textsuperscript{\ensuremath{a}}\fi}
\newunicodechar{ᵉ}{\ifmmode^{e}\else\textsuperscript{\ensuremath{e}}\fi}
\newunicodechar{ᵏ}{\ifmmode^{k}\else\textsuperscript{\ensuremath{k}}\fi}
\newunicodechar{ᵖ}{\ifmmode^{p}\else\textsuperscript{\ensuremath{p}}\fi}
\newunicodechar{ᴺ}{\ifmmode^{N}\else\textsuperscript{\ensuremath{N}}\fi}
\newunicodechar{ᶜ}{\ifmmode^{c}\else\textsuperscript{\ensuremath{c}}\fi}
\newunicodechar{ʰ}{\ifmmode^{h}\else\textsuperscript{\ensuremath{h}}\fi}
\newunicodechar{ᵠ}{\ifmmode^{q}\else\textsuperscript{\ensuremath{q}}\fi}
\newunicodechar{ₙ}{\ifmmode_{n}\else\textsubscript{\ensuremath{n}}\fi}
\newunicodechar{ₐ}{\ifmmode_{a}\else\textsubscript{\ensuremath{a}}\fi}
\newunicodechar{ᵢ}{\ifmmode_{i}\else\textsubscript{\ensuremath{i}}\fi}
\newunicodechar{ᵣ}{\ifmmode_{r}\else\textsubscript{\ensuremath{r}}\fi}
\newunicodechar{ₖ}{\ifmmode_{k}\else\textsubscript{\ensuremath{k}}\fi}
\newunicodechar{ᵦ}{\ifmmode_{\beta}\else\textsubscript{\ensuremath{\beta}}\fi}
\newunicodechar{ₓ}{\ifmmode_{x}\else\textsubscript{\ensuremath{x}}\fi}
\newfontfamily\popdisplay{ArchivoBlack-Regular.ttf}[Path=../../../fonts/]
\newfontfamily\poplogo{SpaceGrotesk-Bold.ttf}[Path=../../../fonts/]
\newfontfamily\popsemi{PlusJakartaSans-SemiBold.ttf}[Path=../../../fonts/]
\newfontfamily\popxbold{PlusJakartaSans-ExtraBold.ttf}[Path=../../../fonts/]
\newfontfamily\popxboldls{PlusJakartaSans-ExtraBold.ttf}[Path=../../../fonts/,LetterSpace=3.0]

\setlist[enumerate]{leftmargin=1.7em,itemsep=5pt,topsep=4pt}
\setlist[itemize]{leftmargin=1.7em,itemsep=4pt,topsep=4pt,label={\color{poporange}\textbullet}}
\setlength{\parindent}{0pt}
\setlength{\parskip}{5pt}
\renewcommand{\arraystretch}{1.25}

% pgfplots : style Pop par défaut
\pgfplotsset{
  every axis/.append style={axis line style={popink,line width=1pt},tick style={popink},
    grid style={popline,line width=0.5pt},label style={font=\small},tick label style={font=\footnotesize}},
  popcurve/.style={poporange,line width=1.8pt,no markers,smooth},
}
% Même style disponible dans un \draw TikZ simple (hors pgfplots)
\tikzset{popcurve/.style={poporange,line width=1.8pt}}

% ============================================================
% Pill / squiggle
% ============================================================
\newcommand{\poppill}[3]{%
  \tikz[baseline=-0.55ex]\node[draw=popink,line width=1.2pt,rounded corners=2.6pt,%
    fill=#2,inner xsep=6.5pt,inner ysep=3pt]{%
    \popxboldls\fontsize{7}{7}\selectfont\color{#3}\MakeUppercase{#1}};%
}
\newcommand{\popsquiggle}[1]{%
  \tikz[baseline=-0.4ex]{
    \draw[#1,line width=2.6pt,line cap=round]
      plot [smooth] coordinates {(0,0.09) (0.34,0.01) (0.68,0.21) (1.02,0.01) (1.36,0.10)};
  }%
}
% Mot-clé surligné (marqueur orange sous le texte)
\newcommand{\cle}[1]{{\popxbold\color{popdark}#1}}

% ============================================================
% En-tête / pied
% ============================================================
\pagestyle{fancy}
\fancyhf{}
\renewcommand{\headrulewidth}{0pt}
\renewcommand{\footrulewidth}{0pt}
\lhead{\raisebox{-0.15ex}{\poplogo\fontsize{13}{13}\selectfont\color{popink}OnBuch\color{poporange}+}}
\rhead{\poppill{\DOCMATIERE\ \textbullet\ \DOCNIVEAU\ \textbullet\ Leçon \DOCLECON}{white}{popink}}
\lfoot{\popxbold\fontsize{7.6}{7.6}\selectfont\color{popmuted}\MakeUppercase{Cours signé OnBuch+}\ \textbullet\ {\popsemi\DOCTITRE}}
\rfoot{\tikz[baseline=-0.6ex]\node[rounded corners=8.5pt,fill=popink,minimum height=17pt,minimum width=17pt,inner xsep=5pt,inner sep=0]{\popxbold\fontsize{8}{8}\selectfont\color{popcream}\thepage\,/\,\pageref*{LastPage}};}

% ============================================================
% Titres
% ============================================================
\newcommand{\chaptitre}[2]{%
  \begin{tcolorbox}[enhanced,colback=poporange,colframe=popink,boxrule=2.6pt,arc=11pt,
      drop shadow=popink,left=15pt,right=15pt,top=12pt,bottom=11pt,before skip=3pt,after skip=18pt]
    \poppill{#2}{popink}{popcream}\par\vspace{9pt}
    {\raggedright\hyphenpenalty=10000\exhyphenpenalty=10000\popdisplay\fontsize{22}{24}\selectfont\color{popcream}\MakeUppercase{#1}\par}
  \end{tcolorbox}
}
% Section numérotée : pastille orange avec le numéro + titre Archivo
\newcounter{popsec}
\newcommand{\coursec}[1]{%
  \stepcounter{popsec}\setcounter{popsub}{0}%
  \par\vspace{16pt}\noindent
  \begin{minipage}{\linewidth}
  \tikz[baseline=-0.7ex]\node[draw=popink,line width=1.6pt,fill=poporange,rounded corners=4pt,minimum size=22pt,inner sep=0]{\popdisplay\fontsize{12}{12}\selectfont\color{popcream}\thepopsec};%
  \hspace{8pt}{\popdisplay\fontsize{15}{17}\selectfont\color{popink}\MakeUppercase{#1}}\par
  \vspace{4pt}\noindent{\color{poporange}\rule{36pt}{3.6pt}}
  \end{minipage}\par\nopagebreak\vspace{8pt}
}
\newcounter{popsub}
\newcommand{\courssub}[1]{%
  \stepcounter{popsub}%
  \par\vspace{9pt}\noindent{\popxbold\fontsize{11.5}{13}\selectfont\color{popink}{\color{poporange}\thepopsec.\thepopsub}\hspace{6pt}#1}\par\nopagebreak\vspace{4pt}
}

% ============================================================
% Boîtes pédagogiques — même recette : fond blanc, bordure épaisse colorée,
% ombre dure encre, badge plein. Titre optionnel [#1].
% ============================================================
\tcbset{popbase/.style={breakable,enhanced,colback=white,boxrule=2.3pt,arc=9pt,
  shadow={3pt}{-3pt}{0pt}{popink},left=14pt,right=14pt,top=11pt,bottom=11pt,before skip=11pt,after skip=15pt,
  fontupper=\popsemi\fontsize{10.6}{14.5}\selectfont\color{popink},
  fontlower=\popsemi\fontsize{10.6}{14.5}\selectfont\color{popink}}}
% Coche dessinée (le glyphe ✓ n'existe pas dans les polices du document)
\newcommand{\popcheck}[1]{\tikz[baseline=-0.5ex]\draw[#1,line width=1.6pt,line cap=round,line join=round](-0.13,0.0)--(-0.04,-0.09)--(0.13,0.1);}
\newcommand{\popboxhead}[3]{\poppill{#1}{#2}{white}\ifx\relax#3\relax\else\hspace{6pt}{\popxbold\fontsize{10.6}{12}\selectfont\color{popink}#3}\fi\par\vspace{7pt}}

\newtcolorbox{aretenir}[1][]{popbase,colframe=popgreen,before upper={\popboxhead{À retenir}{popgreen}{#1}}}
\newtcolorbox{exemplebox}[1][]{popbase,colframe=poporange,before upper={\popboxhead{Exemple}{poporange}{#1}}}
\newtcolorbox{definition}[1][]{popbase,colframe=popblue,before upper={\popboxhead{Définition}{popblue}{#1}}}
\newtcolorbox{propriete}[1][]{popbase,colframe=poppurple,before upper={\popboxhead{Propriété}{poppurple}{#1}}}
\newtcolorbox{methode}[1][]{popbase,colframe=popgold,before upper={\popboxhead{Méthode}{popgold}{#1}}}
\newtcolorbox{attention}[1][]{popbase,colframe=poppink,before upper={\popboxhead{Attention}{poppink}{#1}}}
\newtcolorbox{experience}[1][]{popbase,colframe=popblue,colback=popblueL,before upper={\popboxhead{Expérience}{popblue}{#1}}}
% « Le savais-tu ? » : carte violette pleine (contexte, culture, Cameroun)
\newtcolorbox{savaistu}[1][]{popbase,colback=poppurple,colframe=popink,
  fontupper=\popsemi\fontsize{10.4}{14.2}\selectfont\color{white},
  before upper={\poppill{Le savais-tu\,?}{white}{poppurple}\ifx\relax#1\relax\else\hspace{6pt}{\popxbold\color{white}#1}\fi\par\vspace{7pt}}}
% Exercice résolu : énoncé en haut, \tcblower puis la solution sur fond crème
\newcounter{popexo}\newcounter{popexr}
\newtcolorbox{exoresolu}[1][]{popbase,colframe=poporange,colbacklower=popcream,
  segmentation style={popink,line width=1.2pt,dashed},
  before upper={\stepcounter{popexr}\popboxhead{Exercice résolu \thepopexr}{poporange}{#1}},
  before lower={\poppill{Solution}{popink}{popcream}\par\vspace{6pt}}}
% Exercice (non résolu en place) — la correction est dans la partie Corrigés
\newtcolorbox{exercice}[1][]{popbase,colframe=popink,
  before upper={\stepcounter{popexo}\popboxhead{Exercice \thepopexo}{popink}{#1}}}
% Corrigé
\newtcolorbox{corrige}[1][]{popbase,colframe=popgreen,colback=popgreenL,
  before upper={\popboxhead{Corrigé}{popgreen}{#1}}}
% Objectifs / prérequis / bilan
\newtcolorbox{objectifs}{popbase,colframe=popink,colback=poporangeL,
  before upper={\popboxhead{Objectifs de la leçon}{popink}{}}}
\newtcolorbox{prerequis}{popbase,colframe=popink,colback=popblueL,
  before upper={\popboxhead{Prérequis}{popblue}{}}}
\newtcolorbox{bilan}{popbase,colframe=popink,colback=white,boxrule=2.8pt,
  before upper={\popboxhead{Fiche bilan}{poporange}{L'essentiel à connaître pour l'examen}}}
% Figure encadrée avec légende
\newtcolorbox{popfigure}[1][]{enhanced,breakable=false,colback=white,colframe=popline,boxrule=1.4pt,arc=8pt,
  left=10pt,right=10pt,top=10pt,bottom=8pt,before skip=10pt,after skip=12pt,halign=center,
  lower separated=false,halign lower=center,
  fontlower=\popsemi\fontsize{9}{11.5}\selectfont\color{popmuted},
  #1}
\newcommand{\legende}[1]{\par\vspace{6pt}{\popsemi\fontsize{9}{11.5}\selectfont\color{popmuted}#1}}

% ============================================================
% Signature OnBuch+
% ============================================================
\newcommand{\poplogoinline}[1]{{\poplogo\fontsize{#1}{#1}\selectfont\color{popink}OnBuch\color{poporange}+}}
\newcommand{\signatureonbuch}{%
  \begin{tcolorbox}[enhanced,colback=popink,colframe=popink,boxrule=2.2pt,arc=10pt,
      drop shadow=poporange,left=14pt,right=14pt,top=10pt,bottom=10pt,before skip=0pt,after skip=0pt]
    \tikz[baseline=-0.7ex]\node[circle,fill=popgreen,draw=popcream,line width=1.3pt,minimum size=22pt,inner sep=0]{\popcheck{white}};%
    \hspace{9pt}\begin{minipage}[c]{0.62\linewidth}
      {\popxbold\fontsize{12}{13}\selectfont\color{popcream}Signé\ \textbullet\ {\poplogo OnBuch\color{poporange}+}}\par\vspace{2pt}
      {\raggedright\popsemi\fontsize{8}{10.5}\selectfont\color{popline}\MakeUppercase{Équipe pédagogique OnBuch+}\\\MakeUppercase{Conforme au programme officiel MINESEC}\par}
    \end{minipage}\hfill
    {\poplogo\fontsize{22}{22}\selectfont\color{popcream}OnBuch\color{poporange}+}
  \end{tcolorbox}%
}

% ============================================================
% Page de garde
% #1 titre · #2 matière · #3 niveau · #4 module · #5 n° leçon · #6 durée ·
% #7 description / objectifs (texte ou itemize)
% ============================================================
\newcommand{\pagegarde}[7]{%
  \thispagestyle{empty}
  \newgeometry{a4paper,margin=1.7cm,top=1.6cm,bottom=1.4cm}%
  \begin{tikzpicture}[remember picture,overlay]
    \fill[poporange!18] ([xshift=-3.2cm,yshift=-3.6cm]current page.north east) circle (4.6cm);
    \fill[poppurple!14] ([xshift=2.4cm,yshift=7.6cm]current page.south west) circle (3.8cm);
  \end{tikzpicture}%
  {\poplogo\fontsize{30}{30}\selectfont\color{popink}OnBuch\color{poporange}+}\hfill
  \raisebox{0.6ex}{\poppill{Cours signé \textbullet\ Premium}{popink}{popcream}}\par
  \vspace{1pt}\popsquiggle{poppurple}\par
  \vspace{8pt}
  \begin{tcolorbox}[enhanced,colback=poporange,colframe=popink,boxrule=2.8pt,arc=13pt,
      drop shadow=popink,left=18pt,right=18pt,top=12pt,bottom=13pt,after skip=10pt]
    \poppill{#2 \textbullet\ #3}{popink}{popcream}\hspace{5pt}\poppill{#4}{white}{popink}\par\vspace{8pt}
    {\popdisplay\fontsize{14}{14}\selectfont\color{popink}LEÇON #5}\par\vspace{4pt}
    {\raggedright\hyphenpenalty=10000\exhyphenpenalty=10000\popdisplay\fontsize{25}{27}\selectfont\color{popcream}\MakeUppercase{#1}\par}
    \vspace{8pt}
    {\popxbold\fontsize{9.5}{9.5}\selectfont\color{popink}\MakeUppercase{Durée conseillée : #6}}
  \end{tcolorbox}
  \begin{tcolorbox}[enhanced,colback=white,colframe=popink,boxrule=2.2pt,arc=10pt,
      drop shadow=popink,left=16pt,right=16pt,top=10pt,bottom=10pt,after skip=8pt]
    \poppill{Dans cette leçon}{poppurple}{white}\par\vspace{5pt}
    \popsemi\fontsize{9.3}{12.6}\selectfont\color{popink}#7
  \end{tcolorbox}
  \noindent\begin{minipage}[t]{0.487\linewidth}
    \begin{tcolorbox}[enhanced,colback=popgreen,colframe=popink,boxrule=2.2pt,arc=10pt,
        drop shadow=popink,left=11pt,right=11pt,top=10pt,bottom=10pt,height=3.1cm,valign=center]
      \tcbox[colback=white,colframe=popink,boxrule=1.3pt,arc=5pt,boxsep=0pt,nobeforeafter,
        left=3pt,right=3pt,top=3pt,bottom=3pt,valign=center]{\qrcode[height=1.9cm]{https://play.google.com/store/apps/details?id=com.luvvix.onbuch&hl=fr}}%
      \hspace{8pt}\begin{minipage}[c]{0.5\linewidth}
        {\popdisplay\fontsize{11}{12.5}\selectfont\color{white}\MakeUppercase{Télécharge OnBuch}}\par\vspace{4pt}
        {\raggedright\popsemi\fontsize{8.4}{11}\selectfont\color{white}Gratuit \textbullet\ Google Play\\Tuteur IA, annales, concours, résultats\par}
      \end{minipage}
    \end{tcolorbox}
  \end{minipage}\hfill
  \begin{minipage}[t]{0.487\linewidth}
    \begin{tcolorbox}[enhanced,colback=poppurple,colframe=popink,boxrule=2.2pt,arc=10pt,
        drop shadow=popink,left=11pt,right=11pt,top=10pt,bottom=10pt,height=3.1cm,valign=center]
      \tikz[baseline=-0.7ex]\node[circle,fill=white,draw=popink,line width=1.3pt,minimum size=1.6cm,inner sep=0]{\popdisplay\fontsize{24}{24}\selectfont\color{poppurple}+};%
      \hspace{8pt}\begin{minipage}[c]{0.6\linewidth}
        {\popdisplay\fontsize{11}{12.5}\selectfont\color{white}\MakeUppercase{Passe à OnBuch+}}\par\vspace{4pt}
        {\raggedright\popsemi\fontsize{8.4}{11}\selectfont\color{white}Tous les cours signés, Tuteur IA illimité, fiches PDF — onglet OnBuch+ de l'app\par}
      \end{minipage}
    \end{tcolorbox}
  \end{minipage}
  \vspace{8pt}
  \popcontenu
  \vfill
  \signatureonbuch
  \restoregeometry
  \newpage
}

% Rangée « Ce cours contient » (page de garde)
\newcommand{\poptile}[3]{% #1 couleur #2 gros texte #3 légende
  \begin{tcolorbox}[enhanced,colback=#1,colframe=popink,boxrule=2pt,arc=9pt,shadow={2.5pt}{-2.5pt}{0pt}{popink},
      left=6pt,right=6pt,top=9pt,bottom=9pt,halign=center,valign=center,height=1.75cm,width=\linewidth,nobeforeafter]
    {\popdisplay\fontsize{12.5}{13}\selectfont\color{white}\MakeUppercase{#2}}\par\vspace{3pt}
    {\centering\popsemi\fontsize{7.8}{9.5}\selectfont\color{white}#3\par}
  \end{tcolorbox}}
\newcommand{\popcontenu}{%
  \par\noindent{\popxboldls\fontsize{7.5}{7.5}\selectfont\color{popmuted}CE COURS SIGNÉ CONTIENT}\par\vspace{7pt}
  \noindent\begin{minipage}[t]{0.235\linewidth}\poptile{popblue}{Cours}{complet, illustré, pas à pas}\end{minipage}\hfill
  \begin{minipage}[t]{0.235\linewidth}\poptile{poporange}{Méthodes}{et exercices résolus}\end{minipage}\hfill
  \begin{minipage}[t]{0.235\linewidth}\poptile{poppink}{Exercices}{type examen + corrigés}\end{minipage}\hfill
  \begin{minipage}[t]{0.235\linewidth}\poptile{popgreen}{Fiche bilan}{l'essentiel à retenir}\end{minipage}\par}

% Sommaire
\newtcolorbox{sommaire}{enhanced,colback=white,colframe=popink,boxrule=2.2pt,arc=10pt,
  drop shadow=popink,left=18pt,right=18pt,top=13pt,bottom=13pt,before skip=6pt,after skip=20pt,
  before upper={\poppill{Sommaire}{poppurple}{white}\par\vspace{9pt}\popsemi\fontsize{10.6}{14.5}\selectfont\color{popink}}}

% Signature de fin de cours — poussée tout en bas de la dernière page
\newcommand{\findecours}{%
  \par\vfill\nopagebreak
  \begin{center}\popsquiggle{poporange}\end{center}\vspace{4pt}
  \signatureonbuch
}
\AtBeginDocument{\def\llbracket{\mathopen{[\mkern-3.5mu[}}\def\rrbracket{\mathclose{]\mkern-3.5mu]}}} % intervalles d'entiers (glyphe ⟦ absent de la police)
% Glyphes absents de Fira Math : redéfinis à partir de glyphes disponibles
\AtBeginDocument{%
  \def\setminus{\mathbin{\backslash}}\let\smallsetminus\setminus
  \def\vdots{\mathord{\vbox{\baselineskip3.2pt\lineskiplimit0pt\kern4pt\hbox{.}\hbox{.}\hbox{.}}}}%
  \def\ddots{\mathinner{\mkern1mu\raise6.5pt\hbox{.}\mkern2mu\raise3.6pt\hbox{.}\mkern2mu\raise0.7pt\hbox{.}\mkern1mu}}%
  \def\triangle{\mathord{\scalebox{0.9}{$\symup{\Delta}$}}}%
  \def\nabla{\mathord{\raisebox{\depth}{\scalebox{1}[-1]{$\symup{\Delta}$}}}}%
  \def\checkmark{\popcheck{popdark}}%
  \def\star{\mathbin{\ast}}\def\bigstar{\mathord{\ast}}%
  \def\diamond{\mathbin{\scalebox{0.6}{$\Diamond$}}}%
  \def\bigcup{\mathop{\mathchoice{\raisebox{-0.3em}{\scalebox{1.7}{$\cup$}}}{\scalebox{1.25}{$\cup$}}{\cup}{\cup}}\displaylimits}%
  \def\bigcap{\mathop{\mathchoice{\raisebox{-0.3em}{\scalebox{1.7}{$\cap$}}}{\scalebox{1.25}{$\cap$}}{\cap}{\cap}}\displaylimits}%
  \def\lesseqqgtr{\mathrel{\lessgtr}}\def\bigoplus{\mathop{\oplus}\displaylimits}%
}
\providecommand{\ssi}{si et seulement si} % abréviation fréquente
\AtBeginDocument{\providecommand{\wideparen}{\overparen}} % arc de cercle

%% FIGFIX-BEGIN v3 — correctifs globaux des figures (docs/onbuch-generation/tools/figfix)
\usepackage{adjustbox}\usepackage{etoolbox}
% toute figure TikZ/pgfplots plus large que la ligne est réduite à la largeur disponible
\BeforeBeginEnvironment{tikzpicture}{\begin{adjustbox}{max width=\linewidth}}
\AfterEndEnvironment{tikzpicture}{\end{adjustbox}}
% légendes pgfplots lisibles par-dessus les courbes
\pgfplotsset{every axis/.append style={legend style={fill=white,fill opacity=0.92,text opacity=1,draw=black!15,inner sep=3pt}}}
% étiquettes d'arêtes (syntaxe edge["..."]) avec fond blanc
\tikzset{every edge quotes/.append style={fill=white,inner sep=1.2pt}}
% bulles de cartes mentales : texte à la ligne dans la bulle, bulle un peu plus grande
\tikzset{every concept/.append style={text width=2.0cm,align=center,minimum size=2.3cm,inner sep=2pt}}
%% FIGFIX-END
\begin{document}

\pagegarde{Matrices et applications linéaires d'un plan vectoriel}{Mathématiques}{1ère C}{Configurations et transformations élémentaires du plan}{N17}{14 h}{Cette leçon introduit les matrices comme tableaux de nombres et leurs opérations (somme, produit par un réel, produit, déterminant, inverse), puis les relie aux applications linéaires du plan vectoriel. Elle fournit un outil algébrique puissant pour résoudre des systèmes d'équations et étudier les transformations géométriques du plan.
\vspace{4pt}
\begin{multicols}{2}\small
\begin{itemize}
\item À la fin de cette leçon, je sais présenter des données concrètes sous forme de matrice et interpréter une matrice donnée
\item À la fin de cette leçon, je sais calculer la somme de deux matrices et le produit d'une matrice par un réel
\item À la fin de cette leçon, je sais calculer le produit de deux matrices en respectant les conditions de compatibilité
\item À la fin de cette leçon, je sais calculer le déterminant d'une matrice carrée d'ordre 2 et l'utiliser pour tester son inversibilité
\item À la fin de cette leçon, je sais déterminer l'inverse d'une matrice carrée d'ordre 2 inversible
\item À la fin de cette leçon, je sais résoudre un système de deux équations linéaires à deux inconnues par la méthode matricielle
\end{itemize}\end{multicols}}

\begin{sommaire}
\begin{multicols}{2}
\begin{enumerate}
\item Présentation des matrices
\item Somme de matrices et produit par un réel
\item Produit de deux matrices
\item Déterminant et inverse d'une matrice carrée d'ordre 2
\item Matrice d'une application linéaire dans une base
\item Opérations sur les applications linéaires et matrices
\item Activité d'intégration
\item Méthodes et exercices résolus
\item Exercices
\item Corrigés des exercices
\item Fiche bilan
\end{enumerate}
\end{multicols}
\end{sommaire}

\begin{objectifs}
\begin{itemize}
\item À la fin de cette leçon, je sais présenter des données concrètes sous forme de matrice et interpréter une matrice donnée
\item À la fin de cette leçon, je sais calculer la somme de deux matrices et le produit d'une matrice par un réel
\item À la fin de cette leçon, je sais calculer le produit de deux matrices en respectant les conditions de compatibilité
\item À la fin de cette leçon, je sais calculer le déterminant d'une matrice carrée d'ordre 2 et l'utiliser pour tester son inversibilité
\item À la fin de cette leçon, je sais déterminer l'inverse d'une matrice carrée d'ordre 2 inversible
\item À la fin de cette leçon, je sais résoudre un système de deux équations linéaires à deux inconnues par la méthode matricielle
\item À la fin de cette leçon, je sais écrire la matrice d'une application linéaire du plan dans une base donnée
\item À la fin de cette leçon, je sais calculer les coordonnées de l'image d'un vecteur par un produit matriciel
\item À la fin de cette leçon, je sais déterminer la matrice d'une somme, d'un produit par un réel, d'une composée d'applications linéaires et de la réciproque d'un automorphisme
\end{itemize}
\end{objectifs}

\begin{prerequis}
\begin{itemize}
\item Vecteurs du plan, coordonnées d'un vecteur dans une base (Seconde)
\item Applications linéaires du plan : définition et propriétés (leçon précédente de Première)
\item Transformations du plan : symétries, homothéties, rotations (Seconde et Première)
\item Systèmes de deux équations à deux inconnues : méthodes de substitution et de combinaison (Troisième/Seconde)
\item Calcul littéral et résolution d'équations du premier degré
\end{itemize}
\end{prerequis}

\newpage

\coursec{Présentation des matrices}

À la coopérative de cacao de Sangmélima, le gérant enregistre chaque semaine les quantités achetées à trois villages et les prix pratiqués selon deux qualités de fèves. Pour calculer rapidement les recettes, les stocks et les bénéfices, il range ses données dans des tableaux rectangulaires de nombres. Comment additionner ces tableaux, les multiplier par un coefficient (inflation des prix !), ou les combiner entre eux ? Et comment retrouver les quantités à partir des recettes ? Ces tableaux, appelés \cle{matrices}, sont aussi l'outil privilégié pour décrire les transformations du plan : symétries, rotations, homothéties\ldots{} C'est toute la puissance de cette leçon.

Dans cette première section, nous apprenons à \emph{lire} et à \emph{écrire} ces tableaux : qu'est-ce qu'une matrice ? comment repérer chacun de ses coefficients ? quelles sont les matrices remarquables ? et quel est le lien avec les vecteurs du plan que tu connais déjà ?

\courssub{Une situation concrète : les ventes de la coopérative}

Le gérant de la coopérative de Sangmélima achète du cacao à trois villages : Ambam, Ebolowa et Kribi. Il distingue deux qualités de fèves : la qualité A (fèves fermentées de premier choix) et la qualité B (fèves ordinaires). Voici les quantités achetées cette semaine, en tonnes :

\begin{center}
\begin{tabularx}{\linewidth}{|l|X|X|X|}
\hline
\rowcolor{popblueL} \textbf{Qualité} & \textbf{Ambam} & \textbf{Ebolowa} & \textbf{Kribi} \\
\hline
Qualité A & $12$ & $9$ & $15$ \\
\hline
Qualité B & $20$ & $14$ & $8$ \\
\hline
\end{tabularx}
\end{center}

Observons ce tableau : il contient $2$ lignes (une par qualité) et $3$ colonnes (une par village). Une fois les étiquettes enlevées, il ne reste qu'un \emph{tableau rectangulaire de nombres} :

\[
\begin{pmatrix} 12 & 9 & 15 \\ 20 & 14 & 8 \end{pmatrix}
\]

Ce tableau nu, c'est exactement ce que les mathématiciens appellent une \cle{matrice}. Toute l'information est conservée, à condition de se souvenir de ce que représentent les lignes et les colonnes. Par exemple, le nombre $14$ se lit à l'intersection de la $2^{\text{e}}$ ligne (qualité B) et de la $2^{\text{e}}$ colonne (Ebolowa) : la coopérative a acheté $14$ tonnes de qualité B à Ebolowa.

\courssub{Définition d'une matrice, format et coefficients}

\begin{definition}[Matrice, format, coefficient]
Soient $n$ et $p$ deux entiers naturels non nuls. Une \cle{matrice} de \cle{format} $(n,p)$ est un tableau rectangulaire de nombres réels comportant $n$ \cle{lignes} et $p$ \cle{colonnes}. On la note :
\[
A = \begin{pmatrix} a_{11} & a_{12} & \cdots & a_{1p} \\ a_{21} & a_{22} & \cdots & a_{2p} \\ \vdots & \vdots & \ddots & \vdots \\ a_{n1} & a_{n2} & \cdots & a_{np} \end{pmatrix}
\]
Le nombre $a_{ij}$ est le \cle{coefficient} (ou \emph{terme}) de la matrice situé à l'intersection de la $i$-ème ligne et de la $j$-ème colonne. On écrit aussi $A = (a_{ij})_{1 \le i \le n,\; 1 \le j \le p}$, ou simplement $A = (a_{ij})$.
\end{definition}

Retenons l'essentiel : dans l'écriture $a_{ij}$, le \textbf{premier indice désigne toujours la ligne}, le \textbf{second indice désigne toujours la colonne}. Une matrice de format $(n,p)$ possède donc $n \times p$ coefficients.

\begin{exemplebox}[Lecture de coefficients sur la matrice des ventes]
Considérons la matrice des ventes de la coopérative :
\[
V = \begin{pmatrix} 12 & 9 & 15 \\ 20 & 14 & 8 \end{pmatrix}
\]
C'est une matrice de format $(2,3)$ : $2$ lignes, $3$ colonnes, donc $6$ coefficients.
\begin{itemize}
\item $a_{11} = 12$ : $12$ tonnes de qualité A achetées à Ambam ;
\item $a_{12} = 9$ : $9$ tonnes de qualité A achetées à Ebolowa ;
\item $a_{13} = 15$ : $15$ tonnes de qualité A achetées à Kribi ;
\item $a_{23} = 8$ : $8$ tonnes de qualité B achetées à Kribi.
\end{itemize}
\end{exemplebox}

\begin{attention}[Ligne d'abord, colonne ensuite !]
L'erreur la plus fréquente est de confondre les deux indices de $a_{ij}$. Retiens la phrase : \emph{« ligne d'abord, colonne ensuite »}, ou l'astuce mnémotechnique : dans « $ij$ », $i$ (la ligne) vient avant $j$ (la colonne), comme dans l'alphabet. Ainsi, sur la matrice $V$ ci-dessus, $a_{12} = 9$ alors que le coefficient situé ligne $2$, colonne $1$ est $a_{21} = 20$. Ce ne sont pas les mêmes nombres ! Au Probatoire, échanger les indices fausse tous les calculs qui suivent.
\end{attention}

Pour bien visualiser la position d'un coefficient, observons une matrice carrée à $3$ lignes et $3$ colonnes :

\begin{popfigure}
\begin{center}
\begin{tikzpicture}[scale=1.05]
% Grille de la matrice 3x3
\foreach \i in {0,1,2,3} {
  \draw[popline] (\i,0) -- (\i,-3);
  \draw[popline] (0,-\i) -- (3,-\i);
}
% Coefficients
\node at (0.5,-0.5) {$a_{11}$};
\node at (1.5,-0.5) {$a_{12}$};
\node at (2.5,-0.5) {$a_{13}$};
\node at (0.5,-1.5) {$a_{21}$};
\node at (1.5,-1.5) {$a_{22}$};
\node[popdark] at (2.5,-1.5) {$a_{23}$};
\node at (0.5,-2.5) {$a_{31}$};
\node at (1.5,-2.5) {$a_{32}$};
\node at (2.5,-2.5) {$a_{33}$};
% Mise en évidence de a23
\fill[poporange!35] (2,-1) rectangle (3,-2);
\node[popdark] at (2.5,-1.5) {$a_{23}$};
% Flèches lignes et colonnes
\draw[->,popblue,thick] (-0.4,0) -- (-0.4,-2.7) node[midway,left,popblue] {\footnotesize lignes $i$};
\draw[->,popgreen,thick] (0,0.4) -- (2.7,0.4) node[midway,above,popgreen] {\footnotesize colonnes $j$};
% Annotations ligne 2 et colonne 3
\draw[<->,poporange,thick] (-1.6,-1) -- (-1.6,-2) node[midway,left,poporange] {\footnotesize $2^{\text{e}}$ ligne};
\draw[<->,poporange,thick] (2,0.9) -- (3,0.9) node[midway,above,poporange] {\footnotesize $3^{\text{e}}$ colonne};
\end{tikzpicture}
\end{center}
\legende{Le coefficient $a_{23}$ se trouve à l'intersection de la $2^{\text{e}}$ ligne et de la $3^{\text{e}}$ colonne.}
\end{popfigure}

\courssub{Matrices particulières}

Certaines matrices, par leur forme ou leurs coefficients, jouent un rôle si important qu'on leur a donné des noms.

\begin{definition}[Matrices particulières]
\begin{itemize}
\item Une \cle{matrice ligne} est une matrice de format $(1,p)$ : elle n'a qu'une seule ligne. Exemple : $\begin{pmatrix} 3 & -1 & 7 \end{pmatrix}$ est de format $(1,3)$.
\item Une \cle{matrice colonne} est une matrice de format $(n,1)$ : elle n'a qu'une seule colonne. Exemple : $\begin{pmatrix} 5 \\ -2 \end{pmatrix}$ est de format $(2,1)$.
\item Une \cle{matrice carrée d'ordre $n$} est une matrice de format $(n,n)$ : elle a autant de lignes que de colonnes. Les coefficients $a_{11}, a_{22}, \ldots, a_{nn}$ forment sa \cle{diagonale principale}.
\item La \cle{matrice nulle}, notée $0$ (ou $0_{n,p}$ si l'on veut préciser le format), est la matrice dont tous les coefficients sont nuls.
\item La \cle{matrice unité d'ordre $2$}, notée $I_2$, est la matrice carrée :
\[
I_2 = \begin{pmatrix} 1 & 0 \\ 0 & 1 \end{pmatrix}
\]
Ses coefficients diagonaux valent $1$, tous les autres valent $0$. On définit de même $I_3 = \begin{pmatrix} 1 & 0 & 0 \\ 0 & 1 & 0 \\ 0 & 0 & 1 \end{pmatrix}$, etc.
\end{itemize}
\end{definition}

\begin{aretenir}[À retenir]
\begin{itemize}
\item Le format $(n,p)$ se lit toujours : $n$ lignes, $p$ colonnes.
\item Une matrice carrée d'ordre $2$ a la forme $\begin{pmatrix} a & b \\ c & d \end{pmatrix}$ ; sa diagonale principale est formée de $a$ et $d$.
\item La matrice unité $I_2$ jouera pour les matrices le rôle que joue le nombre $1$ pour les réels : nous verrons qu'elle est « neutre » pour le produit.
\end{itemize}
\end{aretenir}

\begin{exemplebox}[Quelques formats à identifier]
Identifions le format de chacune des matrices suivantes :
\[
A = \begin{pmatrix} 1 & 4 & -2 \\ 0 & 3 & 5 \end{pmatrix}, \qquad
B = \begin{pmatrix} 2 & 1 \\ -3 & 0 \\ 7 & 6 \end{pmatrix}, \qquad
C = \begin{pmatrix} 4 & -1 \\ 2 & 9 \end{pmatrix}, \qquad
D = \begin{pmatrix} 1 & 2 & 3 \\ 0 & 5 & 6 \\ 0 & 0 & 7 \end{pmatrix}.
\]
\begin{itemize}
\item $A$ a $2$ lignes et $3$ colonnes : format $(2,3)$ ;
\item $B$ a $3$ lignes et $2$ colonnes : format $(3,2)$ — attention, ce n'est \emph{pas} le même format que $A$ ;
\item $C$ est une matrice carrée d'ordre $2$ ;
\item $D$ est une matrice carrée d'ordre $3$ (on dit qu'elle est \emph{triangulaire supérieure} car tous ses coefficients sous la diagonale sont nuls).
\end{itemize}
\end{exemplebox}

\courssub{Égalité de deux matrices}

Quand peut-on dire que deux tableaux de nombres sont « les mêmes » ? La réponse est naturelle : lorsqu'ils ont la même forme et les mêmes contenus, case par case.

\begin{definition}[Égalité de deux matrices]
Deux matrices $A = (a_{ij})$ et $B = (b_{ij})$ sont \cle{égales}, et l'on écrit $A = B$, lorsque :
\begin{enumerate}
\item elles ont le \textbf{même format} $(n,p)$ ;
\item leurs coefficients sont \textbf{égaux deux à deux} : pour tous $i$ et $j$, $a_{ij} = b_{ij}$.
\end{enumerate}
\end{definition}

\begin{attention}[Deux conditions, pas une seule]
Il ne suffit pas que deux matrices contiennent « les mêmes nombres » pour être égales. Par exemple :
\[
\begin{pmatrix} 1 & 2 \\ 3 & 4 \end{pmatrix} \neq \begin{pmatrix} 1 & 3 \\ 2 & 4 \end{pmatrix} \qquad \text{et} \qquad \begin{pmatrix} 1 & 2 \end{pmatrix} \neq \begin{pmatrix} 1 \\ 2 \end{pmatrix}.
\]
Dans le premier cas, les coefficients ne sont pas aux mêmes places ; dans le second, les formats $(1,2)$ et $(2,1)$ sont différents.
\end{attention}

L'égalité de matrices est un outil précieux pour résoudre des équations : une seule égalité matricielle entre matrices $(2,2)$ équivaut à \emph{quatre} égalités entre nombres réels.

\begin{exoresolu}[Une équation matricielle]
Énoncé. Déterminer les réels $x$ et $y$ tels que :
\[
\begin{pmatrix} x + 1 & 3 \\ 0 & 2y - 4 \end{pmatrix} = \begin{pmatrix} 5 & 3 \\ 0 & 8 \end{pmatrix}.
\]
\tcblower
Solution détaillée. Les deux matrices ont le même format $(2,2)$. Par définition de l'égalité, leurs coefficients sont égaux deux à deux :
\begin{align*}
x + 1 &= 5 & 3 &= 3 \\
0 &= 0 & 2y - 4 &= 8
\end{align*}
Les équations $3 = 3$ et $0 = 0$ sont automatiquement vérifiées. Il reste :
\[
x = 5 - 1 = 4 \qquad \text{et} \qquad 2y = 8 + 4 = 12 \ \text{donc} \ y = 6.
\]
Conclusion : $x = 4$ et $y = 6$.
\end{exoresolu}

\courssub{Matrices et vecteurs du plan}

Tu sais depuis la Seconde qu'un vecteur du plan muni d'une base $(\vec{i}, \vec{j})$ est repéré par un couple de coordonnées : $\vec{u}(x ; y)$ signifie $\vec{u} = x\vec{i} + y\vec{j}$. Or un couple ordonné de réels peut s'écrire verticalement : c'est une matrice colonne de format $(2,1)$ !

\begin{definition}[Matrice colonne associée à un vecteur]
Le plan étant muni d'une base $(\vec{i}, \vec{j})$, au vecteur $\vec{u}(x ; y)$ on associe la \cle{matrice colonne} de ses coordonnées :
\[
\vec{u} = x\vec{i} + y\vec{j} \quad \longleftrightarrow \quad \begin{pmatrix} x \\ y \end{pmatrix}
\]
La première ligne contient l'abscisse, la deuxième ligne contient l'ordonnée.
\end{definition}

Cette correspondance est le pont entre la géométrie (vecteurs, transformations du plan) et le calcul matriciel : c'est elle qui nous permettra, dans les sections suivantes, de décrire les symétries, rotations et homothéties par des matrices.

\begin{popfigure}
\begin{center}
\begin{tikzpicture}[scale=1]
% Repère
\draw[->,popline] (-0.5,0) -- (4.2,0) node[below left] {\footnotesize $x$};
\draw[->,popline] (0,-0.5) -- (0,3.2) node[below left] {\footnotesize $y$};
\draw[->,popblue,thick] (0,0) -- (1,0) node[midway,below] {\footnotesize $\vec{i}$};
\draw[->,popblue,thick] (0,0) -- (0,1) node[midway,left] {\footnotesize $\vec{j}$};
% Vecteur u
\draw[->,poporange,very thick] (0,0) -- (3,2) node[above right] {$\vec{u}(3\,;2)$};
\draw[dashed,popmuted] (3,0) -- (3,2);
\draw[dashed,popmuted] (0,2) -- (3,2);
\node[below] at (3,0) {\footnotesize $3$};
\node[left] at (0,2) {\footnotesize $2$};
% Flèche de correspondance
\draw[<->,popdark,thick] (4.6,1.4) -- (5.6,1.4);
% Matrice colonne
\node[align=center] at (6.7,1.4) {$\begin{pmatrix} 3 \\ 2 \end{pmatrix}$};
\node[popmuted] at (6.7,0.3) {\footnotesize matrice colonne $(2,1)$};
\end{tikzpicture}
\end{center}
\legende{Au vecteur $\vec{u}(3\,;2)$ correspond la matrice colonne $\begin{pmatrix} 3 \\ 2 \end{pmatrix}$ : l'abscisse en haut, l'ordonnée en bas.}
\end{popfigure}

\begin{exemplebox}[Des vecteurs aux matrices colonnes]
Dans une base $(\vec{i}, \vec{j})$ du plan, donnons les matrices colonnes associées :
\[
\vec{u}(5 ; -1) \leftrightarrow \begin{pmatrix} 5 \\ -1 \end{pmatrix}, \qquad
\vec{v}(0 ; 4) \leftrightarrow \begin{pmatrix} 0 \\ 4 \end{pmatrix}, \qquad
\vec{0}(0 ; 0) \leftrightarrow \begin{pmatrix} 0 \\ 0 \end{pmatrix}.
\]
Réciproquement, la matrice colonne $\begin{pmatrix} -7 \\ 2 \end{pmatrix}$ représente le vecteur $\vec{w} = -7\vec{i} + 2\vec{j}$, de coordonnées $(-7 ; 2)$.
\end{exemplebox}

\begin{methode}[Bien démarrer avec les matrices]
\begin{enumerate}
\item \textbf{Identifier le format} : compter les lignes, puis les colonnes ; écrire $(n,p)$.
\item \textbf{Repérer un coefficient} : dans $a_{ij}$, lire d'abord la ligne $i$, puis la colonne $j$.
\item \textbf{Traduire une égalité} : deux matrices égales donnent une équation par coefficient.
\item \textbf{Passer du vecteur à la matrice} : les coordonnées $(x ; y)$ s'écrivent en colonne $\begin{pmatrix} x \\ y \end{pmatrix}$, dans le même ordre.
\end{enumerate}
\end{methode}

\begin{savaistu}[Les matrices sont partout !]
Les matrices ne sont pas qu'un jeu de mathématiciens : elles font tourner le monde moderne. En \textbf{informatique}, une image numérique n'est autre qu'une matrice de pixels : chaque case contient un nombre qui code la couleur d'un point de l'image. Retoucher une photo sur un smartphone, c'est effectuer des calculs sur une matrice de plusieurs millions de coefficients ! En \textbf{économie}, les tableaux de flux entre branches d'activité (qui achète quoi à qui) sont des matrices, utilisées pour prévoir la production. En \textbf{téléphonie mobile}, les opérateurs comme MTN ou Orange représentent les appels entre antennes relais par des matrices pour optimiser leur réseau. Et c'est le mathématicien britannique Arthur Cayley qui, vers 1858, a le premier défini les matrices et leurs opérations comme nous les étudions aujourd'hui.
\end{savaistu}

\begin{exercice}[Pour s'entraîner]
Le gérant de la coopérative enregistre la semaine suivante les quantités (en tonnes) : qualité A : $10$ à Ambam, $11$ à Ebolowa, $13$ à Kribi ; qualité B : $18$ à Ambam, $16$ à Ebolowa, $10$ à Kribi.
\begin{enumerate}
\item Écrire la matrice $W$ des ventes de cette semaine et préciser son format.
\item Que valent $w_{13}$ et $w_{21}$ ? Interpréter chacun de ces nombres.
\item Déterminer les réels $a$ et $b$ tels que $\begin{pmatrix} a & 11 \\ 18 & b \end{pmatrix} = \begin{pmatrix} w_{11} & w_{12} \\ w_{21} & w_{22} \end{pmatrix}$.
\end{enumerate}
\end{exercice}

\begin{corrige}[Exercice 1]
\begin{enumerate}
\item En rangeant les qualités en lignes et les villages en colonnes (Ambam, Ebolowa, Kribi) :
\[
W = \begin{pmatrix} 10 & 11 & 13 \\ 18 & 16 & 10 \end{pmatrix}, \quad \text{format } (2,3).
\]
\item $w_{13} = 13$ : la coopérative a acheté $13$ tonnes de qualité A à Kribi. $w_{21} = 18$ : elle a acheté $18$ tonnes de qualité B à Ambam.
\item Par égalité des coefficients deux à deux : $a = w_{11} = 10$ et $b = w_{22} = 16$.
\end{enumerate}
\end{corrige}

\begin{aretenir}[L'essentiel de la section]
\begin{itemize}
\item Une \cle{matrice} de format $(n,p)$ est un tableau rectangulaire de $n$ lignes et $p$ colonnes de nombres réels.
\item Le coefficient $a_{ij}$ se lit à l'intersection de la $i$-ème \textbf{ligne} et de la $j$-ème \textbf{colonne} : ligne d'abord, colonne ensuite.
\item Matrices remarquables : matrice ligne $(1,p)$, matrice colonne $(n,1)$, matrice carrée d'ordre $n$, matrice nulle, matrice unité $I_2 = \begin{pmatrix} 1 & 0 \\ 0 & 1 \end{pmatrix}$.
\item Deux matrices sont égales si et seulement si elles ont le même format et les mêmes coefficients deux à deux.
\item Un vecteur $\vec{u}(x ; y)$ du plan se représente par la matrice colonne $\begin{pmatrix} x \\ y \end{pmatrix}$ : c'est le pont entre géométrie et calcul matriciel.
\end{itemize}
\end{aretenir}

\coursec{Somme de matrices et produit par un réel}

Dans la section précédente, nous avons appris à \emph{présenter} des données sous forme de matrice : un tableau de ventes, un système d'équations, les coefficients d'une application linéaire. Mais une matrice n'est pas qu'un simple tableau figé : c'est un objet mathématique sur lequel on peut \emph{calculer}. Toute cette section repose sur une idée enfantine : \textbf{on fait les opérations case par case}. Si le gérant d'une boutique MTN de Douala enregistre ses ventes de la semaine 1 dans un tableau et celles de la semaine 2 dans un autre tableau de même format, le cumul des deux semaines s'obtient en additionnant les cases qui se correspondent. C'est exactement la somme de deux matrices.

\courssub{Somme de deux matrices de même format}

\textbf{Intuition : le cumul des ventes}\par

Imaginons une boutique qui vend deux produits (cartes SIM et téléphones) dans deux agences (Akwa et Bonabéri). Les ventes de la semaine 1 et de la semaine 2 sont résumées par deux matrices de format $(2,2)$ :

\[ A = \begin{pmatrix} 30 & 12 \\ 25 & 8 \end{pmatrix} \qquad B = \begin{pmatrix} 20 & 5 \\ 15 & 10 \end{pmatrix} \]

Pour connaître le total des deux semaines, on additionne \emph{position par position} : $30+20=50$ cartes SIM à Akwa, $12+5=17$ téléphones à Akwa, etc. Le tableau résultat a le \emph{même format} que les deux tableaux de départ. Cette observation naturelle motive la définition officielle.

\begin{definition}[Somme de deux matrices]
Soient $A = (a_{ij})$ et $B = (b_{ij})$ deux matrices de \cle{même format} $(n,p)$. La \cle{somme} de $A$ et $B$, notée $A+B$, est la matrice de format $(n,p)$ définie par :
\[ A + B = (a_{ij} + b_{ij}) \]
Autrement dit, chaque coefficient de $A+B$ est la somme des coefficients situés \textbf{à la même position} dans $A$ et dans $B$.
\end{definition}

\begin{popfigure}
\begin{tikzpicture}[scale=1.05, every node/.style={font=\small}]
% Matrice A
\draw[popblue, thick] (0,0) node[left=2pt] {$A=$};
\draw[popblue, thick] (0.25,0.55) -- (0.35,0.55) -- (0.35,-0.55) -- (0.25,-0.55);
\draw[popblue, thick] (1.75,0.55) -- (1.65,0.55) -- (1.65,-0.55) -- (1.75,-0.55);
\node[popblue] at (0.75,0.28) (a11) {$a_{11}$};
\node[popblue] at (1.25,0.28) (a12) {$a_{12}$};
\node[popblue] at (0.75,-0.28) (a21) {$a_{21}$};
\node[popblue] at (1.25,-0.28) (a22) {$a_{22}$};
% signe +
\node at (2.35,0) {$+$};
% Matrice B
\draw[poporange, thick] (2.95,0.55) -- (3.05,0.55) -- (3.05,-0.55) -- (2.95,-0.55);
\draw[poporange, thick] (4.45,0.55) -- (4.35,0.55) -- (4.35,-0.55) -- (4.45,-0.55);
\node[poporange] at (3.45,0.28) (b11) {$b_{11}$};
\node[poporange] at (3.95,0.28) (b12) {$b_{12}$};
\node[poporange] at (3.45,-0.28) (b21) {$b_{21}$};
\node[poporange] at (3.95,-0.28) (b22) {$b_{22}$};
% signe =
\node at (5.05,0) {$=$};
% Matrice somme
\draw[popgreen, thick] (5.65,0.55) -- (5.75,0.55) -- (5.75,-0.55) -- (5.65,-0.55);
\draw[popgreen, thick] (7.85,0.55) -- (7.75,0.55) -- (7.75,-0.55) -- (7.85,-0.55);
\node[popgreen] at (6.35,0.28) (s11) {$a_{11}{+}b_{11}$};
\node[popgreen] at (7.15,0.28) (s12) {$a_{12}{+}b_{12}$};
\node[popgreen] at (6.35,-0.28) (s21) {$a_{21}{+}b_{21}$};
\node[popgreen] at (7.15,-0.28) (s22) {$a_{22}{+}b_{22}$};
% flèches
\draw[-{Stealth}, poppurple, thick] (a11.east) .. controls (2.6,0.9) and (5.2,0.9) .. (s11.north);
\draw[-{Stealth}, poppurple, thick] (b11.north) .. controls (4.2,0.9) and (5.6,0.8) .. (s11.north);
\draw[-{Stealth}, poppink, thick] (a22.east) .. controls (2.6,-0.9) and (5.6,-0.9) .. (s22.south);
\draw[-{Stealth}, poppink, thick] (b22.south) .. controls (4.4,-0.9) and (6.2,-0.8) .. (s22.south);
\end{tikzpicture}
\legende{Addition coefficient par coefficient : chaque case du résultat est la somme des deux cases de même position.}
\end{popfigure}

\begin{exemplebox}[Cumul des ventes sur deux semaines]
Reprenons les ventes de la boutique :
\[ A + B = \begin{pmatrix} 30 & 12 \\ 25 & 8 \end{pmatrix} + \begin{pmatrix} 20 & 5 \\ 15 & 10 \end{pmatrix} = \begin{pmatrix} 30+20 & 12+5 \\ 25+15 & 8+10 \end{pmatrix} = \begin{pmatrix} 50 & 17 \\ 40 & 18 \end{pmatrix} \]
En deux semaines, la boutique a vendu 50 cartes SIM et 17 téléphones à Akwa, 40 cartes SIM et 18 téléphones à Bonabéri.
\end{exemplebox}

\begin{attention}[Condition impérative : même format]
On ne peut additionner que deux matrices de \cle{même format} $(n,p)$. Additionner une matrice $(2,2)$ avec une matrice $(2,3)$ n'a \textbf{aucun sens} : il n'y a pas de case correspondante pour la troisième colonne. Avant tout calcul, vérifie toujours les formats : c'est le premier réflexe du candidat au Probatoire.
\end{attention}

\courssub{Produit d'une matrice par un réel}

\textbf{Intuition : la hausse des prix}\par

Supposons que les prix (en centaines de FCFA) de trois articles dans deux points de vente soient donnés par une matrice $P$ de format $(2,3)$. Si tous les prix augmentent de $10\,\%$, chaque prix est multiplié par $1{,}1$ : le nouveau tableau s'obtient en multipliant \emph{chaque case} par $1{,}1$. C'est le produit d'une matrice par un réel.

\begin{definition}[Produit d'une matrice par un réel]
Soit $A = (a_{ij})$ une matrice de format $(n,p)$ et $k$ un nombre réel. Le \cle{produit} de $A$ par le réel $k$, noté $kA$, est la matrice de format $(n,p)$ définie par :
\[ kA = (k\,a_{ij}) \]
Autrement dit, on multiplie \textbf{chaque coefficient} de $A$ par $k$.
\end{definition}

\begin{exemplebox}[Hausse de 10 \% des prix]
Les prix (en FCFA) de trois articles dans deux boutiques sont :
\[ P = \begin{pmatrix} 500 & 1200 & 800 \\ 550 & 1150 & 850 \end{pmatrix} \]
Après une hausse générale de $10\,\%$, les nouveaux prix sont donnés par :
\[ 1{,}1\,P = \begin{pmatrix} 1{,}1\times 500 & 1{,}1\times 1200 & 1{,}1\times 800 \\ 1{,}1\times 550 & 1{,}1\times 1150 & 1{,}1\times 850 \end{pmatrix} = \begin{pmatrix} 550 & 1320 & 880 \\ 605 & 1265 & 935 \end{pmatrix} \]
De même, $2A$ représente un doublement des quantités, et $0{,}5A$ une réduction de moitié.
\end{exemplebox}

\courssub{Propriétés de l'addition et du produit par un réel}

L'addition des matrices se faisant coefficient par coefficient, elle hérite de \emph{toutes} les propriétés de l'addition des nombres réels. C'est une excellente nouvelle : tu calcules avec les matrices comme tu calcules avec les nombres, sans changer tes habitudes.

\begin{propriete}[Propriétés de $(\mathcal{M}_{n,p}, +, \cdot)$ — admises]
Pour toutes matrices $A$, $B$, $C$ de même format $(n,p)$ et tous réels $k$ et $k'$ :
\begin{itemize}
\item \textbf{Commutativité} : $A + B = B + A$ ;
\item \textbf{Associativité} : $(A + B) + C = A + (B + C)$ ;
\item \textbf{Élément neutre} : la \cle{matrice nulle} $O_{n,p}$ (tous ses coefficients sont nuls) vérifie $A + O_{n,p} = A$ ;
\item \textbf{Opposée} : toute matrice $A$ admet une opposée $-A = (-1)A$ vérifiant $A + (-A) = O_{n,p}$ ;
\item \textbf{Distributivités} : $k(A+B) = kA + kB$ et $(k+k')A = kA + k'A$ ;
\item \textbf{Compatibilité} : $k(k'A) = (kk')A$ et $1A = A$.
\end{itemize}
Ces propriétés sont admises : leur vérification est immédiate car tout se ramène, coefficient par coefficient, aux propriétés de l'addition et de la multiplication dans $\mathbb{R}$.
\end{propriete}

\begin{experience}[Vérification sur un cas concret]
Vérifions la distributivité $k(A+B) = kA + kB$ sur un exemple. Prenons $k=2$, $A = \begin{pmatrix} 1 & 3 \\ 0 & -2 \end{pmatrix}$ et $B = \begin{pmatrix} 4 & -1 \\ 2 & 5 \end{pmatrix}$.
\begin{align*}
\text{D'une part : } 2(A+B) &= 2\begin{pmatrix} 5 & 2 \\ 2 & 3 \end{pmatrix} = \begin{pmatrix} 10 & 4 \\ 4 & 6 \end{pmatrix} \\
\text{D'autre part : } 2A + 2B &= \begin{pmatrix} 2 & 6 \\ 0 & -4 \end{pmatrix} + \begin{pmatrix} 8 & -2 \\ 4 & 10 \end{pmatrix} = \begin{pmatrix} 10 & 4 \\ 4 & 6 \end{pmatrix}
\end{align*}
Les deux résultats coïncident. La démonstration générale est identique : le coefficient en position $(i,j)$ de $k(A+B)$ est $k(a_{ij}+b_{ij}) = ka_{ij} + kb_{ij}$, qui est le coefficient de $kA + kB$ à la même position.
\end{experience}

\begin{aretenir}[L'essentiel]
\begin{itemize}
\item $A + B$ : on additionne \textbf{case par case} (même format obligatoire) ;
\item $kA$ : on multiplie \textbf{chaque case} par $k$ ;
\item La différence se définit par $A - B = A + (-B)$ ;
\item Toutes les règles de calcul de $\mathbb{R}$ (commutativité, associativité, distributivité) restent valables.
\end{itemize}
\end{aretenir}

\begin{methode}[Combiner des tableaux de données]
\begin{enumerate}
\item \textbf{Vérifier les formats} : les deux matrices doivent avoir le même nombre de lignes et le même nombre de colonnes ;
\item \textbf{Identifier l'opération} : cumul de données $\rightarrow$ somme $A+B$ ; écart $\rightarrow$ différence $A-B$ ; changement d'échelle, hausse ou baisse en pourcentage $\rightarrow$ produit par un réel ($\times 1{,}1$ pour $+10\,\%$, $\times 0{,}9$ pour $-10\,\%$) ;
\item \textbf{Calculer case par case}, en reportant chaque résultat à la même position ;
\item \textbf{Contrôler} : le résultat a le même format que les matrices de départ.
\end{enumerate}
\end{methode}

\begin{exoresolu}[Calcul de $A+B$ et de $3A - 2B$]
On donne $A = \begin{pmatrix} 2 & -1 \\ 3 & 4 \end{pmatrix}$ et $B = \begin{pmatrix} 1 & 5 \\ -2 & 0 \end{pmatrix}$. Calculer $A + B$ et $3A - 2B$.
\tcblower
\textbf{Calcul de $A+B$} (addition case par case) :
\[ A + B = \begin{pmatrix} 2+1 & -1+5 \\ 3+(-2) & 4+0 \end{pmatrix} = \begin{pmatrix} 3 & 4 \\ 1 & 4 \end{pmatrix} \]
\textbf{Calcul de $3A - 2B$} : on calcule d'abord chaque produit par un réel :
\[ 3A = \begin{pmatrix} 6 & -3 \\ 9 & 12 \end{pmatrix} \qquad 2B = \begin{pmatrix} 2 & 10 \\ -4 & 0 \end{pmatrix} \]
puis on soustrait case par case :
\[ 3A - 2B = \begin{pmatrix} 6-2 & -3-10 \\ 9-(-4) & 12-0 \end{pmatrix} = \begin{pmatrix} 4 & -13 \\ 13 & 12 \end{pmatrix} \]
\end{exoresolu}

\begin{exoresolu}[Stocks d'un distributeur Orange]
Un distributeur possède deux dépôts. Les stocks (en unités) de trois articles y sont :
\[ S_1 = \begin{pmatrix} 120 & 80 & 45 \\ 90 & 60 & 30 \end{pmatrix} \quad \text{(janvier)}, \qquad S_2 = \begin{pmatrix} 100 & 70 & 50 \\ 110 & 40 & 25 \end{pmatrix} \quad \text{(février)}. \]
\begin{enumerate}
\item Déterminer le stock total cumulé sur les deux mois.
\item En prévision des fêtes, le distributeur décide de multiplier tous ses stocks de février par $1{,}5$. Écrire la matrice correspondante.
\end{enumerate}
\tcblower
\begin{enumerate}
\item Le cumul est la somme des deux matrices (même format $(2,3)$) :
\[ S_1 + S_2 = \begin{pmatrix} 220 & 150 & 95 \\ 200 & 100 & 55 \end{pmatrix} \]
\item On multiplie chaque coefficient de $S_2$ par $1{,}5$ :
\[ 1{,}5\,S_2 = \begin{pmatrix} 150 & 105 & 75 \\ 165 & 60 & 37{,}5 \end{pmatrix} \]
\end{enumerate}
\end{exoresolu}

\begin{attention}[Pièges fréquents]
\begin{itemize}
\item \textbf{Formats différents} : $\begin{pmatrix} 1 & 2 \end{pmatrix} + \begin{pmatrix} 1 \\ 2 \end{pmatrix}$ n'existe pas, même si les coefficients sont les mêmes : les formats $(1,2)$ et $(2,1)$ diffèrent.
\item \textbf{Produit par un réel} : dans $kA$, \emph{tous} les coefficients sont multipliés par $k$, pas seulement la première ligne ou la première colonne (erreur classique de précipitation).
\item \textbf{Signes} : dans $A - B$, c'est le coefficient de $B$ qui change de signe : $a_{ij} - b_{ij}$. Attention aux coefficients négatifs : $3 - (-2) = 5$ et non $1$.
\end{itemize}
\end{attention}

\begin{exercice}[Pour t'entraîner]
On donne $A = \begin{pmatrix} 4 & -2 & 1 \\ 0 & 3 & -5 \end{pmatrix}$ et $B = \begin{pmatrix} -1 & 2 & 3 \\ 6 & -4 & 0 \end{pmatrix}$.
\begin{enumerate}
\item Calculer $A + B$, $-2A$ et $2A + 3B$.
\item Déterminer la matrice $X$ de format $(2,3)$ telle que $A + X = B$.
\end{enumerate}
\end{exercice}

\begin{corrige}[Exercice proposé]
\begin{enumerate}
\item $A + B = \begin{pmatrix} 3 & 0 & 4 \\ 6 & -1 & -5 \end{pmatrix}$ ; \quad $-2A = \begin{pmatrix} -8 & 4 & -2 \\ 0 & -6 & 10 \end{pmatrix}$ ;
\[ 2A + 3B = \begin{pmatrix} 8 & -4 & 2 \\ 0 & 6 & -10 \end{pmatrix} + \begin{pmatrix} -3 & 6 & 9 \\ 18 & -12 & 0 \end{pmatrix} = \begin{pmatrix} 5 & 2 & 11 \\ 18 & -6 & -10 \end{pmatrix}. \]
\item $X = B - A = \begin{pmatrix} -5 & 4 & 2 \\ 6 & -7 & 5 \end{pmatrix}$. On vérifie que $A + X = B$ case par case.
\end{enumerate}
\end{corrige}

\begin{savaistu}[Pourquoi ces propriétés sont-elles si importantes ?]
Les propriétés établies dans cette section font de l'ensemble $\mathcal{M}_{n,p}(\mathbb{R})$, muni de l'addition et de la multiplication par un réel, ce que les mathématiciens appellent un \emph{espace vectoriel} — exactement la même structure que le plan vectoriel que tu étudies en géométrie. Une matrice $(2,2)$ peut ainsi être vue comme un « vecteur » à quatre composantes. Cette analogie puissante, développée au XIX\textsuperscript{e} siècle par Cayley et Sylvester, est aujourd'hui au cœur de l'informatique : les images de ton téléphone sont des matrices de pixels, et les filtres (luminosité, contraste) ne sont rien d'autre que des combinaisons $kA + B$ de matrices !
\end{savaistu}

Maintenant que l'addition et la multiplication par un réel n'ont plus de secret pour toi, une question naturelle se pose : peut-on multiplier deux matrices \emph{entre elles} ? La réponse, plus subtile et beaucoup plus riche, fera l'objet de la section suivante — et c'est elle qui donnera tout son sens au lien entre matrices et applications linéaires.

\coursec{Produit de deux matrices}

Dans la section précédente, nous avons appris à additionner deux matrices et à multiplier une matrice par un réel : ces opérations se font \og case par case \fg{} et ressemblent beaucoup au calcul sur les nombres. Nous allons maintenant découvrir l'opération la plus riche, la plus utile — et la plus surprenante — de l'algèbre des matrices : le \cle{produit de deux matrices}. Attention tout de suite : ce produit ne se fait \emph{pas} case par case ! C'est lui qui donnera tout son sens à l'écriture matricielle des systèmes d'équations et, dans la suite de la leçon, aux applications linéaires du plan.

\courssub{Le cas fondateur : produit d'une ligne par une colonne}

Avant de définir le produit général, commençons par une situation concrète. Un vendeur du marché Mokolo vend deux articles : des pagnes à \num{2500} FCFA l'unité et des sacs de riz à \num{12000} FCFA l'unité. Lundi, il vend 4 pagnes et 3 sacs de riz. Sa recette est :
\[ 4 \times 2500 + 3 \times 12000 = 10000 + 36000 = 46000 \text{ FCFA.} \]
On a multiplié une \emph{ligne} de quantités $(4 \quad 3)$ par une \emph{colonne} de prix $\begin{pmatrix} 2500 \\ 12000 \end{pmatrix}$, en associant terme à terme puis en additionnant. C'est exactement cela, le produit ligne $\times$ colonne.

\begin{definition}[Produit d'une matrice ligne par une matrice colonne]
Soit $L = \begin{pmatrix} a_1 & a_2 & \cdots & a_p \end{pmatrix}$ une matrice ligne de format $(1,p)$ et $C = \begin{pmatrix} b_1 \\ b_2 \\ \vdots \\ b_p \end{pmatrix}$ une matrice colonne de format $(p,1)$, ayant \textbf{le même nombre $p$ de coefficients}. Leur produit est le réel :
\[ L \times C = a_1 b_1 + a_2 b_2 + \cdots + a_p b_p = \sum_{k=1}^{p} a_k b_k. \]
\end{definition}

\begin{exemplebox}[Calcul d'un produit ligne $\times$ colonne]
\[ \begin{pmatrix} 1 & 2 & -3 \end{pmatrix} \times \begin{pmatrix} 4 \\ 0 \\ 5 \end{pmatrix} = 1 \times 4 + 2 \times 0 + (-3) \times 5 = 4 - 15 = -11. \]
On avance le long de la ligne en descendant le long de la colonne, en multipliant les coefficients qui se font face, puis on additionne.
\end{exemplebox}

\begin{attention}[Condition indispensable]
Le produit $L \times C$ n'existe que si la ligne et la colonne ont le \textbf{même nombre de coefficients}. On ne peut pas multiplier $\begin{pmatrix} 1 & 2 \end{pmatrix}$ par $\begin{pmatrix} 3 \\ 4 \\ 5 \end{pmatrix}$ : il manque un partenaire au troisième coefficient.
\end{attention}

\courssub{Définition générale du produit de deux matrices}

L'idée est maintenant toute simple : pour multiplier deux matrices $A$ et $B$, on multiplie \emph{chaque ligne} de $A$ par \emph{chaque colonne} de $B$, et le résultat de chacun de ces petits produits donne un coefficient de la matrice produit.

\begin{definition}[Produit de deux matrices]
Soit $A$ une matrice de format $(n,p)$ et $B$ une matrice de format $(p,q)$ : le \textbf{nombre de colonnes de $A$ est égal au nombre de lignes de $B$} (condition de compatibilité). Alors le produit $AB$ est la matrice de format $(n,q)$ dont le coefficient situé à la ligne $i$ et à la colonne $j$ est :
\[ (AB)_{i,j} = \sum_{k=1}^{p} a_{i,k}\, b_{k,j} = a_{i,1}b_{1,j} + a_{i,2}b_{2,j} + \cdots + a_{i,p}b_{p,j}, \]
c'est-à-dire le produit de la \cle{ligne $i$ de $A$} par la \cle{colonne $j$ de $B$}.
\end{definition}

Retenons le format du produit, qui se lit comme une simplification de fractions :
\[ (n,\underline{p}) \times (\underline{p},q) = (n,q). \]
Les deux $p$ doivent coïncider (sinon le produit n'existe pas) et ils \og disparaissent \fg{} dans le résultat.

\begin{popfigure}
\begin{center}
\begin{tikzpicture}[scale=1.05, every node/.style={font=\small}]
% Matrice A (2x3)
\node[font=\bfseries] at (-0.9,0.9) {$A$};
\draw[popblue, thick] (-1.9,0.3) -- (-1.9,1.5);
\draw[popblue, thick] (0.1,0.3) -- (0.1,1.5);
\node at (-1.4,1.2) {$1$}; \node at (-0.9,1.2) {$2$}; \node at (-0.4,1.2) {$0$};
\node at (-1.4,0.6) {$3$}; \node at (-0.9,0.6) {$-1$}; \node at (-0.4,0.6) {$4$};
\node[popmuted] at (-0.9,-0.1) {format $(2,3)$};
% symbole fois
\node[font=\large] at (0.7,0.9) {$\times$};
% Matrice B (3x2)
\node[font=\bfseries] at (2.2,0.9) {$B$};
\draw[poporange, thick] (1.2,0.0) -- (1.2,1.8);
\draw[poporange, thick] (3.2,0.0) -- (3.2,1.8);
\node at (1.7,1.5) {$2$}; \node at (2.7,1.5) {$1$};
\node at (1.7,0.9) {$0$}; \node at (2.7,0.9) {$-1$};
\node at (1.7,0.3) {$5$}; \node at (2.7,0.3) {$3$};
\node[popmuted] at (2.2,-0.1) {format $(3,2)$};
% egal
\node[font=\large] at (3.9,0.9) {$=$};
% Matrice produit (2x2)
\node[font=\bfseries] at (5.3,0.9) {$AB$};
\draw[popgreen, thick] (4.6,0.3) -- (4.6,1.5);
\draw[popgreen, thick] (6.0,0.3) -- (6.0,1.5);
\node at (4.9,1.2) {$?$}; \node at (5.7,1.2) {$?$};
\node at (4.9,0.6) {$?$}; \node at (5.7,0.6) {$?$};
\node[popmuted] at (5.3,-0.1) {format $(2,2)$};
% annotation des formats
\draw[->, popdark, thick] (0.2,-0.35) .. controls (1.5,-0.9) .. (1.2,-0.35);
\node[popdark] at (1.6,-1.0) {les deux $3$ coïncident : le produit existe};
\end{tikzpicture}
\end{center}
\legende{Diagramme des formats : $(2,3)\times(3,2)\to(2,2)$. Le nombre de colonnes de $A$ doit égaler le nombre de lignes de $B$.}
\end{popfigure}

\courssub{La disposition en croix : une technique infaillible}

Pour calculer un produit sans se tromper, les meilleurs élèves utilisent la \cle{disposition en croix} : on place $B$ \emph{au-dessus} de l'emplacement du résultat, et $A$ \emph{à gauche}. Chaque coefficient de $AB$ se trouve alors à l'intersection de la ligne de $A$ et de la colonne de $B$ qu'on lit directement.

\begin{methode}[Calculer $AB$ avec la disposition en croix]
\begin{enumerate}
\item Vérifier la compatibilité des formats : $(n,p)\times(p,q)$.
\item Placer $B$ en haut, $A$ à gauche, et préparer une matrice vide $(n,q)$ dans le coin inférieur droit.
\item Pour obtenir le coefficient $(i,j)$ : suivre la ligne $i$ de $A$ et la colonne $j$ de $B$, multiplier les termes qui se font face, additionner.
\item Remplir toutes les cases, puis vérifier un coefficient au hasard.
\end{enumerate}
\end{methode}

\begin{popfigure}
\begin{center}
\begin{tikzpicture}[scale=1.0, every node/.style={font=\small}]
% B au-dessus
\node[font=\bfseries, poporange] at (3.3,3.6) {$B$};
\draw[poporange, thick] (2.3,2.1) -- (2.3,3.3);
\draw[poporange, thick] (4.3,2.1) -- (4.3,3.3);
\node at (2.8,3.0) {$2$}; \node[popink, font=\bfseries] at (3.8,3.0) {$1$};
\node at (2.8,2.4) {$-1$}; \node[popink, font=\bfseries] at (3.8,2.4) {$4$};
% A à gauche
\node[font=\bfseries, popblue] at (-0.5,1.35) {$A$};
\draw[popblue, thick] (-0.1,0.75) -- (-0.1,1.95);
\draw[popblue, thick] (1.5,0.75) -- (1.5,1.95);
\node[popink, font=\bfseries] at (0.35,1.65) {$1$}; \node[popink, font=\bfseries] at (1.05,1.65) {$2$};
\node at (0.35,1.05) {$0$}; \node at (1.05,1.05) {$3$};
% résultat
\node[font=\bfseries, popgreen] at (3.3,0.35) {$AB$};
\draw[popgreen, thick] (2.3,0.75) -- (2.3,1.95);
\draw[popgreen, thick] (4.3,0.75) -- (4.3,1.95);
\node at (2.8,1.65) {$0$}; \node[popink, font=\bfseries] at (3.8,1.65) {$9$};
\node at (2.8,1.05) {$-3$}; \node at (3.8,1.05) {$12$};
% flèches de lecture pour le coefficient (1,2)
\draw[->, popink, thick] (1.6,1.65) -- (3.55,1.65);
\draw[->, popink, thick] (3.8,2.05) -- (3.8,1.9);
\node[popink, align=center] at (6.3,1.7) {coefficient $(1,2)$ :\\ $1\times 1 + 2\times 4 = 9$};
\draw[->, popink] (5.0,1.65) -- (4.35,1.65);
\end{tikzpicture}
\end{center}
\legende{Disposition en croix : $A$ à gauche, $B$ au-dessus. Le coefficient $(1,2)$ de $AB$ se lit à la croisée de la ligne 1 de $A$ et de la colonne 2 de $B$.}
\end{popfigure}

\begin{exemplebox}[Calcul complet de $AB$ avec $A=\begin{pmatrix}1&2\\0&3\end{pmatrix}$ et $B=\begin{pmatrix}2&1\\-1&4\end{pmatrix}$]
Les formats sont $(2,2)\times(2,2)$ : le produit existe et est de format $(2,2)$.
\begin{align*}
(AB)_{1,1} &= 1\times 2 + 2\times(-1) = 2-2 = 0, &
(AB)_{1,2} &= 1\times 1 + 2\times 4 = 1+8 = 9,\\
(AB)_{2,1} &= 0\times 2 + 3\times(-1) = -3, &
(AB)_{2,2} &= 0\times 1 + 3\times 4 = 12.
\end{align*}
Donc $AB = \begin{pmatrix} 0 & 9 \\ -3 & 12 \end{pmatrix}$.
\end{exemplebox}

\courssub{Le produit matriciel n'est pas commutatif}

Voici la grande surprise. Avec les nombres réels, $3\times 5 = 5\times 3$ : on dit que la multiplication est \emph{commutative}. Avec les matrices, c'est \textbf{faux en général} !

Reprenons nos deux matrices et calculons $BA$ cette fois :
\begin{align*}
(BA)_{1,1} &= 2\times 1 + 1\times 0 = 2, &
(BA)_{1,2} &= 2\times 2 + 1\times 3 = 7,\\
(BA)_{2,1} &= -1\times 1 + 4\times 0 = -1, &
(BA)_{2,2} &= -1\times 2 + 4\times 3 = 10.
\end{align*}
Donc $BA = \begin{pmatrix} 2 & 7 \\ -1 & 10 \end{pmatrix}$, alors que $AB = \begin{pmatrix} 0 & 9 \\ -3 & 12 \end{pmatrix}$. On constate que :
\[ AB \neq BA. \]

\begin{attention}[Erreur fréquente : croire que $AB = BA$]
Le produit matriciel \textbf{n'est pas commutatif} : en général $AB \neq BA$. Pire, il arrive que $AB$ existe et que $BA$ n'existe même pas ! Par exemple, si $A$ est de format $(2,3)$ et $B$ de format $(3,4)$, alors $AB$ existe (format $(2,4)$), mais $BA$ serait de type $(3,4)\times(2,3)$ : les formats $4$ et $2$ ne coïncident pas, le produit $BA$ n'a aucun sens. Conséquences pratiques :
\begin{itemize}
\item on ne change jamais l'ordre des facteurs dans un calcul matriciel ;
\item on précise \og multiplier à gauche \fg{} ou \og multiplier à droite \fg{} ;
\item les identités remarquables habituelles sont fausses : $(A+B)^2 = A^2 + AB + BA + B^2$, et non $A^2 + 2AB + B^2$.
\end{itemize}
\end{attention}

\begin{savaistu}[La non-commutativité au cœur de la physique moderne]
En 1925, le jeune physicien allemand Werner \cle{Heisenberg} découvre que les grandeurs physiques d'un atome (position, vitesse, énergie) doivent être représentées par des objets mathématiques dont le produit n'est pas commutatif... ce sont exactement des matrices ! De la relation $AB \neq BA$ naît le célèbre \emph{principe d'incertitude} : on ne peut pas connaître simultanément avec une précision infinie la position et la vitesse d'une particule. Cette révolution, la \cle{mécanique quantique}, est aujourd'hui à la base du laser, de l'imagerie médicale et des puces de nos téléphones. Heisenberg reçut le prix Nobel de physique en 1932.
\end{savaistu}

\courssub{Propriétés du produit matriciel}

Si la commutativité est perdue, heureusement les autres bonnes propriétés sont conservées.

\begin{propriete}[Associativité et distributivité du produit matriciel (admise)]
Pour toutes matrices $A$, $B$, $C$ de formats compatibles, et pour tout réel $k$ :
\begin{itemize}
\item \textbf{Associativité} : $(AB)C = A(BC)$ ; on peut donc écrire $ABC$ sans parenthèses ;
\item \textbf{Distributivité à gauche} : $A(B+C) = AB + AC$ ;
\item \textbf{Distributivité à droite} : $(A+B)C = AC + BC$ ;
\item \textbf{Compatibilité avec le produit par un réel} : $k(AB) = (kA)B = A(kB)$ ;
\item \textbf{Élément neutre} : si $A$ est carrée d'ordre $n$ et $I_n$ la matrice identité d'ordre $n$, alors
\[ A\,I_n = I_n\,A = A. \]
\end{itemize}
Ces résultats sont admis ; au Probatoire, on sait les \emph{utiliser} dans les calculs.
\end{propriete}

\begin{exemplebox}[Vérification de $AI_2 = A$]
Avec $A = \begin{pmatrix} 1 & 2 \\ 0 & 3 \end{pmatrix}$ et $I_2 = \begin{pmatrix} 1 & 0 \\ 0 & 1 \end{pmatrix}$ :
\[ A I_2 = \begin{pmatrix} 1\times 1 + 2\times 0 & 1\times 0 + 2\times 1 \\ 0\times 1 + 3\times 0 & 0\times 0 + 3\times 1 \end{pmatrix} = \begin{pmatrix} 1 & 2 \\ 0 & 3 \end{pmatrix} = A. \]
La matrice identité joue pour le produit le rôle que joue le nombre $1$ pour la multiplication des réels.
\end{exemplebox}

\courssub{Lien avec les systèmes d'équations dans $\mathbb{R}^2$}

Voici enfin la raison profonde de cette définition si particulière du produit. Considérons le système de deux équations à deux inconnues :
\[ \begin{cases} ax + by = c \\ a'x + b'y = c' \end{cases} \]
Posons $A = \begin{pmatrix} a & b \\ a' & b' \end{pmatrix}$ (matrice des coefficients), $X = \begin{pmatrix} x \\ y \end{pmatrix}$ (colonne des inconnues) et $B = \begin{pmatrix} c \\ c' \end{pmatrix}$ (colonne des seconds membres). Calculons le produit $AX$ :
\[ AX = \begin{pmatrix} a & b \\ a' & b' \end{pmatrix} \begin{pmatrix} x \\ y \end{pmatrix} = \begin{pmatrix} ax + by \\ a'x + b'y \end{pmatrix}. \]
Dire que $AX = B$, c'est dire exactement que $ax+by = c$ et $a'x+b'y = c'$. Le système s'écrit donc sous la forme compacte :
\[ \boxed{AX = B} \]
qui ressemble à l'équation du premier degré $ax = b$ ! Cette analogie est précieuse : de même qu'on résout $ax=b$ en multipliant par l'inverse de $a$, on résoudra $AX = B$ en multipliant par l'\emph{inverse} de la matrice $A$ — c'est l'objet de la prochaine section.

\begin{exoresolu}[Écriture matricielle puis résolution d'un système]
Un cybercafé de Yaoundé facture deux services. On sait que 2 h de connexion et 3 impressions coûtent 800 FCFA, et que 1 h de connexion et 5 impressions coûtent 750 FCFA. On note $x$ le prix d'une heure de connexion et $y$ celui d'une impression. Écrire le système sous forme matricielle $AX = B$, puis le résoudre.
\tcblower
Le système est $\begin{cases} 2x + 3y = 800 \\ x + 5y = 750 \end{cases}$. Il s'écrit matriciellement :
\[ \underbrace{\begin{pmatrix} 2 & 3 \\ 1 & 5 \end{pmatrix}}_{A} \underbrace{\begin{pmatrix} x \\ y \end{pmatrix}}_{X} = \underbrace{\begin{pmatrix} 800 \\ 750 \end{pmatrix}}_{B}. \]
Résolvons par combinaison. De la deuxième équation, $x = 750 - 5y$. En reportant dans la première :
\[ 2(750 - 5y) + 3y = 800 \iff 1500 - 10y + 3y = 800 \iff -7y = -700 \iff y = 100. \]
Puis $x = 750 - 5\times 100 = 250$. Vérification par le produit matriciel :
\[ AX = \begin{pmatrix} 2 & 3 \\ 1 & 5 \end{pmatrix}\begin{pmatrix} 250 \\ 100 \end{pmatrix} = \begin{pmatrix} 2\times 250 + 3\times 100 \\ 1\times 250 + 5\times 100 \end{pmatrix} = \begin{pmatrix} 800 \\ 750 \end{pmatrix} = B. \]
L'heure de connexion coûte 250 FCFA et l'impression 100 FCFA.
\end{exoresolu}

\courssub{Puissances d'une matrice carrée}

Lorsque $A$ est une matrice \emph{carrée}, le produit $AA$ existe toujours : on peut donc définir ses puissances.

\begin{definition}[Puissances d'une matrice carrée]
Soit $A$ une matrice carrée d'ordre $n$. On pose :
\[ A^2 = A \times A, \qquad A^3 = A^2 \times A = A \times A \times A, \qquad \text{et plus généralement } A^{k+1} = A^k \times A. \]
Par convention, $A^0 = I_n$ et $A^1 = A$.
\end{definition}

\begin{exemplebox}[Calcul de $A^2$ et $A^3$ pour $A = \begin{pmatrix} 1 & 2 \\ 0 & 3 \end{pmatrix}$]
\[ A^2 = \begin{pmatrix} 1 & 2 \\ 0 & 3 \end{pmatrix}\begin{pmatrix} 1 & 2 \\ 0 & 3 \end{pmatrix} = \begin{pmatrix} 1\times 1 + 2\times 0 & 1\times 2 + 2\times 3 \\ 0\times 1 + 3\times 0 & 0\times 2 + 3\times 3 \end{pmatrix} = \begin{pmatrix} 1 & 8 \\ 0 & 9 \end{pmatrix}. \]
\[ A^3 = A^2 \times A = \begin{pmatrix} 1 & 8 \\ 0 & 9 \end{pmatrix}\begin{pmatrix} 1 & 2 \\ 0 & 3 \end{pmatrix} = \begin{pmatrix} 1 & 2+24 \\ 0 & 27 \end{pmatrix} = \begin{pmatrix} 1 & 26 \\ 0 & 27 \end{pmatrix}. \]
Remarque : ici les coefficients diagonaux de $A^k$ sont $1$ et $3^k$ ($1$, $3$, $9$, $27$...) : ce genre d'observation permet, en Terminale, de conjecturer une formule générale pour $A^k$ et de la prouver par récurrence.
\end{exemplebox}

\begin{attention}[Deux pièges avec les puissances]
\begin{itemize}
\item Les puissances ne sont définies que pour les matrices \textbf{carrées} : si $A$ est de format $(2,3)$, le produit $A\times A$ n'existe pas.
\item En général $(AB)^2 \neq A^2 B^2$, car $(AB)^2 = ABAB$ et on ne peut pas permuter $B$ et $A$.
\end{itemize}
\end{attention}

\begin{exercice}[Pour s'entraîner]
Soit $M = \begin{pmatrix} 2 & -1 \\ 1 & 0 \end{pmatrix}$ et $N = \begin{pmatrix} 1 & 3 \\ 2 & 1 \end{pmatrix}$.
\begin{enumerate}
\item Calculer $MN$ et $NM$. Que constate-t-on ?
\item Calculer $M^2$, puis vérifier que $M^2 = 2M - I_2$.
\item Écrire sous forme matricielle le système $\begin{cases} 2x - y = 5 \\ x = 3 \end{cases}$ et le résoudre.
\end{enumerate}
\end{exercice}

\begin{corrige}[Exercice 1]
\begin{enumerate}
\item $MN = \begin{pmatrix} 2\times 1 + (-1)\times 2 & 2\times 3 + (-1)\times 1 \\ 1\times 1 + 0\times 2 & 1\times 3 + 0\times 1 \end{pmatrix} = \begin{pmatrix} 0 & 5 \\ 1 & 3 \end{pmatrix}$ et $NM = \begin{pmatrix} 1\times 2 + 3\times 1 & 1\times(-1) + 3\times 0 \\ 2\times 2 + 1\times 1 & 2\times(-1) + 1\times 0 \end{pmatrix} = \begin{pmatrix} 5 & -1 \\ 5 & -2 \end{pmatrix}$. On constate que $MN \neq NM$ : le produit n'est pas commutatif.
\item $M^2 = \begin{pmatrix} 2 & -1 \\ 1 & 0 \end{pmatrix}\begin{pmatrix} 2 & -1 \\ 1 & 0 \end{pmatrix} = \begin{pmatrix} 4-1 & -2+0 \\ 2+0 & -1+0 \end{pmatrix} = \begin{pmatrix} 3 & -2 \\ 2 & -1 \end{pmatrix}$. D'autre part $2M - I_2 = \begin{pmatrix} 4 & -2 \\ 2 & 0 \end{pmatrix} - \begin{pmatrix} 1 & 0 \\ 0 & 1 \end{pmatrix} = \begin{pmatrix} 3 & -2 \\ 2 & -1 \end{pmatrix} = M^2$. L'égalité est vérifiée.
\item Le système s'écrit $\begin{pmatrix} 2 & -1 \\ 1 & 0 \end{pmatrix}\begin{pmatrix} x \\ y \end{pmatrix} = \begin{pmatrix} 5 \\ 3 \end{pmatrix}$. La deuxième ligne donne $x = 3$ ; la première donne $2\times 3 - y = 5$, soit $y = 1$. Solution : $(x\,;\,y) = (3\,;\,1)$. Vérification : $\begin{pmatrix} 2 & -1 \\ 1 & 0 \end{pmatrix}\begin{pmatrix} 3 \\ 1 \end{pmatrix} = \begin{pmatrix} 5 \\ 3 \end{pmatrix}$. 
\end{enumerate}
\end{corrige}

\begin{aretenir}[L'essentiel du produit de matrices]
\begin{itemize}
\item Le produit $AB$ existe si et seulement si le \cle{nombre de colonnes de $A$} égale le \cle{nombre de lignes de $B$} : $(n,p)\times(p,q) = (n,q)$.
\item Le coefficient $(i,j)$ de $AB$ est le produit de la ligne $i$ de $A$ par la colonne $j$ de $B$ : $\sum_{k=1}^{p} a_{i,k}b_{k,j}$. La disposition en croix évite les erreurs.
\item Le produit matriciel \textbf{n'est pas commutatif} : $AB \neq BA$ en général, et $BA$ peut même ne pas exister.
\item Il est \textbf{associatif} et \textbf{distributif} sur l'addition ; la matrice identité vérifie $AI = IA = A$.
\item Tout système $\begin{cases} ax+by=c \\ a'x+b'y=c' \end{cases}$ s'écrit $AX = B$ : c'est la clé de la résolution matricielle qui vient.
\item Pour $A$ carrée : $A^2 = A\times A$, $A^3 = A^2 \times A$, et $A^0 = I_n$.
\end{itemize}
\end{aretenir}

\coursec{Déterminant et inverse d'une matrice carrée d'ordre 2}

Dans la section précédente, tu as appris à multiplier deux matrices et tu as constaté que le produit n'est pas commutatif, mais qu'il est associatif. Une question naturelle se pose alors : si la multiplication des matrices ressemble à la multiplication des nombres, peut-on aussi \emph{diviser} ? Autrement dit, une matrice $A$ admet-elle une \og inverse \fg{} $A^{-1}$ telle que $A \times A^{-1} = I_2$ ? C'est exactement le problème que résout cette section, et la clé de tout le chapitre est un petit nombre, le \cle{déterminant}, qui joue pour les matrices le rôle que joue le test \og $a \neq 0$ \fg{} pour les nombres réels.

\courssub{Le déterminant : un nombre qui résume une matrice}

\textbf{Intuition.} Considère le système $\begin{cases} ax + by = e \\ cx + dy = f \end{cases}$. En éliminant $y$ (on multiplie la première ligne par $d$, la seconde par $b$, et on soustrait), on obtient $(ad - bc)\,x = de - bf$. Tout dépend donc du nombre $ad - bc$ : s'il est non nul, le système a une solution unique ; s'il est nul, tout peut arriver. Ce nombre $ad - bc$ mérite un nom.

\begin{definition}[Déterminant d'une matrice carrée d'ordre 2]
Soit $A = \begin{pmatrix} a & b \\ c & d \end{pmatrix}$ une matrice carrée d'ordre 2. On appelle \cle{déterminant} de $A$ le nombre réel
\[ \det(A) = ad - bc. \]
On le note aussi $\begin{vmatrix} a & b \\ c & d \end{vmatrix} = ad - bc$.
\end{definition}

\begin{attention}[Barres verticales ou crochets ?]
La notation $\begin{vmatrix} a & b \\ c & d \end{vmatrix}$ désigne un \textbf{nombre} (le déterminant), tandis que $\begin{pmatrix} a & b \\ c & d \end{pmatrix}$ désigne un \textbf{tableau} (la matrice). Ne confonds jamais les deux : écrire $\begin{vmatrix} 3 & 1 \\ 5 & 2 \end{vmatrix} = \begin{pmatrix} 1 \end{pmatrix}$ serait un contresens.
\end{attention}

Pour retenir la formule, visualise les deux \og diagonales \fg{} : on fait le produit de la diagonale descendante ($a \times d$) \emph{moins} le produit de la diagonale montante ($b \times c$).

\begin{popfigure}
\begin{center}
\begin{tikzpicture}[scale=1.1, every node/.style={font=\large}]
  \node (a) at (0,1.2) {$a$};
  \node (b) at (1.6,1.2) {$b$};
  \node (c) at (0,0) {$c$};
  \node (d) at (1.6,0) {$d$};
  \draw[popblue, very thick, -{Stealth}] (a.south east) -- (d.north west);
  \draw[poporange, very thick, -{Stealth}] (c.north east) -- (b.south west);
  \node[popblue, font=\bfseries] at (3.4,1.2) {$+\,ad$};
  \node[poporange, font=\bfseries] at (3.4,0) {$-\,bc$};
  \node[popdark, font=\bfseries] at (5.6,0.6) {$\det(A)=ad-bc$};
\end{tikzpicture}
\end{center}
\legende{Schéma mnémotechnique : produit de la diagonale descendante moins produit de la diagonale montante.}
\end{popfigure}

\begin{exemplebox}[Premiers calculs de déterminants]
\begin{itemize}
\item $A = \begin{pmatrix} 3 & 1 \\ 5 & 2 \end{pmatrix}$ : $\det(A) = 3\times 2 - 1\times 5 = 6 - 5 = 1$.
\item $B = \begin{pmatrix} 2 & 4 \\ 1 & 2 \end{pmatrix}$ : $\det(B) = 2\times 2 - 4\times 1 = 4 - 4 = 0$.
\item $I_2 = \begin{pmatrix} 1 & 0 \\ 0 & 1 \end{pmatrix}$ : $\det(I_2) = 1\times 1 - 0\times 0 = 1$.
\end{itemize}
\end{exemplebox}

\courssub{Inversibilité : le théorème fondamental}

\begin{definition}[Matrice inversible]
Une matrice carrée $A$ d'ordre 2 est dite \cle{inversible} s'il existe une matrice $A'$, appelée \cle{inverse} de $A$ et notée $A^{-1}$, telle que
\[ A \times A^{-1} = A^{-1} \times A = I_2. \]
\end{definition}

\begin{propriete}[Critère d'inversibilité et formule de l'inverse — démonstration exigible]
Soit $A = \begin{pmatrix} a & b \\ c & d \end{pmatrix}$.
\begin{itemize}
\item $A$ est inversible \textbf{si et seulement si} $\det(A) \neq 0$ ;
\item dans ce cas, son inverse est unique et donnée par
\[ A^{-1} = \frac{1}{\det(A)} \begin{pmatrix} d & -b \\ -c & a \end{pmatrix} = \frac{1}{ad-bc}\begin{pmatrix} d & -b \\ -c & a \end{pmatrix}. \]
\end{itemize}
\end{propriete}

Retenir la formule est facile : on \textbf{échange} les termes de la diagonale principale ($a$ et $d$), on \textbf{change le signe} des deux autres ($b$ et $c$), et on \textbf{divise} par le déterminant. La matrice $\begin{pmatrix} d & -b \\ -c & a \end{pmatrix}$ s'appelle la \textbf{comatrice} de $A$ (ou matrice adjointe).

\begin{experience}[Démonstration guidée : vérifier que $A \times A^{-1} = I_2$]
Posons $\Delta = ad - bc$ et supposons $\Delta \neq 0$. Notons $B = \dfrac{1}{\Delta}\begin{pmatrix} d & -b \\ -c & a \end{pmatrix}$ et calculons le produit $A \times B$ pas à pas.
\begin{align*}
A \times B &= \frac{1}{\Delta}\begin{pmatrix} a & b \\ c & d \end{pmatrix}\begin{pmatrix} d & -b \\ -c & a \end{pmatrix} \\
&= \frac{1}{\Delta}\begin{pmatrix} ad - bc & -ab + ba \\ cd - dc & -bc + da \end{pmatrix} \\
&= \frac{1}{\Delta}\begin{pmatrix} \Delta & 0 \\ 0 & \Delta \end{pmatrix} = \begin{pmatrix} 1 & 0 \\ 0 & 1 \end{pmatrix} = I_2.
\end{align*}
Un calcul identique montre que $B \times A = I_2$. Donc $B$ est bien l'inverse de $A$.

Réciproquement, si $\Delta = 0$, montrons qu'aucune matrice inverse n'existe. On a $A \times \begin{pmatrix} d & -b \\ -c & a \end{pmatrix} = \begin{pmatrix} \Delta & 0 \\ 0 & \Delta \end{pmatrix} = \begin{pmatrix} 0 & 0 \\ 0 & 0 \end{pmatrix}$ (matrice nulle). Si $A$ avait un inverse $A^{-1}$, en multipliant cette égalité à gauche par $A^{-1}$, on obtiendrait $\begin{pmatrix} d & -b \\ -c & a \end{pmatrix} = A^{-1} \times (\text{matrice nulle}) = \text{matrice nulle}$, donc $a=b=c=d=0$, ce qui signifie $A = 0_2$. Or la matrice nulle n'est pas inversible. Par contraposée, si $\det(A)=0$, $A$ n'est pas inversible. Conclusion : $A$ inversible $\iff \det(A) \neq 0$.
\end{experience}

\begin{attention}[L'erreur classique du Probatoire]
L'erreur la plus fréquente est d'écrire $A^{-1} = \begin{pmatrix} d & -b \\ -c & a \end{pmatrix}$ \textbf{en oubliant de diviser par le déterminant}. La matrice $\begin{pmatrix} d & -b \\ -c & a \end{pmatrix}$ seule n'est \emph{pas} l'inverse : c'est l'inverse \emph{multipliée par} $\det(A)$. Vérifie toujours ton résultat en calculant $A \times A^{-1}$ : tu dois obtenir $I_2$.
\end{attention}

\begin{exoresolu}[Calculer un inverse et vérifier]
Énoncé. Soit $A = \begin{pmatrix} 3 & 1 \\ 5 & 2 \end{pmatrix}$. Montrer que $A$ est inversible, déterminer $A^{-1}$ et vérifier.
\tcblower
\textbf{Solution.} On calcule le déterminant :
\[ \det(A) = 3\times 2 - 1\times 5 = 6 - 5 = 1 \neq 0, \]
donc $A$ est inversible. On échange $3$ et $2$, on change les signes de $1$ et $5$, et on divise par $1$ :
\[ A^{-1} = \frac{1}{1}\begin{pmatrix} 2 & -1 \\ -5 & 3 \end{pmatrix} = \begin{pmatrix} 2 & -1 \\ -5 & 3 \end{pmatrix}. \]
Vérification :
\[ A \times A^{-1} = \begin{pmatrix} 3 & 1 \\ 5 & 2 \end{pmatrix}\begin{pmatrix} 2 & -1 \\ -5 & 3 \end{pmatrix} = \begin{pmatrix} 6-5 & -3+3 \\ 10-10 & -5+6 \end{pmatrix} = \begin{pmatrix} 1 & 0 \\ 0 & 1 \end{pmatrix} = I_2. \quad\checkmark \]
\end{exoresolu}

\begin{propriete}[Déterminant d'un produit — admise]
Pour toutes matrices carrées $A$ et $B$ d'ordre 2 :
\[ \det(A \times B) = \det(A) \times \det(B). \]
En particulier, si $A$ est inversible, $\det(A^{-1}) = \dfrac{1}{\det(A)}$.
\end{propriete}

\begin{exemplebox}[Vérification sur un exemple]
Prenons $A = \begin{pmatrix} 3 & 1 \\ 5 & 2 \end{pmatrix}$ et $B = \begin{pmatrix} 1 & 2 \\ 0 & 3 \end{pmatrix}$. On a $\det(A) = 1$ et $\det(B) = 1\times 3 - 2\times 0 = 3$.
\[ A\times B = \begin{pmatrix} 3 & 6+3 \\ 5 & 10+6 \end{pmatrix} = \begin{pmatrix} 3 & 9 \\ 5 & 16 \end{pmatrix}, \qquad \det(A\times B) = 48 - 45 = 3 = \det(A)\times\det(B). \quad\checkmark \]
\end{exemplebox}

\courssub{Application : résolution des systèmes de deux équations dans $\mathbb{R}^2$}

Voici l'application la plus utile au Probatoire. Le système $\begin{cases} ax + by = e \\ cx + dy = f \end{cases}$ s'écrit matriciellement $A X = B$ avec
\[ A = \begin{pmatrix} a & b \\ c & d \end{pmatrix}, \qquad X = \begin{pmatrix} x \\ y \end{pmatrix}, \qquad B = \begin{pmatrix} e \\ f \end{pmatrix}. \]
Si $\det(A) \neq 0$, on multiplie à gauche par $A^{-1}$ : $A^{-1}AX = A^{-1}B$, c'est-à-dire
\[ AX = B \iff X = A^{-1}B. \]
Le système admet alors une \textbf{unique} solution.

\begin{methode}[Résoudre un système 2$\times$2 par inversion de matrice]
\begin{enumerate}
\item Écrire le système sous la forme $AX = B$ (attention à l'ordre des inconnues et aux signes des coefficients).
\item Calculer $\det(A)$. Si $\det(A) = 0$, la méthode ne s'applique pas (voir plus bas).
\item Calculer $A^{-1} = \dfrac{1}{\det(A)}\begin{pmatrix} d & -b \\ -c & a \end{pmatrix}$.
\item Calculer le produit $X = A^{-1}B$ : la première ligne donne $x$, la deuxième donne $y$.
\item Vérifier en reportant les valeurs trouvées dans les deux équations de départ.
\end{enumerate}
\end{methode}

\begin{exoresolu}[Résolution complète d'un système]
Énoncé. Résoudre dans $\mathbb{R}^2$ le système $\begin{cases} 3x + y = 7 \\ 5x + 2y = 12 \end{cases}$.
\tcblower
\textbf{Solution.} Le système s'écrit $AX = B$ avec $A = \begin{pmatrix} 3 & 1 \\ 5 & 2 \end{pmatrix}$, $X = \begin{pmatrix} x \\ y \end{pmatrix}$, $B = \begin{pmatrix} 7 \\ 12 \end{pmatrix}$.

On a vu que $\det(A) = 1 \neq 0$ et $A^{-1} = \begin{pmatrix} 2 & -1 \\ -5 & 3 \end{pmatrix}$. Donc
\[ X = A^{-1}B = \begin{pmatrix} 2 & -1 \\ -5 & 3 \end{pmatrix}\begin{pmatrix} 7 \\ 12 \end{pmatrix} = \begin{pmatrix} 14 - 12 \\ -35 + 36 \end{pmatrix} = \begin{pmatrix} 2 \\ 1 \end{pmatrix}. \]
La solution unique est $(x ; y) = (2 ; 1)$. Vérification : $3\times 2 + 1 = 7$ \checkmark{} et $5\times 2 + 2\times 1 = 12$ \checkmark.
\end{exoresolu}

\textbf{Et si $\det(A) = 0$ ?} La matrice n'a pas d'inverse et le système n'a pas de solution unique. Deux cas se présentent, que l'on distingue en comparant les équations :
\begin{itemize}
\item \textbf{Infinité de solutions} : les deux équations sont proportionnelles (elles représentent la même droite). Exemple : $\begin{cases} x + 2y = 3 \\ 2x + 4y = 6 \end{cases}$ : la deuxième équation est le double de la première ; tous les couples $(3 - 2t\,;\,t)$ avec $t \in \mathbb{R}$ conviennent.
\item \textbf{Aucune solution} : les équations sont proportionnelles pour les coefficients mais pas pour les seconds membres (droites parallèles distinctes). Exemple : $\begin{cases} x + 2y = 3 \\ 2x + 4y = 7 \end{cases}$ : impossible, car le membre de gauche de la deuxième équation vaut le double du premier, mais $7 \neq 6$.
\end{itemize}

\begin{attention}[Ne pas conclure trop vite]
$\det(A) = 0$ ne signifie pas automatiquement \og pas de solution \fg{} : cela peut aussi vouloir dire \og une infinité de solutions \fg{}. Il faut examiner les équations pour trancher.
\end{attention}

\begin{savaistu}[La règle de Cramer]
La règle de Cramer, que tu as peut-être rencontrée, donne directement
\[ x = \frac{\begin{vmatrix} e & b \\ f & d \end{vmatrix}}{\begin{vmatrix} a & b \\ c & d \end{vmatrix}}, \qquad y = \frac{\begin{vmatrix} a & e \\ c & f \end{vmatrix}}{\begin{vmatrix} a & b \\ c & d \end{vmatrix}}. \]
Ce n'est rien d'autre que le calcul de $X = A^{-1}B$ écrit composante par composante ! Vérifie-le sur l'exemple précédent : $x = \dfrac{7\times 2 - 1\times 12}{1} = 2$ et $y = \dfrac{3\times 12 - 7\times 5}{1} = 1$. Le mathématicien suisse Gabriel Cramer publia cette règle en 1750, bien avant que le mot \og matrice \fg{} n'existe.
\end{savaistu}

\courssub{Complément pour la Terminale : le déterminant mesure une aire}

\begin{savaistu}[Interprétation géométrique du déterminant]
Si $\vec{u}\begin{pmatrix} a \\ c \end{pmatrix}$ et $\vec{v}\begin{pmatrix} b \\ d \end{pmatrix}$ sont les deux vecteurs colonnes de $A$, alors $\left|\det(A)\right|$ est l'\textbf{aire du parallélogramme} construit sur $\vec{u}$ et $\vec{v}$. Ainsi $\det(A) = 0$ signifie géométriquement que le parallélogramme est aplati : les vecteurs sont colinéaires. C'est pourquoi une matrice de déterminant nul \og écrase \fg{} le plan sur une droite et ne peut pas être inversée. Tu retrouveras cette idée en Terminale avec le produit mixte et les transformations du plan.
\end{savaistu}

\begin{popfigure}
\begin{center}
\begin{tikzpicture}[scale=1.3, >=Stealth]
  \coordinate (O) at (0,0);
  \coordinate (U) at (2.6,0.6);
  \coordinate (V) at (0.8,1.8);
  \coordinate (W) at (3.4,2.4);
  \fill[popblueL, opacity=0.55] (O) -- (U) -- (W) -- (V) -- cycle;
  \draw[popblue, thick] (O) -- (U) -- (W) -- (V) -- cycle;
  \draw[->, very thick, popdark] (O) -- (U) node[midway, below, align=center] {$\vec{u}\begin{pmatrix} a \\ c \end{pmatrix}$};
  \draw[->, very thick, poporange] (O) -- (V) node[midway, left, align=center] {$\vec{v}\begin{pmatrix} b \\ d \end{pmatrix}$};
  \node[popdark, font=\bfseries] at (1.7,1.15) {$\mathcal{A} = |ad - bc|$};
  \fill[popdark] (O) circle (1.5pt) node[below left] {$O$};
\end{tikzpicture}
\end{center}
\legende{L'aire du parallélogramme construit sur les colonnes de $A$ vaut $|\det(A)|$.}
\end{popfigure}

\begin{exercice}[Pour t'entraîner]
\begin{enumerate}
\item Calculer le déterminant de $M = \begin{pmatrix} 4 & 3 \\ 2 & 5 \end{pmatrix}$ et de $N = \begin{pmatrix} 6 & 3 \\ 4 & 2 \end{pmatrix}$. Lesquelles sont inversibles ?
\item Déterminer $M^{-1}$ et vérifier que $M \times M^{-1} = I_2$.
\item Un vendeur du marché Mokolo vend des mangues et des avocats. $x$ mangues et $y$ avocats coûtent : $4x + 3y = 2900$ FCFA et $2x + 5y = 2700$ FCFA. Résoudre par inversion de matrice pour trouver le prix unitaire de chaque fruit.
\end{enumerate}
\end{exercice}

\begin{corrige}[Exercice]
\begin{enumerate}
\item $\det(M) = 4\times 5 - 3\times 2 = 20 - 6 = 14 \neq 0$ : $M$ est inversible. $\det(N) = 6\times 2 - 3\times 4 = 12 - 12 = 0$ : $N$ n'est pas inversible.
\item $M^{-1} = \dfrac{1}{14}\begin{pmatrix} 5 & -3 \\ -2 & 4 \end{pmatrix}$. Vérification : $M\times M^{-1} = \dfrac{1}{14}\begin{pmatrix} 20-6 & -12+12 \\ 10-10 & -6+20 \end{pmatrix} = \dfrac{1}{14}\begin{pmatrix} 14 & 0 \\ 0 & 14 \end{pmatrix} = I_2$. \checkmark
\item Le système s'écrit $MX = B$ avec $M = \begin{pmatrix} 4 & 3 \\ 2 & 5 \end{pmatrix}$ et $B = \begin{pmatrix} 2900 \\ 2700 \end{pmatrix}$.
On a $\det(M) = 14 \neq 0$, donc $M$ est inversible et $M^{-1} = \dfrac{1}{14}\begin{pmatrix} 5 & -3 \\ -2 & 4 \end{pmatrix}$.
\[ X = M^{-1}B = \frac{1}{14}\begin{pmatrix} 5 & -3 \\ -2 & 4 \end{pmatrix}\begin{pmatrix} 2900 \\ 2700 \end{pmatrix} = \frac{1}{14}\begin{pmatrix} 5\times 2900 - 3\times 2700 \\ -2\times 2900 + 4\times 2700 \end{pmatrix} = \frac{1}{14}\begin{pmatrix} 14500 - 8100 \\ -5800 + 10800 \end{pmatrix} = \frac{1}{14}\begin{pmatrix} 6400 \\ 5000 \end{pmatrix}. \]
Ainsi $x = \dfrac{6400}{14} = \dfrac{3200}{7} \approx 457,14$ FCFA et $y = \dfrac{5000}{14} = \dfrac{2500}{7} \approx 357,14$ FCFA.
Vérification : $4\times \frac{3200}{7} + 3\times \frac{2500}{7} = \frac{12800+7500}{7} = \frac{20300}{7} = 2900$ \checkmark{} et $2\times \frac{3200}{7} + 5\times \frac{2500}{7} = \frac{6400+12500}{7} = \frac{18900}{7} = 2700$ \checkmark.
\end{enumerate}
\end{corrige}

\begin{aretenir}[L'essentiel de la section]
\begin{itemize}
\item $\det\begin{pmatrix} a & b \\ c & d \end{pmatrix} = ad - bc$ : diagonale descendante \emph{moins} diagonale montante.
\item $A$ inversible $\iff \det(A) \neq 0$, et alors $A^{-1} = \dfrac{1}{\det(A)}\begin{pmatrix} d & -b \\ -c & a \end{pmatrix}$ : on échange $a$ et $d$, on change les signes de $b$ et $c$, on divise par le déterminant.
\item $\det(AB) = \det(A)\det(B)$ et $\det(A^{-1}) = \dfrac{1}{\det(A)}$.
\item Système $AX = B$ : si $\det(A) \neq 0$, solution unique $X = A^{-1}B$ ; si $\det(A) = 0$, aucune solution ou une infinité.
\item Géométriquement, $|\det(A)|$ est l'aire du parallélogramme construit sur les colonnes de $A$.
\end{itemize}
\end{aretenir}

\coursec{Matrice d'une application linéaire dans une base}

Dans la section précédente, nous avons appris à manipuler les matrices carrées d'ordre 2 : déterminant, inverse, lien avec les systèmes de deux équations à deux inconnues. Une question naturelle se pose alors : d'où viennent ces tableaux de nombres, et à quoi servent-ils vraiment ? La réponse est magnifique : une matrice n'est pas un simple tableau, c'est la \emph{carte d'identité} d'une transformation géométrique. Symétrie, rotation, projection, homothétie... chaque transformation du plan qui « respecte l'addition et la multiplication par un réel » se code entièrement par quatre nombres. C'est ce codage que nous allons construire dans cette section.

\courssub{Rappels : applications linéaires et bases du plan}

\begin{definition}[Application linéaire du plan]
Soit $E$ le plan vectoriel. Une application $f : E \to E$ est dite \cle{linéaire} si elle vérifie les deux propriétés suivantes, pour tous vecteurs $\vec{u}, \vec{v}$ et tout réel $\lambda$ :
\begin{enumerate}
\item $f(\vec{u} + \vec{v}) = f(\vec{u}) + f(\vec{v})$ ;
\item $f(\lambda \vec{u}) = \lambda f(\vec{u})$.
\end{enumerate}
De façon équivalente, $f$ est linéaire si et seulement si, pour tous $\vec{u}, \vec{v}$ et tout réel $\lambda$ :
\[ f(\lambda \vec{u} + \vec{v}) = \lambda f(\vec{u}) + f(\vec{v}). \]
Une application linéaire bijective du plan sur lui-même est appelée un \cle{automorphisme} du plan.
\end{definition}

\begin{definition}[Base du plan vectoriel]
Un couple $(\vec{i}, \vec{j})$ de vecteurs non colinéaires du plan est appelé une \cle{base} du plan vectoriel. Dire que $(\vec{i}, \vec{j})$ est une base signifie que tout vecteur $\vec{u}$ du plan s'écrit de manière \emph{unique} sous la forme :
\[ \vec{u} = x\vec{i} + y\vec{j}, \qquad (x, y) \in \mathbb{R}^2. \]
Les réels $x$ et $y$ sont les \cle{coordonnées} de $\vec{u}$ dans la base $(\vec{i}, \vec{j})$. On note alors $\vec{u}\begin{pmatrix} x \\ y \end{pmatrix}$ ou simplement $\vec{u}(x ; y)$.
\end{definition}

\begin{exemplebox}[Applications linéaires bien connues]
Sont des applications linéaires du plan : l'identité, les homothéties vectorielles, les symétries par rapport à une droite passant par l'origine, les projections sur une droite (parallèlement à une autre), les rotations de centre $O$. En revanche, une translation de vecteur non nul n'est \emph{pas} linéaire (elle ne fixe pas $\vec{0}$ : or toute application linéaire vérifie $f(\vec{0}) = \vec{0}$).
\end{exemplebox}

\courssub{Théorème fondamental : une application linéaire est déterminée par les images des vecteurs de base}

Voici l'idée qui rend tout le reste possible. Imagine que tu connaisses seulement les images de deux vecteurs de base, $\vec{i}$ et $\vec{j}$. Connais-tu alors l'image de \emph{n'importe quel} vecteur ? Oui ! Car tout vecteur $\vec{u} = x\vec{i} + y\vec{j}$ est une combinaison de $\vec{i}$ et $\vec{j}$, et la linéarité « transporte » cette combinaison :
\[ f(\vec{u}) = f(x\vec{i} + y\vec{j}) = x f(\vec{i}) + y f(\vec{j}). \]

\begin{propriete}[Théorème fondamental]
Soit $(\vec{i}, \vec{j})$ une base du plan vectoriel. Une application linéaire $f$ du plan est \cle{entièrement déterminée} par la donnée des deux vecteurs $f(\vec{i})$ et $f(\vec{j})$. Autrement dit : deux applications linéaires qui coïncident sur $\vec{i}$ et sur $\vec{j}$ sont égales partout.
\end{propriete}

\begin{experience}[Démonstration guidée du théorème fondamental]
Soient $f$ et $g$ deux applications linéaires telles que $f(\vec{i}) = g(\vec{i})$ et $f(\vec{j}) = g(\vec{j})$. Montrons que $f = g$.
\begin{enumerate}
\item Soit $\vec{u}$ un vecteur quelconque du plan. Puisque $(\vec{i}, \vec{j})$ est une base, il existe deux réels $x$ et $y$ tels que $\vec{u} = x\vec{i} + y\vec{j}$.
\item Appliquons $f$ en utilisant sa linéarité : $f(\vec{u}) = f(x\vec{i} + y\vec{j}) = x f(\vec{i}) + y f(\vec{j})$.
\item Appliquons de même $g$ : $g(\vec{u}) = x g(\vec{i}) + y g(\vec{j})$.
\item Par hypothèse, $f(\vec{i}) = g(\vec{i})$ et $f(\vec{j}) = g(\vec{j})$, donc $f(\vec{u}) = g(\vec{u})$.
\item Le vecteur $\vec{u}$ étant quelconque, on conclut que $f = g$. \hfill $\square$
\end{enumerate}
Ce résultat est capital : pour décrire complètement une application linéaire, il suffit de connaître \emph{quatre nombres} : les deux coordonnées de $f(\vec{i})$ et les deux coordonnées de $f(\vec{j})$. Ces quatre nombres, nous allons les ranger dans une matrice.
\end{experience}

\courssub{Définition de la matrice d'une application linéaire dans une base}

\begin{definition}[Matrice d'une application linéaire dans une base]
Soit $f$ une application linéaire du plan et $(\vec{i}, \vec{j})$ une base. Supposons que :
\[ f(\vec{i}) = a\vec{i} + c\vec{j} \qquad \text{et} \qquad f(\vec{j}) = b\vec{i} + d\vec{j}. \]
La \cle{matrice de $f$ dans la base $(\vec{i}, \vec{j})$}, notée $\mathrm{Mat}(f)$ ou $M_{(\vec{i},\vec{j})}(f)$, est la matrice carrée d'ordre 2 dont \textbf{les colonnes sont les coordonnées de $f(\vec{i})$ et de $f(\vec{j})$ dans la base $(\vec{i}, \vec{j})$} :
\[ \mathrm{Mat}(f) = \begin{pmatrix} a & b \\ c & d \end{pmatrix}. \]
La première colonne contient les coordonnées de $f(\vec{i})$, la deuxième colonne celles de $f(\vec{j})$.
\end{definition}

\begin{attention}[Colonnes, pas lignes !]
L'erreur la plus fréquente au Probatoire est d'écrire les coordonnées de $f(\vec{i})$ et $f(\vec{j})$ \emph{en lignes} au lieu de les écrire \emph{en colonnes}. Si $f(\vec{i}) = a\vec{i} + c\vec{j}$ et $f(\vec{j}) = b\vec{i} + d\vec{j}$, la bonne matrice est $\begin{pmatrix} a & b \\ c & d \end{pmatrix}$ et non $\begin{pmatrix} a & c \\ b & d \end{pmatrix}$. Moyen mnémotechnique : « les \textbf{C}olonnes sont les images des ve\textbf{C}teurs de base ». Vérifie systématiquement ta matrice en calculant l'image de $\vec{i}$ par produit : la première colonne doit ressortir.
\end{attention}

\begin{methode}[Écrire la matrice d'une application linéaire]
Pour déterminer la matrice d'une application linéaire $f$ dans une base $(\vec{i}, \vec{j})$ :
\begin{enumerate}
\item Calculer le vecteur $f(\vec{i})$ et l'exprimer dans la base : $f(\vec{i}) = a\vec{i} + c\vec{j}$.
\item Calculer le vecteur $f(\vec{j})$ et l'exprimer dans la base : $f(\vec{j}) = b\vec{i} + d\vec{j}$.
\item Placer les coordonnées de $f(\vec{i})$ en \textbf{première colonne} et celles de $f(\vec{j})$ en \textbf{deuxième colonne} :
\[ \mathrm{Mat}(f) = \begin{pmatrix} a & b \\ c & d \end{pmatrix}. \]
\item Vérifier : le produit de la matrice par la colonne $\begin{pmatrix} 1 \\ 0 \end{pmatrix}$ doit redonner la première colonne, et le produit par $\begin{pmatrix} 0 \\ 1 \end{pmatrix}$ la deuxième.
\end{enumerate}
\end{methode}

\begin{popfigure}
\begin{center}
\begin{tikzpicture}[scale=1.05]
% Case 1 : f(i)
\draw[fill=popblueL] (0,0) rectangle (0.9,1.8);
\node at (0.45,1.5) {$a$};
\node at (0.45,0.6) {$c$};
\node[below, popdark] at (0.45,-0.1) {$f(\vec{i})$};
% Case 2 : f(j)
\draw[fill=poporangeL] (1.5,0) rectangle (2.4,1.8);
\node at (1.95,1.5) {$b$};
\node at (1.95,0.6) {$d$};
\node[below, popdark] at (1.95,-0.1) {$f(\vec{j})$};
% Arrow
\draw[->, very thick, popdark] (2.9,0.9) -- (3.7,0.9);
% Matrix
\draw[thick, popink] (4.0,0) -- (4.0,1.8);
\draw[thick, popink] (5.6,0) -- (5.6,1.8);
\draw[thick, popink] (4.0,0) -- (3.85,0); \draw[thick, popink] (4.0,1.8) -- (3.85,1.8);
\draw[thick, popink] (5.6,0) -- (5.75,0); \draw[thick, popink] (5.6,1.8) -- (5.75,1.8);
\node[fill=popblueL, minimum width=0.5cm] at (4.45,1.35) {$a$};
\node[fill=poporangeL, minimum width=0.5cm] at (5.15,1.35) {$b$};
\node[fill=popblueL, minimum width=0.5cm] at (4.45,0.45) {$c$};
\node[fill=poporangeL, minimum width=0.5cm] at (5.15,0.45) {$d$};
\node[below, popdark] at (4.8,-0.1) {$\mathrm{Mat}(f)$};
\end{tikzpicture}
\end{center}
\legende{La règle d'or : les colonnes de la matrice sont les coordonnées des images des vecteurs de base.}
\end{popfigure}

\courssub{Exemples fondamentaux à connaître}

\begin{exemplebox}[Matrices des applications linéaires usuelles]
Dans une base $(\vec{i}, \vec{j})$ orthonormée directe :
\begin{enumerate}
\item \textbf{Identité} : $f(\vec{i}) = \vec{i}$ et $f(\vec{j}) = \vec{j}$, donc
\[ \mathrm{Mat}(\mathrm{id}) = \begin{pmatrix} 1 & 0 \\ 0 & 1 \end{pmatrix} = I_2. \]
\item \textbf{Homothétie de rapport $k$} : $f(\vec{i}) = k\vec{i}$ et $f(\vec{j}) = k\vec{j}$, donc
\[ \mathrm{Mat}(f) = \begin{pmatrix} k & 0 \\ 0 & k \end{pmatrix} = k I_2. \]
\item \textbf{Symétrie d'axe $(Ox)$} (axe dirigé par $\vec{i}$) : $f(\vec{i}) = \vec{i}$ et $f(\vec{j}) = -\vec{j}$, donc
\[ \mathrm{Mat}(f) = \begin{pmatrix} 1 & 0 \\ 0 & -1 \end{pmatrix}. \]
\item \textbf{Projection sur l'axe $(Ox)$ parallèlement à l'axe $(Oy)$} : $f(\vec{i}) = \vec{i}$ et $f(\vec{j}) = \vec{0}$, donc
\[ \mathrm{Mat}(f) = \begin{pmatrix} 1 & 0 \\ 0 & 0 \end{pmatrix}. \]
Cette matrice n'est pas inversible ($\det = 0$) : normal, une projection « écrase » le plan sur une droite, elle n'est pas bijective.
\end{enumerate}
\end{exemplebox}

\begin{popfigure}
\begin{center}
\begin{tikzpicture}[scale=1.6, >=Stealth]
% Axes
\draw[->, thick, popmuted] (-2.2,0) -- (2.4,0) node[below left] {$x$};
\draw[->, thick, popmuted] (0,-1.9) -- (0,2.1) node[below left] {$y$};
\node[below left] at (0,0) {$O$};
% Base vectors
\draw[->, very thick, popblue] (0,0) -- (1,0) node[midway, below] {$\vec{i}$};
\draw[->, very thick, popblue] (0,0) -- (0,1) node[midway, left] {$\vec{j}$};
% Images of base vectors
\draw[->, very thick, poporange] (0,0) -- (1,0) node[pos=0.75, above] {$f(\vec{i})$};
\draw[->, very thick, poporange] (0,0) -- (0,-1) node[midway, left] {$f(\vec{j})$};
% Vector u and its image
\draw[->, very thick, popgreen] (0,0) -- (1.6,1.2) node[above right] {$\vec{u}(1{,}6\,;\,1{,}2)$};
\draw[->, very thick, poppurple] (0,0) -- (1.6,-1.2) node[below right] {$f(\vec{u})(1{,}6\,;\,-1{,}2)$};
% Dashed symmetry lines
\draw[dashed, popmuted] (1.6,1.2) -- (1.6,-1.2);
\draw[dashed, popmuted] (0,1) -- (0,-1);
% Points
\fill[popgreen] (1.6,1.2) circle (1.3pt);
\fill[poppurple] (1.6,-1.2) circle (1.3pt);
% Matrix annotation
\node[align=center, popdark] at (3.6,0.8) {$\mathrm{Mat}(f) = \begin{pmatrix} 1 & 0 \\ 0 & -1 \end{pmatrix}$};
\node[align=center, popmuted] at (3.6,-0.4) {\small 1\textsuperscript{re} colonne : $f(\vec{i})$\\ \small 2\textsuperscript{me} colonne : $f(\vec{j})$};
\end{tikzpicture}
\end{center}
\legende{Symétrie d'axe $(Ox)$ : $\vec{i}$ est invariant, $\vec{j}$ est envoyé sur $-\vec{j}$. Les colonnes de la matrice sont les coordonnées de ces deux images.}
\end{popfigure}

\courssub{Coordonnées de l'image d'un vecteur : le produit $AX$}

Nous arrivons au résultat pratique le plus important de cette section : une fois la matrice connue, le calcul de l'image de \emph{n'importe quel} vecteur se réduit à un simple produit de matrices.

\begin{propriete}[Théorème : calcul de l'image par produit matriciel]
Soit $f$ une application linéaire du plan, de matrice $A$ dans une base $(\vec{i}, \vec{j})$. Soit $\vec{u}$ un vecteur de coordonnées $(x ; y)$ dans cette base, et notons $X = \begin{pmatrix} x \\ y \end{pmatrix}$ la colonne de ses coordonnées. Alors les coordonnées $(x' ; y')$ du vecteur image $f(\vec{u})$ sont données par la colonne :
\[ X' = AX, \qquad \text{c'est-à-dire} \qquad \begin{pmatrix} x' \\ y' \end{pmatrix} = \begin{pmatrix} a & b \\ c & d \end{pmatrix} \begin{pmatrix} x \\ y \end{pmatrix} = \begin{pmatrix} ax + by \\ cx + dy \end{pmatrix}. \]
Cette démonstration est exigible : elle utilise uniquement la linéarité et la définition de la matrice.
\end{propriete}

\begin{experience}[Démonstration guidée : pourquoi $X' = AX$]
Soit $f$ de matrice $A = \begin{pmatrix} a & b \\ c & d \end{pmatrix}$ dans la base $(\vec{i}, \vec{j})$.
\begin{enumerate}
\item Par définition de la matrice, les colonnes donnent les images des vecteurs de base :
\[ f(\vec{i}) = a\vec{i} + c\vec{j} \qquad \text{et} \qquad f(\vec{j}) = b\vec{i} + d\vec{j}. \]
\item Soit $\vec{u}$ de coordonnées $(x ; y)$. Par définition des coordonnées : $\vec{u} = x\vec{i} + y\vec{j}$.
\item Appliquons $f$ et utilisons la linéarité :
\[ f(\vec{u}) = f(x\vec{i} + y\vec{j}) = x f(\vec{i}) + y f(\vec{j}). \]
\item Remplaçons $f(\vec{i})$ et $f(\vec{j})$ par leurs expressions :
\[ f(\vec{u}) = x(a\vec{i} + c\vec{j}) + y(b\vec{i} + d\vec{j}). \]
\item Regroupons les termes en $\vec{i}$ et en $\vec{j}$ :
\[ f(\vec{u}) = (ax + by)\vec{i} + (cx + dy)\vec{j}. \]
\item Les coordonnées de $f(\vec{u})$ sont donc $(ax + by \,;\, cx + dy)$. Or le produit matriciel donne exactement :
\[ AX = \begin{pmatrix} a & b \\ c & d \end{pmatrix}\begin{pmatrix} x \\ y \end{pmatrix} = \begin{pmatrix} ax + by \\ cx + dy \end{pmatrix}. \]
On a bien $X' = AX$. \hfill $\square$
\end{enumerate}
\end{experience}

\begin{aretenir}[Le mécanisme à retenir]
Pour calculer l'image d'un vecteur $\vec{u}(x ; y)$ par une application linéaire $f$ de matrice $A$ :
\[ \underbrace{\begin{pmatrix} x' \\ y' \end{pmatrix}}_{\text{coordonnées de } f(\vec{u})} = \underbrace{\begin{pmatrix} a & b \\ c & d \end{pmatrix}}_{\mathrm{Mat}(f)} \underbrace{\begin{pmatrix} x \\ y \end{pmatrix}}_{\text{coordonnées de } \vec{u}}. \]
On en déduit l'expression analytique de $f$ : $f(x ; y) = (ax + by \,;\, cx + dy)$.
\end{aretenir}

\begin{exoresolu}[Écrire une matrice et calculer une image]
Le plan est muni d'une base $(\vec{i}, \vec{j})$. Soit $f$ l'application linéaire définie par :
\[ f(x ; y) = (2x - y \,;\, x + 3y). \]
\begin{enumerate}
\item Écrire la matrice de $f$ dans la base $(\vec{i}, \vec{j})$.
\item Calculer les coordonnées de l'image du vecteur $\vec{u}(1 ; -2)$ par un produit matriciel, puis vérifier par l'expression de $f$.
\end{enumerate}
\tcblower
\textbf{1. Matrice de $f$.} On calcule les images des vecteurs de base. Le vecteur $\vec{i}$ a pour coordonnées $(1 ; 0)$ et $\vec{j}$ a pour coordonnées $(0 ; 1)$ :
\begin{align*}
f(\vec{i}) &= f(1 ; 0) = (2\times 1 - 0 \,;\, 1 + 3\times 0) = (2 ; 1), \\
f(\vec{j}) &= f(0 ; 1) = (2\times 0 - 1 \,;\, 0 + 3\times 1) = (-1 ; 3).
\end{align*}
On place ces coordonnées \textbf{en colonnes} :
\[ \mathrm{Mat}(f) = \begin{pmatrix} 2 & -1 \\ 1 & 3 \end{pmatrix}. \]
Remarque : on pouvait aussi lire directement les coefficients dans l'expression $f(x;y) = (2x - y \,;\, x + 3y)$ : les coefficients de $x$ forment la première colonne, ceux de $y$ la deuxième.

\textbf{2. Image de $\vec{u}(1 ; -2)$.} On effectue le produit $AX$ :
\[ \begin{pmatrix} x' \\ y' \end{pmatrix} = \begin{pmatrix} 2 & -1 \\ 1 & 3 \end{pmatrix} \begin{pmatrix} 1 \\ -2 \end{pmatrix} = \begin{pmatrix} 2\times 1 + (-1)\times(-2) \\ 1\times 1 + 3\times(-2) \end{pmatrix} = \begin{pmatrix} 2 + 2 \\ 1 - 6 \end{pmatrix} = \begin{pmatrix} 4 \\ -5 \end{pmatrix}. \]
Donc $f(\vec{u})$ a pour coordonnées $(4 ; -5)$.

\textbf{Vérification} par l'expression directe : $f(1 ; -2) = (2\times 1 - (-2) \,;\, 1 + 3\times(-2)) = (4 ; -5)$. Les deux méthodes concordent.
\end{exoresolu}

\courssub{Lecture réciproque : retrouver $f$ à partir de sa matrice}

Le théorème $X' = AX$ fonctionne dans les deux sens. Si l'on te donne une matrice $A = \begin{pmatrix} a & b \\ c & d \end{pmatrix}$ et qu'on te dit que c'est la matrice d'une application linéaire $f$ dans la base $(\vec{i}, \vec{j})$, alors tu peux immédiatement écrire :
\[ f(x ; y) = (ax + by \,;\, cx + dy). \]
Autrement dit : la \textbf{première ligne} de la matrice donne la première coordonnée de l'image, la \textbf{deuxième ligne} donne la deuxième coordonnée.

\begin{exemplebox}[De la matrice vers l'application]
Soit $g$ l'application linéaire de matrice $B = \begin{pmatrix} 0 & -1 \\ 1 & 0 \end{pmatrix}$ dans une base orthonormée directe $(\vec{i}, \vec{j})$. Alors :
\[ g(x ; y) = (0\times x + (-1)\times y \,;\, 1\times x + 0\times y) = (-y \,;\, x). \]
On reconnaît la rotation d'angle $\dfrac{\pi}{2}$ (quart de tour direct) : en effet $g(\vec{i}) = \vec{j}$ et $g(\vec{j}) = -\vec{i}$.
\end{exemplebox}

\begin{attention}[Bien distinguer lignes et colonnes]
Les deux lectures de la matrice coexistent et il ne faut pas les confondre :
\begin{itemize}
\item \textbf{Colonnes} : coordonnées des images des vecteurs de base ($f(\vec{i})$ en colonne 1, $f(\vec{j})$ en colonne 2) ;
\item \textbf{Lignes} : coefficients de l'expression analytique $f(x ; y) = (ax + by \,;\, cx + dy)$.
\end{itemize}
Dans un sujet de Probatoire, précise toujours dans quelle base tu travailles et annonce clairement : « la matrice de $f$ dans la base $(\vec{i}, \vec{j})$ est... ».
\end{attention}

\begin{exercice}[À toi de jouer]
Le plan est muni d'une base $(\vec{i}, \vec{j})$. Soit $h$ l'application linéaire définie par $h(x ; y) = (3x + 2y \,;\, x - y)$.
\begin{enumerate}
\item Écrire la matrice $A$ de $h$ dans la base $(\vec{i}, \vec{j})$.
\item Calculer, par produit matriciel, les coordonnées de l'image du vecteur $\vec{v}(2 ; 5)$.
\item $h$ est-elle un automorphisme du plan ? Justifier à l'aide du déterminant de $A$.
\end{enumerate}
\end{exercice}

\begin{corrige}[Exercice : application linéaire $h$]
\textbf{1.} $h(\vec{i}) = h(1 ; 0) = (3 ; 1)$ et $h(\vec{j}) = h(0 ; 1) = (2 ; -1)$. En plaçant ces coordonnées en colonnes :
\[ A = \begin{pmatrix} 3 & 2 \\ 1 & -1 \end{pmatrix}. \]
\textbf{2.} $A \begin{pmatrix} 2 \\ 5 \end{pmatrix} = \begin{pmatrix} 3\times 2 + 2\times 5 \\ 1\times 2 + (-1)\times 5 \end{pmatrix} = \begin{pmatrix} 16 \\ -3 \end{pmatrix}$. Donc $h(\vec{v})$ a pour coordonnées $(16 ; -3)$.
\textbf{3.} $\det(A) = 3\times(-1) - 2\times 1 = -3 - 2 = -5 \neq 0$. La matrice $A$ est inversible, donc $h$ est bijective : c'est un automorphisme du plan.
\end{corrige}

\begin{savaistu}[Des matrices par millions : l'infographie]
Chaque fois que tu tournes, zoomes ou déformes une photo sur ton téléphone, ton appareil applique en réalité des produits de matrices ! Une image numérique est un tableau de millions de pixels, chacun repéré par des coordonnées $(x ; y)$. Appliquer une rotation, une symétrie ou un redimensionnement revient à multiplier la colonne des coordonnées de chaque pixel par une matrice $2\times 2$ (ou $3\times 3$ en infographie avancée). Un téléphone moderne effectue ainsi des milliards de produits matriciels par seconde. Les jeux vidéo, les effets spéciaux au cinéma et même la cartographie numérique utilisée pour tracer les routes entre Douala et Yaoundé reposent sur exactement les mathématiques que tu viens d'apprendre : $X' = AX$.
\end{savaistu}

Dans la section suivante, nous verrons comment les opérations sur les matrices (somme, produit, inverse) traduisent les opérations sur les applications linéaires elles-mêmes : additionner deux applications, les composer, ou trouver la transformation réciproque d'un automorphisme. Tu découvriras que le produit de matrices, qui semblait si étrange, est en fait parfaitement naturel : c'est celui de la composition des transformations.

\coursec{Opérations sur les applications linéaires et matrices}

Dans la section précédente, tu as appris à associer à chaque application linéaire $f$ du plan sa matrice $\mathrm{Mat}(f)$ dans une base $(\vec{i},\vec{j})$, et à calculer les coordonnées de l'image d'un vecteur par un simple produit matriciel. Cette correspondance \og application $\leftrightarrow$ matrice \fg{} est bien plus qu'une simple traduction : c'est un véritable \emph{dictionnaire}. Additionner deux applications linéaires, les multiplier par un réel, les composer, les inverser... chaque opération sur les applications correspond exactement à une opération sur les matrices que tu as étudiée dans les sections 2, 3 et 4. C'est ce dictionnaire complet que nous construisons maintenant — et c'est lui qui fait toute la puissance du calcul matriciel.

\courssub{Somme de deux applications linéaires et produit par un réel}

\begin{definition}[Somme de deux applications linéaires]
Soient $f$ et $g$ deux applications linéaires du plan vectoriel $E$. On appelle \cle{somme} de $f$ et $g$, notée $f+g$, l'application définie pour tout vecteur $\vec{u}$ de $E$ par :
\[ (f+g)(\vec{u}) = f(\vec{u}) + g(\vec{u}). \]
De même, pour tout réel $k$, on définit l'application $kf$ par $(kf)(\vec{u}) = k\,f(\vec{u})$.
\end{definition}

On vérifie sans peine que $f+g$ et $kf$ sont encore des applications linéaires. La question naturelle est alors : quelle est la matrice de $f+g$ ? L'intuition est la bonne réponse.

\begin{propriete}[Matrice d'une somme et d'un produit par un réel]
Soient $f$ et $g$ deux applications linéaires du plan, de matrices respectives $A$ et $B$ dans une base $(\vec{i},\vec{j})$, et $k$ un réel. Alors :
\[ \mathrm{Mat}(f+g) = A + B \qquad \text{et} \qquad \mathrm{Mat}(kf) = kA. \]
Autrement dit : $\mathrm{Mat}(f+g) = \mathrm{Mat}(f) + \mathrm{Mat}(g)$ et $\mathrm{Mat}(kf) = k\,\mathrm{Mat}(f)$.
\end{propriete}

\textbf{Justification.} La première colonne de $\mathrm{Mat}(f+g)$ est formée des coordonnées de $(f+g)(\vec{i}) = f(\vec{i}) + g(\vec{i})$, c'est-à-dire la somme des premières colonnes de $A$ et de $B$. Même raisonnement pour la deuxième colonne avec $\vec{j}$. On obtient donc exactement la matrice somme $A+B$.

\begin{exemplebox}[Somme et produit par un réel]
Dans la base $(\vec{i},\vec{j})$, soient $f$ et $g$ de matrices $A = \begin{pmatrix} 1 & 2 \\ 0 & 1 \end{pmatrix}$ et $B = \begin{pmatrix} 0 & 1 \\ 1 & 0 \end{pmatrix}$. Alors :
\[ \mathrm{Mat}(f+g) = \begin{pmatrix} 1 & 3 \\ 1 & 1 \end{pmatrix}, \qquad \mathrm{Mat}(3f) = \begin{pmatrix} 3 & 6 \\ 0 & 3 \end{pmatrix}, \qquad \mathrm{Mat}(2f - g) = \begin{pmatrix} 2 & 3 \\ -1 & 2 \end{pmatrix}. \]
\end{exemplebox}

\courssub{Composée de deux applications linéaires et produit de matrices}

Voici le cœur de cette section, et sans doute le résultat le plus important de toute la leçon. Imaginons une chaîne de traitement : un vecteur $\vec{u}$ subit d'abord la transformation $f$, puis le résultat $f(\vec{u})$ subit la transformation $g$. Le vecteur final est $g\bigl(f(\vec{u})\bigr) = (g \circ f)(\vec{u})$.

\begin{attention}[L'ordre de la composition]
Dans l'écriture $g \circ f$, c'est $f$ qui agit \textbf{en premier} : $(g \circ f)(\vec{u}) = g\bigl(f(\vec{u})\bigr)$. On lit et on applique de droite à gauche. En général $g \circ f \neq f \circ g$ : la composition n'est \textbf{pas commutative}.
\end{attention}

\begin{propriete}[Matrice de la composée — démonstration exigible]
Soient $f$ et $g$ deux applications linéaires du plan, de matrices respectives $A$ et $B$ dans la base $(\vec{i},\vec{j})$. Alors :
\[ \mathrm{Mat}(g \circ f) = B \times A = \mathrm{Mat}(g) \times \mathrm{Mat}(f). \]
La matrice de la composée est le produit des matrices, \textbf{dans le même ordre} que la composition.
\end{propriete}

\begin{experience}[Démonstration guidée : pourquoi $\mathrm{Mat}(g \circ f) = BA$ ?]
Posons $A = \begin{pmatrix} a & b \\ c & d \end{pmatrix}$ et $B = \begin{pmatrix} a' & b' \\ c' & d' \end{pmatrix}$. Rappelons que les colonnes de $A$ sont les coordonnées de $f(\vec{i})$ et $f(\vec{j})$ : $f(\vec{i}) = a\vec{i} + c\vec{j}$ et $f(\vec{j}) = b\vec{i} + d\vec{j}$.

\textbf{Étape 1 : image de $\vec{i}$ par $g \circ f$.} Par linéarité de $g$ :
\[ (g \circ f)(\vec{i}) = g\bigl(a\vec{i} + c\vec{j}\bigr) = a\,g(\vec{i}) + c\,g(\vec{j}). \]
Or $g(\vec{i}) = a'\vec{i} + c'\vec{j}$ et $g(\vec{j}) = b'\vec{i} + d'\vec{j}$ (colonnes de $B$). Donc :
\begin{align*}
(g \circ f)(\vec{i}) &= a\bigl(a'\vec{i} + c'\vec{j}\bigr) + c\bigl(b'\vec{i} + d'\vec{j}\bigr) \\
&= (a'a + b'c)\,\vec{i} + (c'a + d'c)\,\vec{j}.
\end{align*}

\textbf{Étape 2 : image de $\vec{j}$ par $g \circ f$.} Le même calcul donne :
\[ (g \circ f)(\vec{j}) = (a'b + b'd)\,\vec{i} + (c'b + d'd)\,\vec{j}. \]

\textbf{Étape 3 : lecture de la matrice.} Les coordonnées trouvées forment les colonnes de $\mathrm{Mat}(g \circ f)$ :
\[ \mathrm{Mat}(g \circ f) = \begin{pmatrix} a'a + b'c & a'b + b'd \\ c'a + d'c & c'b + d'd \end{pmatrix} = \begin{pmatrix} a' & b' \\ c' & d' \end{pmatrix} \begin{pmatrix} a & b \\ c & d \end{pmatrix} = B \times A. \]
C'est exactement le produit $B \times A$, dans cet ordre. $\blacksquare$
\end{experience}

\begin{popfigure}
\begin{center}
\begin{tikzpicture}[>=Stealth, node distance=1.1cm]
  \node[draw=popblue, fill=popblueL, rounded corners, minimum width=1.6cm, minimum height=0.9cm] (u) {$\vec{u}$};
  \node[draw=poporange, fill=poporangeL, rounded corners, minimum width=1.9cm, minimum height=0.9cm, right=2.6cm of u] (fu) {$f(\vec{u}) = AX$};
  \node[draw=popgreen, fill=popgreenL, rounded corners, minimum width=2.6cm, minimum height=0.9cm, right=2.9cm of fu] (gfu) {$(g\circ f)(\vec{u}) = BAX$};
  \draw[->, very thick, poporange] (u) -- (fu) node[midway, above, popdark]{$f$ : matrice $A$};
  \draw[->, very thick, popgreen] (fu) -- (gfu) node[midway, above, popdark]{$g$ : matrice $B$};
  \draw[->, very thick, poppurple, dashed] (u.south) to[bend right=35] node[midway, below, popdark]{$g \circ f$ : matrice $B \times A$} (gfu.south);
\end{tikzpicture}
\end{center}
\legende{Chaîne de composition : le vecteur $\vec{u}$ de coordonnées $X$ subit d'abord $f$ (matrice $A$), puis $g$ (matrice $B$). La transformation globale $g \circ f$ a pour matrice le produit $B \times A$.}
\end{popfigure}

\begin{attention}[Piège n°1 : l'ordre des facteurs]
L'erreur la plus fréquente au Probatoire est d'écrire $\mathrm{Mat}(g \circ f) = A \times B$. C'est \textbf{faux} : la matrice de $g \circ f$ est $\mathrm{Mat}(g) \times \mathrm{Mat}(f)$, dans \textbf{le même ordre} que la composition. Comme le produit matriciel n'est pas commutatif, intervertir les facteurs donne en général un résultat différent. Retiens : \og première appliquée, deuxième dans le produit, à droite \fg.
\end{attention}

\begin{exoresolu}[Composées dans les deux sens et inverse]
Dans la base $(\vec{i},\vec{j})$, soient $f$ et $g$ les applications linéaires de matrices respectives
\[ A = \begin{pmatrix} 1 & 2 \\ 0 & 1 \end{pmatrix} \qquad \text{et} \qquad B = \begin{pmatrix} 0 & 1 \\ 1 & 0 \end{pmatrix}. \]
\begin{enumerate}
\item Déterminer les matrices de $g \circ f$ et de $f \circ g$. Que constate-t-on ?
\item Montrer que $f$ est un automorphisme et déterminer la matrice de $f^{-1}$.
\end{enumerate}
\tcblower
\textbf{1.} On applique la règle $\mathrm{Mat}(g \circ f) = B \times A$ :
\[ \mathrm{Mat}(g \circ f) = \begin{pmatrix} 0 & 1 \\ 1 & 0 \end{pmatrix} \begin{pmatrix} 1 & 2 \\ 0 & 1 \end{pmatrix} = \begin{pmatrix} 0\times 1 + 1\times 0 & 0\times 2 + 1\times 1 \\ 1\times 1 + 0\times 0 & 1\times 2 + 0\times 1 \end{pmatrix} = \begin{pmatrix} 0 & 1 \\ 1 & 2 \end{pmatrix}. \]
\[ \mathrm{Mat}(f \circ g) = A \times B = \begin{pmatrix} 1 & 2 \\ 0 & 1 \end{pmatrix} \begin{pmatrix} 0 & 1 \\ 1 & 0 \end{pmatrix} = \begin{pmatrix} 2 & 1 \\ 1 & 0 \end{pmatrix}. \]
On constate que $\mathrm{Mat}(g \circ f) \neq \mathrm{Mat}(f \circ g)$, donc $g \circ f \neq f \circ g$ : la composition n'est pas commutative, ce qui confirme l'importance de l'ordre.

\textbf{2.} $\det(A) = 1\times 1 - 2\times 0 = 1 \neq 0$, donc $A$ est inversible et $f$ est un automorphisme. La formule de l'inverse donne :
\[ \mathrm{Mat}(f^{-1}) = A^{-1} = \frac{1}{\det(A)}\begin{pmatrix} 1 & -2 \\ 0 & 1 \end{pmatrix} = \begin{pmatrix} 1 & -2 \\ 0 & 1 \end{pmatrix}. \]
\textbf{Vérification} : $A \times A^{-1} = \begin{pmatrix} 1 & 2 \\ 0 & 1 \end{pmatrix}\begin{pmatrix} 1 & -2 \\ 0 & 1 \end{pmatrix} = \begin{pmatrix} 1 & 0 \\ 0 & 1 \end{pmatrix} = I_2$. La composée $f^{-1} \circ f$ a bien pour matrice l'identité : c'est l'application identique.
\end{exoresolu}

\courssub{Automorphismes et matrice de la réciproque}

\begin{definition}[Automorphisme]
On appelle \cle{automorphisme} du plan vectoriel toute application linéaire \textbf{bijective} du plan dans lui-même.
\end{definition}

Géométriquement, un automorphisme est une transformation linéaire qui ne \og écrase \fg rien : deux vecteurs distincts ont des images distinctes, et tout vecteur est l'image d'un unique vecteur. Le critère algébrique est remarquablement simple.

\begin{propriete}[Théorème : critère d'inversibilité et matrice de la réciproque]
Soit $f$ une application linéaire du plan, de matrice $A$ dans une base $(\vec{i},\vec{j})$. Alors :
\[ f \text{ est un automorphisme} \iff A \text{ est inversible} \iff \det(A) \neq 0. \]
Dans ce cas, la réciproque $f^{-1}$ est aussi une application linéaire, et :
\[ \mathrm{Mat}(f^{-1}) = A^{-1} = \bigl[\mathrm{Mat}(f)\bigr]^{-1}. \]
\end{propriete}

\textbf{Justification.} Si $f$ est bijective, sa réciproque $f^{-1}$ est linéaire et vérifie $f^{-1} \circ f = f \circ f^{-1} = \mathrm{Id}$. En passant aux matrices : $\mathrm{Mat}(f^{-1}) \times A = A \times \mathrm{Mat}(f^{-1}) = I_2$, ce qui prouve que $A$ est inversible et que $\mathrm{Mat}(f^{-1}) = A^{-1}$. Réciproquement, si $A$ est inversible, l'application linéaire de matrice $A^{-1}$ est exactement la réciproque de $f$, donc $f$ est bijective.

\begin{methode}[Matrice d'une transformation composée et de sa réciproque]
Pour déterminer la matrice d'une transformation composée puis celle de sa réciproque :
\begin{enumerate}
\item \textbf{Identifier} les transformations élémentaires et écrire la matrice de chacune dans la base donnée.
\item \textbf{Composer} : la matrice de $g \circ f$ est le produit $B \times A$, en respectant l'ordre (la première transformation appliquée est à droite dans le produit).
\item \textbf{Tester l'inversibilité} : calculer le déterminant du produit. S'il est non nul, la composée est un automorphisme.
\item \textbf{Inverser} : appliquer la formule $M^{-1} = \dfrac{1}{\det(M)}\begin{pmatrix} d & -b \\ -c & a \end{pmatrix}$ pour $M = \begin{pmatrix} a & b \\ c & d \end{pmatrix}$.
\item \textbf{Vérifier} (réflexe de rigueur) : contrôler que $M \times M^{-1} = I_2$.
\end{enumerate}
\end{methode}

\begin{exemplebox}[Application géométrique : symétrie puis homothétie]
Plaçons-nous dans la base orthonormée $(\vec{i},\vec{j})$. Soit $s$ la symétrie orthogonale d'axe $(O;\vec{i})$ (l'axe des abscisses), de matrice $S = \begin{pmatrix} 1 & 0 \\ 0 & -1 \end{pmatrix}$, et $h$ l'homothétie de centre $O$ et de rapport $2$, de matrice $H = \begin{pmatrix} 2 & 0 \\ 0 & 2 \end{pmatrix}$.

La transformation $t = h \circ s$ (on réfléchit d'abord, puis on agrandit) a pour matrice :
\[ \mathrm{Mat}(t) = H \times S = \begin{pmatrix} 2 & 0 \\ 0 & 2 \end{pmatrix}\begin{pmatrix} 1 & 0 \\ 0 & -1 \end{pmatrix} = \begin{pmatrix} 2 & 0 \\ 0 & -2 \end{pmatrix}. \]
Comme $\det(H S) = -4 \neq 0$, $t$ est un automorphisme. Sa réciproque — intuitivement : rétrécir de rapport $\tfrac{1}{2}$ puis réfléchir à nouveau — a pour matrice :
\[ \mathrm{Mat}(t^{-1}) = (HS)^{-1} = \frac{1}{-4}\begin{pmatrix} -2 & 0 \\ 0 & 2 \end{pmatrix} = \begin{pmatrix} \tfrac{1}{2} & 0 \\ 0 & -\tfrac{1}{2} \end{pmatrix}. \]
On retrouve bien $S^{-1} \times H^{-1}$ : la réciproque d'une composée est la composée des réciproques \textbf{dans l'ordre inverse}, exactement comme pour les matrices : $(HS)^{-1} = S^{-1}H^{-1}$.
\end{exemplebox}

\begin{attention}[Deux pièges classiques sur la réciproque]
\begin{itemize}
\item Ne jamais écrire $(AB)^{-1} = A^{-1}B^{-1}$ : la bonne formule est $(AB)^{-1} = B^{-1}A^{-1}$ (ordre inversé).
\item Avant de calculer une matrice inverse, \textbf{toujours vérifier} que le déterminant est non nul. Si $\det(A) = 0$, l'application n'est pas bijective et $f^{-1}$ n'existe pas.
\end{itemize}
\end{attention}

\courssub{Synthèse : le dictionnaire applications linéaires $\leftrightarrow$ matrices}

Tout ce que nous avons établi tient en un tableau, à connaître par cœur : c'est la clé de voûte de la leçon.

\begin{center}
\begin{tabularx}{\linewidth}{|l|X|X|}
\hline
\rowcolor{popblueL} \textbf{Opération} & \textbf{Côté applications linéaires} & \textbf{Côté matrices} \\
\hline
Image d'un vecteur & $\vec{u}' = f(\vec{u})$ & $X' = AX$ \\
\hline
Somme & $f + g$ & $A + B$ \\
\hline
Produit par un réel & $kf$ & $kA$ \\
\hline
Composée & $g \circ f$ & $B \times A$ (même ordre) \\
\hline
Identité & $\mathrm{Id}$ : $\vec{u} \mapsto \vec{u}$ & $I_2 = \begin{pmatrix} 1 & 0 \\ 0 & 1 \end{pmatrix}$ \\
\hline
Bijectivité & $f$ automorphisme & $\det(A) \neq 0$ \\
\hline
Réciproque & $f^{-1}$ & $A^{-1}$ \\
\hline
\end{tabularx}
\end{center}

Ce tableau résume la philosophie de tout le chapitre : pour étudier une transformation du plan, on la remplace par sa matrice, on fait les calculs (plus simples, purement algébriques), puis on traduit le résultat en langage géométrique. Les matrices sont la \og langue de calcul \fg{} des transformations géométriques.

\begin{exercice}[Pour t'entraîner]
Dans la base $(\vec{i},\vec{j})$, soient $f$ et $g$ les applications linéaires de matrices $A = \begin{pmatrix} 2 & 1 \\ 1 & 1 \end{pmatrix}$ et $B = \begin{pmatrix} 1 & -1 \\ 0 & 3 \end{pmatrix}$.
\begin{enumerate}
\item Déterminer les matrices de $f + 2g$, de $g \circ f$ et de $f \circ g$.
\item $f$ est-elle un automorphisme ? Si oui, déterminer $\mathrm{Mat}(f^{-1})$.
\item $g$ est-elle un automorphisme ? Justifier.
\end{enumerate}
\end{exercice}

\begin{corrige}[Exercice 1]
\textbf{1.} $A + 2B = \begin{pmatrix} 2 & 1 \\ 1 & 1 \end{pmatrix} + \begin{pmatrix} 2 & -2 \\ 0 & 6 \end{pmatrix} = \begin{pmatrix} 4 & -1 \\ 1 & 7 \end{pmatrix}$. Ensuite :
\[ \mathrm{Mat}(g \circ f) = BA = \begin{pmatrix} 1 & -1 \\ 0 & 3 \end{pmatrix}\begin{pmatrix} 2 & 1 \\ 1 & 1 \end{pmatrix} = \begin{pmatrix} 1 & 0 \\ 3 & 3 \end{pmatrix}, \quad \mathrm{Mat}(f \circ g) = AB = \begin{pmatrix} 2 & 1 \\ 1 & 1 \end{pmatrix}\begin{pmatrix} 1 & -1 \\ 0 & 3 \end{pmatrix} = \begin{pmatrix} 2 & 1 \\ 1 & 2 \end{pmatrix}. \]
\textbf{2.} $\det(A) = 2 - 1 = 1 \neq 0$, donc $f$ est un automorphisme et $A^{-1} = \begin{pmatrix} 1 & -1 \\ -1 & 2 \end{pmatrix}$. Vérification : $A A^{-1} = I_2$.
\textbf{3.} $\det(B) = 3 - 0 = 3 \neq 0$, donc $g$ est aussi un automorphisme, avec $B^{-1} = \dfrac{1}{3}\begin{pmatrix} 3 & 1 \\ 0 & 1 \end{pmatrix}$.
\end{corrige}

\begin{savaistu}[Des matrices 2D aux mondes 3D des jeux vidéo]
Quand un personnage tourne, saute ou qu'une caméra pivote dans un jeu vidéo, la console calcule en permanence des \textbf{composées de transformations} : rotations, translations, homothéties, projections. Chacune est codée par une matrice (de taille $3\times 3$ ou $4\times 4$), et la transformation totale s'obtient exactement comme ici : \textbf{par un produit de matrices, dans le bon ordre}. Les cartes graphiques modernes effectuent des milliards de produits matriciels par seconde ! La règle $\mathrm{Mat}(g \circ f) = \mathrm{Mat}(g) \times \mathrm{Mat}(f)$ que tu viens de démontrer est donc littéralement à l'œuvre dans chaque image de tes jeux préférés — et aussi dans les logiciels de cartographie utilisés pour l'aménagement du territoire camerounais.
\end{savaistu}

\newpage

\coursec{Activité d'intégration}

\courssub{Objectifs de l'activité}

À l'issue de cette activité d'intégration, tu seras capable de :
\begin{itemize}
\item traduire la présentation de données commerciales réelles par des matrices (matrice des quantités, matrice colonne des prix) ;
\item calculer la somme de deux matrices et le produit d'une matrice par un réel dans un contexte de gestion de stocks et de tarifs ;
\item calculer un produit de matrices pour obtenir des recettes hebdomadaires ;
\item modéliser un problème de tarification par un système de deux équations linéaires dans $\mathbb{R}^2$, l'écrire sous forme matricielle $AX=B$, vérifier l'inversibilité de $A$ par son déterminant et le résoudre par le calcul de $A^{-1}B$ ;
\item écrire la matrice d'une application linéaire simple du plan et calculer les coordonnées de l'image d'un vecteur par produit matriciel ;
\item déterminer les matrices de la somme, de la composée et de la réciproque d'applications linéaires.
\end{itemize}

\courssub{Situation}

Au quartier Mfoundi, non loin du marché central de Yaoundé, Mme \textbf{Ngo Bell} tient une cybercafétéria très fréquentée par les élèves et les étudiants. Elle propose trois services payants :

\begin{itemize}
\item la \cle{connexion internet}, facturée $300$ FCFA l'heure ;
\item l'\cle{impression} de documents, facturée $100$ FCFA la page ;
\item la \cle{photocopie}, facturée $50$ FCFA la page.
\end{itemize}

Soucieuse de bien gérer son affaire, Mme Ngo Bell a relevé les quantités vendues pendant deux semaines consécutives :

\begin{center}
\begin{tabularx}{\linewidth}{|l|X|X|X|}
\hline
\rowcolor{popblueL} \textbf{Semaine} & \textbf{Connexion (heures)} & \textbf{Impressions (pages)} & \textbf{Photocopies (pages)} \\
\hline
Semaine 1 & $120$ & $80$ & $200$ \\
\hline
Semaine 2 & $150$ & $60$ & $180$ \\
\hline
\end{tabularx}
\end{center}

Par ailleurs, Mme Ngo Bell propose deux \cle{forfaits} combinés très appréciés, dont elle a oublié les prix unitaires exacts $x$ (forfait « Étudiant ») et $y$ (forfait « Secrétariat »), en FCFA. Elle se souvient seulement de deux journées de vente :
\begin{itemize}
\item lundi : $2$ forfaits « Étudiant » et $1$ forfait « Secrétariat » ont rapporté $1\,100$ FCFA ;
\item mardi : $1$ forfait « Étudiant » et $3$ forfaits « Secrétariat » ont rapporté $1\,900$ FCFA.
\end{itemize}

Enfin, pour fidéliser sa clientèle, Mme Ngo Bell envisage une \cle{remise fidélité} qui transforme tout couple de prix $(a\,;\,b)$ — où $a$ est un prix de connexion et $b$ un prix d'impression — en le couple $(0{,}9a\,;\,0{,}85b)$ : c'est une application linéaire du plan dans lui-même.

\courssub{Consignes}

\begin{enumerate}
\item \textbf{Organiser les données.} Écrire la matrice $Q$ de type $(2\,;\,3)$ des quantités vendues (lignes : semaines ; colonnes : connexion, impression, photocopie), puis la matrice colonne $P$ des prix unitaires.
\item \textbf{Ventes cumulées.} On note $Q_1$ et $Q_2$ les matrices lignes des semaines 1 et 2. Calculer $Q_1+Q_2$ et interpréter le résultat.
\item \textbf{Hausse des prix.} Face à l'augmentation du coût de l'électricité, Mme Ngo Bell augmente tous ses prix de \SI{15}{\percent}. Déterminer la nouvelle matrice des prix $P'$ à l'aide d'un produit d'une matrice par un réel.
\item \textbf{Recettes hebdomadaires.} Justifier que le produit $QP$ est défini, le calculer, et en déduire la recette de chaque semaine ainsi que la recette totale des deux semaines.
\item \textbf{Retrouver les prix des forfaits.}
\begin{enumerate}
\item Traduire les ventes de lundi et de mardi par un système de deux équations à deux inconnues $x$ et $y$.
\item Écrire ce système sous la forme matricielle $AX=B$ en précisant les matrices $A$, $X$ et $B$.
\item Calculer $\det(A)$ et justifier que $A$ est inversible.
\item Calculer $A^{-1}$, puis $X=A^{-1}B$. Conclure sur les prix des deux forfaits.
\end{enumerate}
\item \textbf{Remise fidélité (prolongement).} On considère l'application linéaire $f$ du plan dans lui-même définie par $f(a\,;\,b)=(0{,}9a\,;\,0{,}85b)$.
\begin{enumerate}
\item Écrire la matrice $M$ de $f$ dans la base canonique $(\vec{i}\,;\,\vec{j})$ du plan.
\item Un « panier-type » de prix est représenté par le vecteur $\vec{u}\,(300\,;\,100)$. Calculer, par un produit matriciel, les coordonnées de $f(\vec{u})$ et interpréter.
\end{enumerate}
\end{enumerate}

\courssub{Ressources à mobiliser}

\begin{aretenir}[Boîte à outils de l'activité]
\begin{itemize}
\item \textbf{Somme} de deux matrices de même type : on additionne les coefficients de même position. \textbf{Produit par un réel} $k$ : on multiplie chaque coefficient par $k$.
\item \textbf{Produit} $A\times B$ : défini si le nombre de colonnes de $A$ égale le nombre de lignes de $B$ ; le coefficient d'indice $(i\,;\,j)$ est la somme des produits des coefficients de la ligne $i$ de $A$ par ceux de la colonne $j$ de $B$.
\item \textbf{Déterminant} de $A=\begin{pmatrix} a & b \\ c & d \end{pmatrix}$ : $\det(A)=ad-bc$.
\item $A$ est \textbf{inversible} si et seulement si $\det(A)\neq 0$, et alors $A^{-1}=\dfrac{1}{ad-bc}\begin{pmatrix} d & -b \\ -c & a \end{pmatrix}$.
\item Le système $AX=B$ avec $A$ inversible admet l'unique solution $X=A^{-1}B$.
\item La \textbf{matrice d'une application linéaire} $f$ dans une base $(\vec{i}\,;\,\vec{j})$ a pour colonnes les coordonnées de $f(\vec{i})$ et $f(\vec{j})$ ; l'image d'un vecteur de colonne $X$ a pour colonne $MX$.
\item \textbf{Somme d'applications linéaires} : la matrice de $f+g$ est la somme des matrices de $f$ et $g$.
\item \textbf{Composée d'applications linéaires} : la matrice de $f\circ g$ est le produit des matrices de $f$ et $g$ (dans cet ordre).
\item \textbf{Réciproque d'un automorphisme} : si $f$ est bijective, la matrice de $f^{-1}$ est l'inverse de la matrice de $f$.
\end{itemize}
\end{aretenir}

\courssub{Démarche et corrigé commenté}

\begin{exoresolu}[La cybercafétéria de Mfoundi : corrigé complet]
Énoncé : reprendre les six consignes ci-dessus à partir des données de Mme Ngo Bell.
\tcblower
\textbf{1. Organisation des données.} En rangeant les semaines en lignes et les services en colonnes :
\[ Q=\begin{pmatrix} 120 & 80 & 200 \\ 150 & 60 & 180 \end{pmatrix}, \qquad P=\begin{pmatrix} 300 \\ 100 \\ 50 \end{pmatrix}. \]
$Q$ est de type $(2\,;\,3)$ et $P$ de type $(3\,;\,1)$.

\textbf{2. Ventes cumulées.} Les matrices lignes sont $Q_1=\begin{pmatrix} 120 & 80 & 200 \end{pmatrix}$ et $Q_2=\begin{pmatrix} 150 & 60 & 180 \end{pmatrix}$. Elles ont le même type, donc :
\[ Q_1+Q_2=\begin{pmatrix} 270 & 140 & 380 \end{pmatrix}. \]
Sur les deux semaines, la cybercafétéria a vendu $270$ heures de connexion, $140$ impressions et $380$ photocopies.

\textbf{3. Hausse de \SI{15}{\percent} des prix.} Augmenter de \SI{15}{\percent} revient à multiplier par $1+\dfrac{15}{100}=1{,}15$ :
\[ P'=1{,}15\,P=1{,}15\begin{pmatrix} 300 \\ 100 \\ 50 \end{pmatrix}=\begin{pmatrix} 345 \\ 115 \\ 57{,}5 \end{pmatrix}. \]
L'heure de connexion passe à $345$ FCFA, l'impression à $115$ FCFA, la photocopie à $57{,}5$ FCFA.

\textbf{4. Recettes hebdomadaires.} $Q$ a $3$ colonnes et $P$ a $3$ lignes : le produit $QP$ est défini et donne une matrice de type $(2\,;\,1)$.
\begin{align*}
QP&=\begin{pmatrix} 120\times 300+80\times 100+200\times 50 \\ 150\times 300+60\times 100+180\times 50 \end{pmatrix}\\
&=\begin{pmatrix} 36\,000+8\,000+10\,000 \\ 45\,000+6\,000+9\,000 \end{pmatrix}=\begin{pmatrix} 54\,000 \\ 60\,000 \end{pmatrix}.
\end{align*}
La recette de la semaine 1 est de $54\,000$ FCFA, celle de la semaine 2 de $60\,000$ FCFA, soit un total de $114\,000$ FCFA. Remarque : on retrouve ce total en multipliant la ligne cumulée $Q_1+Q_2$ par $P$, ce qui illustre la distributivité du produit matriciel : $(Q_1+Q_2)P = Q_1P + Q_2P$.

\textbf{5. Prix des forfaits.}
a) Les ventes se traduisent par le système :
\[ \begin{cases} 2x+y=1\,100 \\ x+3y=1\,900 \end{cases} \]
b) Forme matricielle $AX=B$ avec :
\[ A=\begin{pmatrix} 2 & 1 \\ 1 & 3 \end{pmatrix},\qquad X=\begin{pmatrix} x \\ y \end{pmatrix},\qquad B=\begin{pmatrix} 1\,100 \\ 1\,900 \end{pmatrix}. \]
c) $\det(A)=2\times 3-1\times 1=6-1=5$. Comme $\det(A)=5\neq 0$, la matrice $A$ est inversible et le système admet une unique solution.
d) On calcule :
\[ A^{-1}=\dfrac{1}{5}\begin{pmatrix} 3 & -1 \\ -1 & 2 \end{pmatrix}. \]
Puis :
\begin{align*}
X=A^{-1}B&=\dfrac{1}{5}\begin{pmatrix} 3 & -1 \\ -1 & 2 \end{pmatrix}\begin{pmatrix} 1\,100 \\ 1\,900 \end{pmatrix}\\
&=\dfrac{1}{5}\begin{pmatrix} 3\times 1\,100-1\times 1\,900 \\ -1\times 1\,100+2\times 1\,900 \end{pmatrix}=\dfrac{1}{5}\begin{pmatrix} 1\,400 \\ 2\,700 \end{pmatrix}=\begin{pmatrix} 280 \\ 540 \end{pmatrix}.
\end{align*}
Le forfait « Étudiant » coûte $280$ FCFA et le forfait « Secrétariat » $540$ FCFA. Vérification : $2\times 280+540=1\,100$ et $280+3\times 540=1\,900$. Le compte est bon !

\textbf{6. Remise fidélité.}
a) On calcule les images des vecteurs de base : $f(\vec{i})=f(1\,;\,0)=(0{,}9\,;\,0)$ et $f(\vec{j})=f(0\,;\,1)=(0\,;\,0{,}85)$. En plaçant ces coordonnées en colonnes :
\[ M=\begin{pmatrix} 0{,}9 & 0 \\ 0 & 0{,}85 \end{pmatrix}. \]
b) L'image du panier-type $\vec{u}\,(300\,;\,100)$ s'obtient par le produit :
\[ M\begin{pmatrix} 300 \\ 100 \end{pmatrix}=\begin{pmatrix} 0{,}9\times 300+0\times 100 \\ 0\times 300+0{,}85\times 100 \end{pmatrix}=\begin{pmatrix} 270 \\ 85 \end{pmatrix}. \]
Après remise, le client fidèle paie $270$ FCFA l'heure de connexion et $85$ FCFA la page imprimée.
\end{exoresolu}

\begin{exemplebox}[Somme, composée et réciproque d'applications linéaires (approfondissement)]
Reprenons l'application $f$ de la remise fidélité (matrice $M_f = \begin{pmatrix} 0{,}9 & 0 \\ 0 & 0{,}85 \end{pmatrix}$) et définissons l'application $g$ correspondant à la hausse des prix de \SI{15}{\percent} : $g(a\,;\,b) = (1{,}15a\,;\,1{,}15b)$. Sa matrice dans la base canonique est $M_g = 1{,}15\,I_2 = \begin{pmatrix} 1{,}15 & 0 \\ 0 & 1{,}15 \end{pmatrix}$.

\begin{enumerate}
\item \textbf{Somme $f+g$ :} $(f+g)(a\,;\,b) = f(a\,;\,b)+g(a\,;\,b) = (0{,}9a+1{,}15a\,;\,0{,}85b+1{,}15b) = (2{,}05a\,;\,2{,}00b)$.
La matrice de $f+g$ est $M_f+M_g = \begin{pmatrix} 0{,}9+1{,}15 & 0 \\ 0 & 0{,}85+1{,}15 \end{pmatrix} = \begin{pmatrix} 2{,}05 & 0 \\ 0 & 2{,}00 \end{pmatrix}$.
\item \textbf{Composée $f \circ g$ (hausse puis remise) :} $(f \circ g)(a\,;\,b) = f(1{,}15a\,;\,1{,}15b) = (0{,}9\times1{,}15a\,;\,0{,}85\times1{,}15b) = (1{,}035a\,;\,0{,}9775b)$.
La matrice de $f \circ g$ est le produit $M_f M_g = \begin{pmatrix} 0{,}9 & 0 \\ 0 & 0{,}85 \end{pmatrix}\begin{pmatrix} 1{,}15 & 0 \\ 0 & 1{,}15 \end{pmatrix} = \begin{pmatrix} 1{,}035 & 0 \\ 0 & 0{,}9775 \end{pmatrix}$.
(Remarque : comme les matrices sont diagonales, $M_f M_g = M_g M_f$, l'ordre n'importe pas ici).
\item \textbf{Réciproque de $f$ (automorphisme) :} $\det(M_f) = 0{,}9 \times 0{,}85 = 0{,}765 \neq 0$, donc $f$ est bijective.
La matrice de $f^{-1}$ est l'inverse de $M_f$ :
\[ M_f^{-1} = \frac{1}{0{,}765}\begin{pmatrix} 0{,}85 & 0 \\ 0 & 0{,}9 \end{pmatrix} = \begin{pmatrix} \frac{1}{0{,}9} & 0 \\ 0 & \frac{1}{0{,}85} \end{pmatrix} \approx \begin{pmatrix} 1{,}111 & 0 \\ 0 & 1{,}176 \end{pmatrix}. \]
On vérifie : $f^{-1}(a\,;\,b) = \left(\frac{a}{0{,}9}\,;\,\frac{b}{0{,}85}\right)$, ce qui retrouve bien les prix initiaux avant remise.
\end{enumerate}
\end{exemplebox}

\begin{attention}[Pièges rencontrés dans cette activité]
\begin{itemize}
\item Le produit $PQ$ n'est \textbf{pas défini} (une matrice $(3\,;\,1)$ ne se multiplie pas par une matrice $(2\,;\,3)$ dans cet ordre) : vérifie toujours la compatibilité des dimensions avant de calculer.
\item Augmenter de \SI{15}{\percent}, c'est multiplier par $1{,}15$, et non par $0{,}15$ ni par $15$.
\item Dans la formule de l'inverse, n'oublie pas le facteur $\dfrac{1}{\det(A)}$ : oublier le $\dfrac{1}{5}$ donnerait ici des forfaits à $1\,400$ et $2\,700$ FCFA, ce que la vérification dans le système initial dément immédiatement.
\item La matrice d'une application linéaire se construit avec les images des vecteurs de base \textbf{en colonnes}, pas en lignes.
\item Pour la composée $f \circ g$, la matrice est $M_f M_g$ (matrice de $f$ fois matrice de $g$), l'ordre est l'inverse de l'ordre d'application usuel des fonctions : on applique $g$ puis $f$, mais on multiplie $M_f$ par $M_g$.
\end{itemize}
\end{attention}

\courssub{Bilan de l'activité}

Cette activité a permis de mettre en évidence que les matrices ne sont pas un objet abstrait, mais un \cle{outil de rangement et de calcul} parfaitement adapté à la vie économique quotidienne :

\begin{itemize}
\item une matrice \textbf{présente} des données (quantités vendues, prix) de manière compacte et lisible ;
\item la \textbf{somme} et le \textbf{produit par un réel} traduisent des opérations naturelles de gestion : cumul des ventes, hausse générale des tarifs ;
\item le \textbf{produit de deux matrices} effectue en une seule écriture tous les calculs « quantité $\times$ prix » qui donnent les recettes ;
\item le \textbf{déterminant} décide si un système commercial admet une solution unique, et l'\textbf{inverse} fournit cette solution directement par $X=A^{-1}B$ ;
\item enfin, une remise proportionnelle est une \textbf{application linéaire}, dont la matrice permet de calculer instantanément le nouveau prix de n'importe quel panier ; les opérations sur les applications (somme, composée, réciproque) se traduisent exactement par les opérations correspondantes sur leurs matrices.
\end{itemize}

Tu as ainsi mobilisé l'ensemble des savoir-faire de la leçon — présentation, opérations, déterminant, inversion, lien avec les systèmes et applications linéaires — dans une situation authentique de la vie yaoundéenne. La prochaine fois que tu entreras dans une cybercafétéria, tu verras peut-être des matrices partout !

\begin{savaistu}[Et à plus grande échelle ?]
Les grandes entreprises de téléphonie mobile au Cameroun (MTN, Orange) gèrent chaque jour des millions de transactions (forfaits voix, data, Mobile Money) rangées dans d'immenses tableaux de données. Les algorithmes qui calculent les facturations, détectent les fraudes ou recommandent des offres reposent précisément sur les opérations matricielles que tu viens de pratiquer : sommes, produits, inversions. Les mêmes mathématiques, à une autre échelle !
\end{savaistu}

\newpage

\coursec{Méthodes et exercices résolus}

\coursec{Matrices et applications linéaires d'un plan vectoriel}

\courssub{Présentation et vocabulaire}

\begin{definition}[Matrice]
Une \cle{matrice} est un tableau rectangulaire de nombres (réels le plus souvent) organisés en \textbf{lignes} et en \textbf{colonnes}.
Une matrice à $m$ lignes et $n$ colonnes est de \textbf{format} $(m,n)$ ou de \textbf{type} $m\times n$.
On note $a_{ij}$ le coefficient situé à la $i^{\text{ème}}$ ligne et à la $j^{\text{ème}}$ colonne.
L'ensemble des matrices de format $(m,n)$ à coefficients dans $\mathbb{R}$ se note $\mathcal{M}_{m,n}(\mathbb{R})$.
Une matrice de format $(n,n)$ est dite \textbf{carrée d'ordre $n$}.
\end{definition}

\begin{propriete}[Matrices particulières]
\begin{itemize}
\item \textbf{Matrice nulle} $O_{m,n}$ : tous les coefficients sont nuls.
\item \textbf{Matrice identité} $I_n$ (d'ordre $n$) : matrice carrée avec des $1$ sur la diagonale principale et des $0$ ailleurs.
\[ I_2=\begin{pmatrix} 1 & 0 \\ 0 & 1 \end{pmatrix},\qquad I_3=\begin{pmatrix} 1 & 0 & 0 \\ 0 & 1 & 0 \\ 0 & 0 & 1 \end{pmatrix}. \]
\item \textbf{Matrice transposée} : la transposée de $A=(a_{ij})$ est ${}^t\!A=(a_{ji})$ (lignes deviennent colonnes).
\end{itemize}
\end{propriete}

\courssub{Traduire la présentation des données par une matrice}

\begin{methode}[Organiser des données en matrice]
Pour traduire une situation concrète par une matrice :
\begin{enumerate}
\item \cle{Identifier} les deux « dimensions » du problème : ce qui sera en \textbf{lignes} (ex. : magasins, jours, élèves) et ce qui sera en \textbf{colonnes} (ex. : produits, matières, notes).
\item \cle{Fixer l'ordre} de lecture une fois pour toutes et l'annoncer dans la rédaction : « la ligne $i$ correspond à..., la colonne $j$ correspond à... ».
\item \cle{Remplir} les coefficients : le nombre situé à la ligne $i$ et à la colonne $j$ est noté $a_{ij}$.
\item \cle{Vérifier} la cohérence : la matrice a autant de lignes que de catégories du premier type, autant de colonnes que de catégories du second type. Annoncer son format $(m,n)$.
\end{enumerate}
\textbf{Réflexe Probatoire :} annoncer toujours le format de la matrice obtenue et la signification d'au moins un coefficient, par exemple $a_{21}$.
\end{methode}

\begin{exoresolu}[Stocks de deux boutiques]
Une entreprise de vente de crédit téléphonique possède deux boutiques, $B_1$ à Yaoundé et $B_2$ à Douala. La boutique $B_1$ dispose de $150$ cartes MTN et de $230$ cartes Orange ; la boutique $B_2$ dispose de $180$ cartes MTN et de $95$ cartes Orange.
\begin{enumerate}
\item Présenter ces données par une matrice $A$ dont les lignes correspondent aux boutiques et les colonnes aux opérateurs (MTN puis Orange).
\item Que représente le coefficient $a_{21}$ ? Quel est le format de $A$ ?
\end{enumerate}
\tcblower
\begin{enumerate}
\item On fixe l'ordre : ligne 1 $\leftrightarrow$ boutique $B_1$, ligne 2 $\leftrightarrow$ boutique $B_2$ ; colonne 1 $\leftrightarrow$ MTN, colonne 2 $\leftrightarrow$ Orange. On obtient :
\[ A=\begin{pmatrix} 150 & 230 \\ 180 & 95 \end{pmatrix}. \]
\item Le coefficient $a_{21}$ est situé à la $2^{\text{ème}}$ ligne et à la $1^{\text{ère}}$ colonne : il représente le nombre de cartes MTN en stock dans la boutique $B_2$ de Douala, soit $a_{21}=180$.
La matrice $A$ possède $2$ lignes et $2$ colonnes : c'est une matrice \textbf{carrée d'ordre $2$} (format $(2,2)$).
\end{enumerate}
\textbf{Conclusion :} $A=\begin{pmatrix} 150 & 230 \\ 180 & 95 \end{pmatrix}$ est une matrice carrée d'ordre $2$, et $a_{21}=180$ désigne le stock de cartes MTN à Douala.
\end{exoresolu}

\courssub{Opérations sur les matrices : somme, produit par un réel, produit}

\begin{methode}[Effectuer les opérations sur les matrices]
Soient $A=(a_{ij})$ et $B=(b_{ij})$ deux matrices de même format $(m,n)$, et $k\in\mathbb{R}$.
\begin{enumerate}
\item \textbf{Somme} $A+B$ : on additionne \cle{coefficient par coefficient}. Condition : $A$ et $B$ ont le \textbf{même format}.
\[ A+B=(a_{ij}+b_{ij})_{1\le i\le m,\,1\le j\le n}. \]
\item \textbf{Produit par un réel} $kA$ : on multiplie \cle{chaque coefficient} par $k$.
\[ kA=(k\,a_{ij})_{1\le i\le m,\,1\le j\le n}. \]
\item \textbf{Produit} $A\times B$ (ou $AB$) : soit $A\in\mathcal{M}_{m,p}(\mathbb{R})$ et $B\in\mathcal{M}_{p,n}(\mathbb{R})$. Le produit $C=AB$ est une matrice de format $(m,n)$ dont le coefficient $c_{ij}$ est le \cle{produit scalaire} de la ligne $i$ de $A$ par la colonne $j$ de $B$ :
\[ c_{ij}=\sum_{k=1}^{p} a_{ik}b_{kj}. \]
Condition impérative : le nombre de \textbf{colonnes de $A$} égale le nombre de \textbf{lignes de $B$}.
Pour des matrices carrées d'ordre $2$ :
\[ \begin{pmatrix} a & b \\ c & d \end{pmatrix}\begin{pmatrix} a' & b' \\ c' & d' \end{pmatrix}=\begin{pmatrix} aa'+bc' & ab'+bd' \\ ca'+dc' & cb'+dd' \end{pmatrix}. \]
\end{enumerate}
\textbf{Propriétés (admises) :} associativité $(AB)C=A(BC)$, distributivité $A(B+C)=AB+AC$, $(A+B)C=AC+BC$, $I_n$ neutre pour le produit.
\textbf{Avertissement :} le produit de matrices n'est \textbf{pas commutatif} : en général $AB\neq BA$.
\textbf{Réflexe :} pour vérifier un calcul, refaire mentalement un seul coefficient, par exemple $c_{11}$.
\end{methode}

\begin{exoresolu}[Opérations sur deux matrices]
On donne $A=\begin{pmatrix} 2 & -1 \\ 3 & 4 \end{pmatrix}$ et $B=\begin{pmatrix} 1 & 5 \\ -2 & 0 \end{pmatrix}$.
Calculer $A+B$, $3A$, $A\times B$ et $B\times A$. Que constate-t-on ?
\tcblower
\textbf{Somme} (coefficient par coefficient) :
\[ A+B=\begin{pmatrix} 2+1 & -1+5 \\ 3+(-2) & 4+0 \end{pmatrix}=\begin{pmatrix} 3 & 4 \\ 1 & 4 \end{pmatrix}. \]
\textbf{Produit par le réel $3$} (chaque coefficient est multiplié par $3$) :
\[ 3A=\begin{pmatrix} 3\times 2 & 3\times (-1) \\ 3\times 3 & 3\times 4 \end{pmatrix}=\begin{pmatrix} 6 & -3 \\ 9 & 12 \end{pmatrix}. \]
\textbf{Produit $A\times B$} (ligne de $A$ contre colonne de $B$) :
\begin{align*}
A\times B &=\begin{pmatrix} 2\times 1+(-1)\times(-2) & 2\times 5+(-1)\times 0 \\ 3\times 1+4\times(-2) & 3\times 5+4\times 0 \end{pmatrix}\\
&=\begin{pmatrix} 2+2 & 10+0 \\ 3-8 & 15+0 \end{pmatrix}=\begin{pmatrix} 4 & 10 \\ -5 & 15 \end{pmatrix}.
\end{align*}
\textbf{Produit $B\times A$} :
\begin{align*}
B\times A &=\begin{pmatrix} 1\times 2+5\times 3 & 1\times(-1)+5\times 4 \\ -2\times 2+0\times 3 & -2\times(-1)+0\times 4 \end{pmatrix}\\
&=\begin{pmatrix} 17 & 19 \\ -4 & 2 \end{pmatrix}.
\end{align*}
\textbf{Constat :} $A\times B=\begin{pmatrix} 4 & 10 \\ -5 & 15 \end{pmatrix}\neq\begin{pmatrix} 17 & 19 \\ -4 & 2 \end{pmatrix}=B\times A$ : le produit matriciel \textbf{n'est pas commutatif}.
\end{exoresolu}

\courssub{Déterminant et inverse d'une matrice carrée d'ordre 2}

\begin{methode}[Calculer un déterminant et inverser une matrice $2\times 2$]
Soit $A=\begin{pmatrix} a & b \\ c & d \end{pmatrix}$.
\begin{enumerate}
\item \textbf{Déterminant} : $\det(A)=ad-bc$ (produit des coefficients de la diagonale principale \textbf{moins} produit des coefficients de l'autre diagonale).
\item \textbf{Test d'inversibilité} : $A$ est inversible \cle{si et seulement si} $\det(A)\neq 0$.
\item \textbf{Formule de l'inverse} : si $\det(A)\neq 0$,
\[ A^{-1}=\frac{1}{\det(A)}\begin{pmatrix} d & -b \\ -c & a \end{pmatrix}. \]
Moyen mnémotechnique : on \textbf{échange} $a$ et $d$, on change le \textbf{signe} de $b$ et $c$, et on divise tout par le déterminant.
\item \textbf{Vérification} : contrôler que $A\times A^{-1}=I_2$ où $I_2=\begin{pmatrix} 1 & 0 \\ 0 & 1 \end{pmatrix}$.
\end{enumerate}
\end{methode}

\begin{exoresolu}[Inversibilité et calcul de l'inverse]
On considère $M=\begin{pmatrix} 3 & 2 \\ 5 & 4 \end{pmatrix}$ et $N=\begin{pmatrix} 2 & 4 \\ 1 & 2 \end{pmatrix}$.
\begin{enumerate}
\item Montrer que $M$ est inversible et déterminer $M^{-1}$.
\item La matrice $N$ est-elle inversible ? Justifier.
\item Vérifier que $M\times M^{-1}=I_2$.
\end{enumerate}
\tcblower
\begin{enumerate}
\item $\det(M)=3\times 4-2\times 5=12-10=2$. Comme $\det(M)=2\neq 0$, la matrice $M$ \textbf{est inversible} et :
\[ M^{-1}=\frac{1}{2}\begin{pmatrix} 4 & -2 \\ -5 & 3 \end{pmatrix}=\begin{pmatrix} 2 & -1 \\ -\dfrac{5}{2} & \dfrac{3}{2} \end{pmatrix}. \]
\item $\det(N)=2\times 2-4\times 1=4-4=0$. Comme $\det(N)=0$, la matrice $N$ \textbf{n'est pas inversible}.
\item Vérification :
\begin{align*}
M\times M^{-1}&=\begin{pmatrix} 3 & 2 \\ 5 & 4 \end{pmatrix}\times\frac{1}{2}\begin{pmatrix} 4 & -2 \\ -5 & 3 \end{pmatrix}\\
&=\frac{1}{2}\begin{pmatrix} 3\times 4+2\times(-5) & 3\times(-2)+2\times 3 \\ 5\times 4+4\times(-5) & 5\times(-2)+4\times 3 \end{pmatrix}\\
&=\frac{1}{2}\begin{pmatrix} 12-10 & -6+6 \\ 20-20 & -10+12 \end{pmatrix}=\frac{1}{2}\begin{pmatrix} 2 & 0 \\ 0 & 2 \end{pmatrix}=\begin{pmatrix} 1 & 0 \\ 0 & 1 \end{pmatrix}=I_2.
\end{align*}
\end{enumerate}
\textbf{Conclusion :} $M$ est inversible, $M^{-1}=\begin{pmatrix} 2 & -1 \\ -\frac{5}{2} & \frac{3}{2} \end{pmatrix}$, et $N$ n'est pas inversible car son déterminant est nul.
\end{exoresolu}

\courssub{Résolution de systèmes linéaires par les matrices}

\begin{methode}[Résoudre un système de deux équations à deux inconnues]
Pour résoudre $\begin{cases} ax+by=e \\ cx+dy=f \end{cases}$ :
\begin{enumerate}
\item \cle{Écrire} le système sous forme matricielle $AX=B$ avec $A=\begin{pmatrix} a & b \\ c & d \end{pmatrix}$, $X=\begin{pmatrix} x \\ y \end{pmatrix}$, $B=\begin{pmatrix} e \\ f \end{pmatrix}$.
\item \cle{Calculer} $\det(A)=ad-bc$.
\item Si $\det(A)\neq 0$ : le système admet un \textbf{unique couple solution}, donné par $X=A^{-1}B$.
\item Si $\det(A)=0$ : le système n'a pas de solution unique (aucune solution ou une infinité) ; on revient à la méthode par combinaison pour trancher.
\item \cle{Vérifier} le couple trouvé en le substituant dans les deux équations de départ.
\end{enumerate}
\end{methode}

\begin{exoresolu}[Problème de marché à Mokolo]
Au marché Mokolo, $3$ kg de tomates et $2$ kg d'oignons coûtent $\SI{2500}{\text{F}}$ ; $5$ kg de tomates et $4$ kg d'oignons coûtent $\SI{4500}{\text{F}}$. On note $x$ le prix d'un kg de tomates et $y$ celui d'un kg d'oignons (en francs).
\begin{enumerate}
\item Traduire la situation par un système, puis par une équation matricielle $AX=B$.
\item Résoudre le système en calculant $A^{-1}$.
\end{enumerate}
\tcblower
\begin{enumerate}
\item Le système est $\begin{cases} 3x+2y=2500 \\ 5x+4y=4500 \end{cases}$, soit matriciellement :
\[ \underbrace{\begin{pmatrix} 3 & 2 \\ 5 & 4 \end{pmatrix}}_{A}\underbrace{\begin{pmatrix} x \\ y \end{pmatrix}}_{X}=\underbrace{\begin{pmatrix} 2500 \\ 4500 \end{pmatrix}}_{B}. \]
\item $\det(A)=3\times 4-2\times 5=2\neq 0$, donc $A$ est inversible et $A^{-1}=\dfrac{1}{2}\begin{pmatrix} 4 & -2 \\ -5 & 3 \end{pmatrix}$. Alors :
\begin{align*}
X=A^{-1}B&=\frac{1}{2}\begin{pmatrix} 4 & -2 \\ -5 & 3 \end{pmatrix}\begin{pmatrix} 2500 \\ 4500 \end{pmatrix}\\
&=\frac{1}{2}\begin{pmatrix} 4\times 2500-2\times 4500 \\ -5\times 2500+3\times 4500 \end{pmatrix}=\frac{1}{2}\begin{pmatrix} 10000-9000 \\ -12500+13500 \end{pmatrix}=\frac{1}{2}\begin{pmatrix} 1000 \\ 1000 \end{pmatrix}=\begin{pmatrix} 500 \\ 500 \end{pmatrix}.
\end{align*}
\textbf{Vérification :} $3\times 500+2\times 500=1500+1000=2500$ \checkmark\ et $5\times 500+4\times 500=2500+2000=4500$ \checkmark.
\end{enumerate}
\textbf{Conclusion :} le kilogramme de tomates coûte $\SI{500}{\text{F}}$ et celui d'oignons $\SI{500}{\text{F}}$.
\end{exoresolu}

\courssub{Matrice d'une application linéaire dans une base}

\begin{definition}[Application linéaire]
Une application $f:E\to E$ (où $E$ est un plan vectoriel) est \textbf{linéaire} si pour tous vecteurs $\vec{u},\vec{v}\in E$ et tout réel $k$ :
\[ f(\vec{u}+\vec{v})=f(\vec{u})+f(\vec{v}) \quad\text{et}\quad f(k\vec{u})=k\,f(\vec{u}). \]
Conséquence immédiate : $f(\vec{0})=\vec{0}$ et $f(\vec{u}-\vec{v})=f(\vec{u})-f(\vec{v})$.
\end{definition}

\begin{methode}[Construire la matrice d'une application linéaire]
Soit $f$ une application linéaire du plan vectoriel $E$ dans lui-même et $\mathcal{B}=(\vec{i},\vec{j})$ une base de $E$.
\begin{enumerate}
\item \cle{Calculer} les images $f(\vec{i})$ et $f(\vec{j})$.
\item \cle{Décomposer} chaque image dans la base $\mathcal{B}$ : $f(\vec{i})=a\vec{i}+c\vec{j}$ et $f(\vec{j})=b\vec{i}+d\vec{j}$.
\item \cle{Écrire} la matrice en plaçant les coordonnées \textbf{en colonnes} : la première colonne contient les coordonnées de $f(\vec{i})$, la deuxième celles de $f(\vec{j})$ :
\[ \mathrm{Mat}_{\mathcal{B}}(f)=\begin{pmatrix} a & b \\ c & d \end{pmatrix}. \]
\end{enumerate}
\textbf{Phrase de rédaction à recopier :} « Les colonnes de la matrice de $f$ dans $\mathcal{B}$ sont les coordonnées dans $\mathcal{B}$ des images des vecteurs de $\mathcal{B}$. »
\end{methode}

\begin{exoresolu}[Matrice d'une application linéaire]
Le plan est muni d'une base $\mathcal{B}=(\vec{i},\vec{j})$. Soit $f$ l'application linéaire définie par :
\[ f(\vec{i})=2\vec{i}-\vec{j} \qquad \text{et} \qquad f(\vec{j})=3\vec{i}+4\vec{j}. \]
\begin{enumerate}
\item Écrire la matrice $A$ de $f$ dans la base $\mathcal{B}$.
\item Soit $\vec{u}=\vec{i}+2\vec{j}$. Déterminer $f(\vec{u})$ par linéarité, puis retrouver ses coordonnées par un produit matriciel.
\end{enumerate}
\tcblower
\begin{enumerate}
\item Les coordonnées de $f(\vec{i})$ dans $\mathcal{B}$ sont $(2\,;-1)$ : c'est la \textbf{première colonne}. Celles de $f(\vec{j})$ sont $(3\,;4)$ : c'est la \textbf{deuxième colonne}. Donc :
\[ A=\mathrm{Mat}_{\mathcal{B}}(f)=\begin{pmatrix} 2 & 3 \\ -1 & 4 \end{pmatrix}. \]
\item \textbf{Par linéarité :} $f(\vec{u})=f(\vec{i}+2\vec{j})=f(\vec{i})+2f(\vec{j})=(2\vec{i}-\vec{j})+2(3\vec{i}+4\vec{j})=2\vec{i}-\vec{j}+6\vec{i}+8\vec{j}=8\vec{i}+7\vec{j}$.

\textbf{Par produit matriciel :} les coordonnées de $\vec{u}$ sont $X=\begin{pmatrix} 1 \\ 2 \end{pmatrix}$. Celles de $f(\vec{u})$ sont :
\[ AX=\begin{pmatrix} 2 & 3 \\ -1 & 4 \end{pmatrix}\begin{pmatrix} 1 \\ 2 \end{pmatrix}=\begin{pmatrix} 2\times 1+3\times 2 \\ -1\times 1+4\times 2 \end{pmatrix}=\begin{pmatrix} 8 \\ 7 \end{pmatrix}. \]
On retrouve bien $f(\vec{u})=8\vec{i}+7\vec{j}$.
\end{enumerate}
\textbf{Conclusion :} $A=\begin{pmatrix} 2 & 3 \\ -1 & 4 \end{pmatrix}$ et $f(\vec{u})$ a pour coordonnées $(8\,;7)$ dans $\mathcal{B}$, obtenues par les deux méthodes.
\end{exoresolu}

\courssub{Matrices des transformations élémentaires du plan}

\begin{savaistu}[Lien avec la géométrie : configurations et transformations]
Toute transformation géométrique du plan qui conserve l'origine, le parallélisme et les rapports de longueurs (homotheties, rotations, symétries orthogonales, projections orthogonales) est une \textbf{application linéaire}. Sa matrice dans la base orthonormée directe $(\vec{i},\vec{j})$ se construit exactement par la méthode ci-dessus. C'est le cœur de la famille de situations « Représentations et transformations des configurations planes ».
\end{savaistu}

\begin{propriete}[Matrices usuelles dans une base orthonormée directe $(\vec{i},\vec{j})$]
\begin{center}
\begin{tabularx}{\linewidth}{|l|X|X|}\hline
\textbf{Transformation} & \textbf{Description} & \textbf{Matrice} \\ \hline
Hométhétie $h_{O,k}$ & Rapport $k$ & $\begin{pmatrix} k & 0 \\ 0 & k \end{pmatrix}=kI_2$ \\ \hline
Rotation $r_{O,\theta}$ & Angle $\theta$ (trigonométrique) & $\begin{pmatrix} \cos\theta & -\sin\theta \\ \sin\theta & \cos\theta \end{pmatrix}$ \\ \hline
Symétrie centrale $s_O$ & Centre $O$ (rotation $\pi$) & $\begin{pmatrix} -1 & 0 \\ 0 & -1 \end{pmatrix}=-I_2$ \\ \hline
Symétrie axiale $s_{(Ox)}$ & Axe $(Ox)$ & $\begin{pmatrix} 1 & 0 \\ 0 & -1 \end{pmatrix}$ \\ \hline
Symétrie axiale $s_{(Oy)}$ & Axe $(Oy)$ & $\begin{pmatrix} -1 & 0 \\ 0 & 1 \end{pmatrix}$ \\ \hline
Symétrie axiale $s_{y=x}$ & Axe droite $y=x$ & $\begin{pmatrix} 0 & 1 \\ 1 & 0 \end{pmatrix}$ \\ \hline
Projection orthogonale sur $(Ox)$ & & $\begin{pmatrix} 1 & 0 \\ 0 & 0 \end{pmatrix}$ \\ \hline
Projection orthogonale sur $(Oy)$ & & $\begin{pmatrix} 0 & 0 \\ 0 & 1 \end{pmatrix}$ \\ \hline
\end{tabularx}
\end{center}
\textbf{À retenir :} Pour une rotation de $90^\circ$, $\cos90^\circ=0$, $\sin90^\circ=1$, matrice $\begin{pmatrix} 0 & -1 \\ 1 & 0 \end{pmatrix}$.
\end{propriete}

\begin{experience}[Découverte : matrice d'une rotation de $90^\circ$]
Dans la base orthonormée directe $(\vec{i},\vec{j})$, soit $r$ la rotation de centre $O$ et d'angle $90^\circ$.
\begin{enumerate}
\item Déterminer $r(\vec{i})$ et $r(\vec{j})$ par lecture graphique.
\item En déduire la matrice de $r$.
\item Vérifier que $r^2$ (deux rotations successives) est la symétrie centrale $s_O$.
\end{enumerate}
\tcblower
\begin{enumerate}
\item Graphiquement, $r(\vec{i})=\vec{j}$ (coordonnées $(0,1)$) et $r(\vec{j})=-\vec{i}$ (coordonnées $(-1,0)$).
\item La matrice est donc $\begin{pmatrix} 0 & -1 \\ 1 & 0 \end{pmatrix}$ (première colonne : image de $\vec{i}$, deuxième : image de $\vec{j}$).
\item Matrice de $r^2$ : $\begin{pmatrix} 0 & -1 \\ 1 & 0 \end{pmatrix}^2 = \begin{pmatrix} -1 & 0 \\ 0 & -1 \end{pmatrix} = -I_2$, ce qui est bien la matrice de la symétrie centrale $s_O$.
\end{enumerate}
\end{experience}

\begin{popfigure}
\begin{tikzpicture}[>=Stealth, scale=1.8]
% Axes
\draw[->, thin, popmuted] (-1.5,0) -- (1.5,0) node[right] {$x$};
\draw[->, thin, popmuted] (0,-1.5) -- (0,1.5) node[above] {$y$};
% Base vectors
\draw[->, very thick, popblue] (0,0) -- (1,0) node[below right] {$\vec{i}$};
\draw[->, very thick, popblue] (0,0) -- (0,1) node[above left] {$\vec{j}$};
% Rotated vectors
\draw[->, very thick, poporange, dashed] (0,0) -- (0,1) node[above left] {$r(\vec{i})=\vec{j}$};
\draw[->, very thick, poporange, dashed] (0,0) -- (-1,0) node[above left] {$r(\vec{j})=-\vec{i}$};
% Unit square original
\draw[popblue, thin, fill=popblue!10] (0,0) -- (1,0) -- (1,1) -- (0,1) -- cycle;
% Unit square rotated
\draw[poporange, thin, dashed, fill=poporange!10] (0,0) -- (0,1) -- (-1,1) -- (-1,0) -- cycle;
% Angle mark
\draw[poppurple, thick] (0.3,0) arc (0:90:0.3) node[midway, right=2pt] {$90^\circ$};
\end{tikzpicture}
\legende{Effet de la rotation de $90^\circ$ sur la base $(\vec{i},\vec{j})$ et le carré unité. Les colonnes de la matrice sont les images des vecteurs de la base.}
\end{popfigure}

\courssub{Opérations sur les applications linéaires et matrices associées}

\begin{methode}[Opérations sur les applications linéaires et leurs matrices]
Soit $f$ et $g$ deux applications linéaires du plan $E$, de matrices respectives $A$ et $B$ dans une même base $\mathcal{B}$, et $k$ un réel.
\begin{enumerate}
\item \textbf{Somme} : $\mathrm{Mat}_{\mathcal{B}}(f+g)=A+B$.
\item \textbf{Produit par un réel} : $\mathrm{Mat}_{\mathcal{B}}(k\,f)=kA$.
\item \textbf{Composée} : $\mathrm{Mat}_{\mathcal{B}}(g\circ f)=B\times A$. \cle{Attention à l'ordre} : la matrice de l'application appliquée \textbf{en premier} ($f$) est écrite \textbf{à droite} dans le produit.
\end{enumerate}
\textbf{Réflexe :} pour calculer $g\circ f$, on applique d'abord $f$ puis $g$ ; côté matrices, on calcule $B\times A$ (et non $A\times B$).
\end{methode}

\begin{exoresolu}[Somme et composée d'applications linéaires]
Dans une base $\mathcal{B}=(\vec{i},\vec{j})$, les applications linéaires $f$ et $g$ ont pour matrices :
\[ A=\begin{pmatrix} 1 & 2 \\ 0 & -1 \end{pmatrix} \qquad \text{et} \qquad B=\begin{pmatrix} 2 & 1 \\ 1 & 3 \end{pmatrix}. \]
\begin{enumerate}
\item Déterminer la matrice de $f+g$ et celle de $2f$.
\item Déterminer la matrice de $g\circ f$, puis celle de $f\circ g$. Conclure.
\end{enumerate}
\tcblower
\begin{enumerate}
\item $\mathrm{Mat}_{\mathcal{B}}(f+g)=A+B=\begin{pmatrix} 1+2 & 2+1 \\ 0+1 & -1+3 \end{pmatrix}=\begin{pmatrix} 3 & 3 \\ 1 & 2 \end{pmatrix}$.

$\mathrm{Mat}_{\mathcal{B}}(2f)=2A=\begin{pmatrix} 2 & 4 \\ 0 & -2 \end{pmatrix}$.
\item $\mathrm{Mat}_{\mathcal{B}}(g\circ f)=B\times A$ :
\begin{align*}
B\times A&=\begin{pmatrix} 2 & 1 \\ 1 & 3 \end{pmatrix}\begin{pmatrix} 1 & 2 \\ 0 & -1 \end{pmatrix}=\begin{pmatrix} 2\times 1+1\times 0 & 2\times 2+1\times(-1) \\ 1\times 1+3\times 0 & 1\times 2+3\times(-1) \end{pmatrix}=\begin{pmatrix} 2 & 3 \\ 1 & -1 \end{pmatrix}.
\end{align*}
$\mathrm{Mat}_{\mathcal{B}}(f\circ g)=A\times B$ :
\begin{align*}
A\times B&=\begin{pmatrix} 1 & 2 \\ 0 & -1 \end{pmatrix}\begin{pmatrix} 2 & 1 \\ 1 & 3 \end{pmatrix}=\begin{pmatrix} 1\times 2+2\times 1 & 1\times 1+2\times 3 \\ 0\times 2+(-1)\times 1 & 0\times 1+(-1)\times 3 \end{pmatrix}=\begin{pmatrix} 4 & 7 \\ -1 & -3 \end{pmatrix}.
\end{align*}
Comme $B\times A\neq A\times B$, on a $g\circ f\neq f\circ g$ : \textbf{la composition des applications linéaires n'est pas commutative}.
\end{enumerate}
\textbf{Conclusion :} $\mathrm{Mat}(f+g)=\begin{pmatrix} 3 & 3 \\ 1 & 2 \end{pmatrix}$, $\mathrm{Mat}(2f)=\begin{pmatrix} 2 & 4 \\ 0 & -2 \end{pmatrix}$, $\mathrm{Mat}(g\circ f)=\begin{pmatrix} 2 & 3 \\ 1 & -1 \end{pmatrix}$ et $\mathrm{Mat}(f\circ g)=\begin{pmatrix} 4 & 7 \\ -1 & -3 \end{pmatrix}$.
\end{exoresolu}

\courssub{Matrice de la réciproque d'un automorphisme}

\begin{methode}[Automorphisme et matrice de la réciproque]
Un \cle{automorphisme} du plan est une application linéaire \textbf{bijective} de $E$ dans $E$.
\begin{enumerate}
\item \textbf{Test} : $f$ est un automorphisme si et seulement si sa matrice $A$ dans une base $\mathcal{B}$ est inversible, c'est-à-dire $\det(A)\neq 0$.
\item \textbf{Matrice de la réciproque} : si $f$ est un automorphisme, alors $f^{-1}$ est aussi un automorphisme et
\[ \mathrm{Mat}_{\mathcal{B}}(f^{-1})=A^{-1}. \]
\item \textbf{Contrôle} : vérifier que $A\times A^{-1}=I_2$, ce qui traduit $f\circ f^{-1}=\mathrm{Id}_E$.
\end{enumerate}
\end{methode}

\begin{exoresolu}[Réciproque d'un automorphisme]
Dans la base $\mathcal{B}=(\vec{i},\vec{j})$, l'application linéaire $f$ a pour matrice $A=\begin{pmatrix} 2 & 1 \\ 1 & 1 \end{pmatrix}$.
\begin{enumerate}
\item Montrer que $f$ est un automorphisme du plan.
\item Déterminer la matrice de $f^{-1}$ dans $\mathcal{B}$.
\item Un vecteur $\vec{v}$ a pour image par $f$ le vecteur de coordonnées $\begin{pmatrix} 5 \\ 3 \end{pmatrix}$. Déterminer les coordonnées de $\vec{v}$.
\end{enumerate}
\tcblower
\begin{enumerate}
\item $\det(A)=2\times 1-1\times 1=2-1=1\neq 0$. La matrice $A$ est donc inversible, ce qui prouve que $f$ est \textbf{un automorphisme} du plan.
\item La matrice de $f^{-1}$ est l'inverse de $A$ :
\[ A^{-1}=\frac{1}{\det(A)}\begin{pmatrix} 1 & -1 \\ -1 & 2 \end{pmatrix}=\begin{pmatrix} 1 & -1 \\ -1 & 2 \end{pmatrix}. \]
Vérification : $A\times A^{-1}=\begin{pmatrix} 2 & 1 \\ 1 & 1 \end{pmatrix}\begin{pmatrix} 1 & -1 \\ -1 & 2 \end{pmatrix}=\begin{pmatrix} 2-1 & -2+2 \\ 1-1 & -1+2 \end{pmatrix}=\begin{pmatrix} 1 & 0 \\ 0 & 1 \end{pmatrix}=I_2$ \checkmark.
\item Notons $X$ la colonne des coordonnées de $\vec{v}$. On a $AX=\begin{pmatrix} 5 \\ 3 \end{pmatrix}$, donc :
\[ X=A^{-1}\begin{pmatrix} 5 \\ 3 \end{pmatrix}=\begin{pmatrix} 1 & -1 \\ -1 & 2 \end{pmatrix}\begin{pmatrix} 5 \\ 3 \end{pmatrix}=\begin{pmatrix} 5-3 \\ -5+6 \end{pmatrix}=\begin{pmatrix} 2 \\ 1 \end{pmatrix}. \]
\end{enumerate}
\textbf{Conclusion :} $f$ est un automorphisme, $\mathrm{Mat}_{\mathcal{B}}(f^{-1})=\begin{pmatrix} 1 & -1 \\ -1 & 2 \end{pmatrix}$, et $\vec{v}=2\vec{i}+\vec{j}$.
\end{exoresolu}

\begin{attention}[Les 5 erreurs qui coûtent des points au Probatoire]
\begin{enumerate}
\item \textbf{Additionner des matrices de formats différents.} La somme $A+B$ n'existe que si $A$ et $B$ ont le même format ; le produit $A\times B$ n'existe que si le nombre de colonnes de $A$ égale le nombre de lignes de $B$. Toujours annoncer les formats avant de calculer.
\item \textbf{Croire que $A\times B=B\times A$.} Le produit matriciel n'est pas commutatif : dans la matrice d'une composée, $\mathrm{Mat}(g\circ f)=B\times A$, la matrice de l'application appliquée \emph{en premier} se met \emph{à droite}. Inverser l'ordre est l'erreur la plus fréquente.
\item \textbf{Se tromper de signe dans le déterminant.} $\det(A)=ad-bc$ : c'est un \textbf{moins} entre les deux produits diagonaux. Écrire $ad+bc$ fausse le test d'inversibilité et tout l'exercice.
\item \textbf{Inverser une matrice sans vérifier que $\det(A)\neq 0$.} La condition $\det(A)\neq 0$ doit être \emph{écrite et justifiée} avant d'appliquer la formule $A^{-1}=\dfrac{1}{\det(A)}\begin{pmatrix} d & -b \\ -c & a \end{pmatrix}$ ; de plus, ne pas oublier le facteur $\dfrac{1}{\det(A)}$ devant la matrice.
\item \textbf{Placer les images en lignes au lieu de colonnes.} Dans la matrice d'une application linéaire, les coordonnées de $f(\vec{i})$ et $f(\vec{j})$ s'écrivent \textbf{en colonnes}, pas en lignes. Une matrice transposée par erreur donne des images fausses dans toutes les questions suivantes.
\end{enumerate}
\end{attention}

\newpage

\coursec{Exercices}

\courssub{Je vérifie mes connaissances}

\begin{exercice}[QCM : une seule réponse exacte]
Pour chaque question, choisir la bonne réponse (aucune justification n'est demandée ici).

\textbf{1.} Soit $A = \begin{pmatrix} 2 & -1 \\ 3 & 4 \end{pmatrix}$. Le déterminant de $A$ vaut :
\begin{itemize}
\item[a)] $5$
\item[b)] $11$
\item[c)] $-11$
\item[d)] $8$
\end{itemize}

\textbf{2.} La matrice $M = \begin{pmatrix} 1 & 2 \\ 3 & 6 \end{pmatrix}$ est :
\begin{itemize}
\item[a)] inversible
\item[b)] non inversible
\item[c)] égale à l'identité
\item[d)] nulle
\end{itemize}

\textbf{3.} Si $A$ et $B$ sont deux matrices carrées d'ordre $2$, alors en général :
\begin{itemize}
\item[a)] $AB = BA$
\item[b)] $AB \neq BA$
\item[c)] $AB = A + B$
\item[d)] $AB$ n'existe pas
\end{itemize}

\textbf{4.} Si $f$ est un automorphisme du plan de matrice $A$ dans une base, alors la matrice de $f^{-1}$ dans cette base est :
\begin{itemize}
\item[a)] $A$
\item[b)] $-A$
\item[c)] $A^{-1}$
\item[d)] $2A$
\end{itemize}
\end{exercice}

\begin{exercice}[Vrai ou faux, en justifiant]
Répondre par \textbf{vrai} ou \textbf{faux} en justifiant chaque réponse par un argument ou un contre-exemple.

\textbf{1.} La somme de deux matrices carrées d'ordre $2$ est une matrice carrée d'ordre $2$.

\textbf{2.} Pour toute matrice carrée $A$ d'ordre $2$, on a $\det(-A) = -\det(A)$.

\textbf{3.} Une matrice carrée d'ordre $2$ dont le déterminant est nul admet un inverse.

\textbf{4.} Si $A = \begin{pmatrix} 1 & 1 \\ 0 & 1 \end{pmatrix}$ et $I = \begin{pmatrix} 1 & 0 \\ 0 & 1 \end{pmatrix}$, alors $AI = A$.
\end{exercice}

\begin{exercice}[Définitions et complétion]
\textbf{1.} Donner la définition du déterminant d'une matrice carrée d'ordre $2$, $A = \begin{pmatrix} a & b \\ c & d \end{pmatrix}$.

\textbf{2.} Compléter : une matrice carrée $A$ d'ordre $2$ est inversible si et seulement si $\det(A) \neq 0$ ; dans ce cas $A^{-1} = \dfrac{1}{\det(A)}\begin{pmatrix} d & -b \\ -c & a \end{pmatrix}$.

\textbf{3.} Compléter : la matrice d'une application linéaire $f$ dans une base $(\vec{\imath}, \vec{\jmath})$ a pour colonnes les coordonnées des vecteurs $f(\vec{\imath})$ et $f(\vec{\jmath})$ dans cette base.

\textbf{4.} Qu'appelle-t-on \cle{automorphisme} d'un plan vectoriel ?
\end{exercice}

\begin{exercice}[Calculs directs]
Soient $A = \begin{pmatrix} 2 & -1 \\ 3 & 4 \end{pmatrix}$ et $B = \begin{pmatrix} 1 & 0 \\ -2 & 5 \end{pmatrix}$.
\begin{enumerate}
\item Calculer $A + B$ et $3A$.
\item Calculer $\det(A)$ et $\det(B)$.
\item La matrice $A$ est-elle inversible ? Si oui, calculer $A^{-1}$.
\end{enumerate}
\end{exercice}

\courssub{Je m'entraîne}

\begin{exercice}[Produits de matrices]
On reprend les matrices $A = \begin{pmatrix} 2 & -1 \\ 3 & 4 \end{pmatrix}$ et $B = \begin{pmatrix} 1 & 0 \\ -2 & 5 \end{pmatrix}$.
\begin{enumerate}
\item Calculer $AB$ et $BA$. Que constate-t-on ?
\item Calculer $A^2 = A \times A$.
\item Vérifier que $A \times A^{-1} = I_2$ où $I_2 = \begin{pmatrix} 1 & 0 \\ 0 & 1 \end{pmatrix}$.
\end{enumerate}
\end{exercice}

\begin{exercice}[Condition d'inversibilité avec paramètre]
Soit la matrice $M = \begin{pmatrix} m & 2 \\ 3 & m-1 \end{pmatrix}$ où $m$ est un nombre réel.
\begin{enumerate}
\item Calculer $\det(M)$ en fonction de $m$.
\item Déterminer les valeurs de $m$ pour lesquelles $M$ n'est pas inversible.
\item Pour $m = 4$, montrer que $M$ est inversible et calculer $M^{-1}$.
\end{enumerate}
\end{exercice}

\begin{exercice}[Matrice d'une application linéaire]
Le plan est muni d'une base $(\vec{\imath}, \vec{\jmath})$. Soit $f$ l'application linéaire définie par :
\[ f(\vec{\imath}) = 2\vec{\imath} + \vec{\jmath} \qquad \text{et} \qquad f(\vec{\jmath}) = -\vec{\imath} + 3\vec{\jmath}. \]
\begin{enumerate}
\item Écrire la matrice $M$ de $f$ dans la base $(\vec{\imath}, \vec{\jmath})$.
\item Soit $\vec{u}$ le vecteur de coordonnées $(3\,; -1)$. À l'aide d'un produit de matrices, calculer les coordonnées de $f(\vec{u})$.
\item Déterminer l'expression de $f(x\,; y)$ pour tout vecteur de coordonnées $(x\,; y)$.
\item $f$ est-elle un automorphisme du plan ? Justifier.
\end{enumerate}
\end{exercice}

\begin{exercice}[Opérations sur les applications linéaires]
Dans une base $(\vec{\imath}, \vec{\jmath})$ du plan, $f$ et $g$ sont deux applications linéaires de matrices respectives :
\[ A = \begin{pmatrix} 1 & 1 \\ 0 & 2 \end{pmatrix} \qquad \text{et} \qquad B = \begin{pmatrix} 2 & 0 \\ 1 & 1 \end{pmatrix}. \]
\begin{enumerate}
\item Déterminer la matrice de $f + g$ et celle de $2f$.
\item Déterminer la matrice de $g \circ f$, puis celle de $f \circ g$. Comparer.
\item Montrer que $f$ est un automorphisme et déterminer la matrice de $f^{-1}$.
\end{enumerate}
\end{exercice}

\courssub{Je me prépare au Probatoire}

\begin{exercice}[Résolution d'un système par la méthode matricielle]
On considère le système $(S)$ d'inconnues $(x\,; y) \in \mathbb{R}^2$ :
\[ (S) : \begin{cases} 2x + 3y = 12 \\ x - y = 1 \end{cases} \]
\begin{enumerate}
\item Écrire le système $(S)$ sous la forme matricielle $AX = B$, où $X = \begin{pmatrix} x \\ y \end{pmatrix}$, en précisant les matrices $A$ et $B$.
\item Calculer $\det(A)$ et en déduire que le système $(S)$ admet une unique solution.
\item Calculer $A^{-1}$.
\item En déduire la résolution de $(S)$ par le calcul de $X = A^{-1}B$.
\item Vérifier la solution obtenue en la substituant dans les deux équations du système.
\end{enumerate}
\end{exercice}

\begin{exercice}[Les mélanges d'arachides du commerçant de Douala]
Un commerçant du marché Mboppi à Douala vend deux mélanges d'arachides. Le mélange A contient des arachides grillées et des arachides bouillies ; le prix d'un kilogramme d'arachides grillées est $x$ FCFA et celui d'un kilogramme d'arachides bouillies est $y$ FCFA.

Lundi, il a vendu $4$ kg d'arachides grillées et $2$ kg d'arachides bouillies pour une recette totale de $\num{5200}$ FCFA. Mardi, il a vendu $3$ kg d'arachides grillées et $5$ kg d'arachides bouillies pour une recette totale de $\num{7400}$ FCFA.
\begin{enumerate}
\item Traduire les recettes des deux jours par un système $(S)$ de deux équations d'inconnues $x$ et $y$.
\item Écrire $(S)$ sous forme matricielle $AX = B$ avec $X = \begin{pmatrix} x \\ y \end{pmatrix}$.
\item Calculer $\det(A)$. Que peut-on en conclure sur l'existence et l'unicité des prix ?
\item Calculer $A^{-1}$, puis résoudre le système et donner le prix d'un kilogramme de chaque type d'arachides.
\item Le commerçant souhaite préparer un panier contenant $2$ kg d'arachides grillées et $3$ kg d'arachides bouillies. Quel sera le prix de ce panier ?
\end{enumerate}
\end{exercice}

\begin{exercice}[Étude complète d'applications linéaires du plan]
Le plan vectoriel est muni d'une base $(\vec{\imath}, \vec{\jmath})$. On considère les applications linéaires $f$ et $g$ définies par :
\[ f(x\,; y) = (x + 2y\,; \; 3x + 6y) \qquad \text{et} \qquad g(x\,; y) = (2x - y\,; \; x + y). \]

\textbf{Partie A : étude de $f$}
\begin{enumerate}
\item Écrire la matrice $A$ de $f$ dans la base $(\vec{\imath}, \vec{\jmath})$.
\item Calculer $\det(A)$. L'application $f$ est-elle un automorphisme du plan ? Justifier.
\item Déterminer tous les vecteurs $(x\,; y)$ tels que $f(x\,; y) = (0\,; 0)$. Interpréter géométriquement l'ensemble obtenu.
\item Montrer que l'image de tout vecteur du plan par $f$ est un vecteur de la forme $(t\,; 3t)$ avec $t \in \mathbb{R}$. Que représente géométriquement l'ensemble des images ?
\end{enumerate}

\textbf{Partie B : étude de $g$}
\begin{enumerate}
\item Écrire la matrice $B$ de $g$ dans la base $(\vec{\imath}, \vec{\jmath})$.
\item Montrer que $g$ est un automorphisme et déterminer la matrice de $g^{-1}$.
\item En déduire les coordonnées du vecteur $\vec{u}$ tel que $g(\vec{u}) = (5\,; 4)$.
\end{enumerate}

\textbf{Partie C : composition}
\begin{enumerate}
\item Déterminer la matrice de $g \circ f$ dans la base $(\vec{\imath}, \vec{\jmath})$.
\item $g \circ f$ est-elle un automorphisme ? Justifier sans calcul supplémentaire de déterminant, en utilisant les résultats des parties A et B.
\end{enumerate}
\end{exercice}

\newpage

\coursec{Corrigés des exercices}

\begin{corrige}[Exercice 1 — QCM : une seule réponse exacte]
\textbf{1. Réponse b) : $11$.} En effet, $\det(A) = 2\times 4 - (-1)\times 3 = 8 + 3 = 11$.

\textbf{2. Réponse b) : non inversible.} On calcule $\det(M) = 1\times 6 - 2\times 3 = 6 - 6 = 0$. Une matrice carrée d'ordre $2$ est inversible si et seulement si son déterminant est non nul ; or ici $\det(M)=0$, donc $M$ n'est pas inversible. (On remarque d'ailleurs que la deuxième ligne est le triple de la première.)

\textbf{3. Réponse b) : $AB \neq BA$.} Le produit matriciel n'est \emph{pas commutatif} en général. Contre-exemple : avec $A = \begin{pmatrix} 1 & 1 \\ 0 & 1 \end{pmatrix}$ et $B = \begin{pmatrix} 1 & 0 \\ 1 & 1 \end{pmatrix}$, on obtient $AB = \begin{pmatrix} 2 & 1 \\ 1 & 1 \end{pmatrix}$ et $BA = \begin{pmatrix} 1 & 1 \\ 1 & 2 \end{pmatrix}$, qui sont différentes.

\textbf{4. Réponse c) : $A^{-1}$.} C'est un résultat fondamental du cours : la matrice de l'automorphisme réciproque $f^{-1}$ dans une base donnée est l'inverse de la matrice de $f$ dans cette même base.
\end{corrige}

\begin{corrige}[Exercice 2 — Vrai ou faux, en justifiant]
\textbf{1. Vrai.} La somme de deux matrices de même format est définie terme à terme et conserve le format : la somme de deux matrices $2\times 2$ est bien une matrice $2\times 2$.

\textbf{2. Faux.} Pour une matrice carrée d'ordre $2$, $A = \begin{pmatrix} a & b \\ c & d \end{pmatrix}$, on a $-A = \begin{pmatrix} -a & -b \\ -c & -d \end{pmatrix}$, donc
\[ \det(-A) = (-a)(-d) - (-b)(-c) = ad - bc = \det(A). \]
Ainsi $\det(-A) = \det(A)$ (et non $-\det(A)$), car multiplier les \emph{deux} lignes par $-1$ fait intervenir $(-1)^2 = 1$. Contre-exemple : $A = I_2$, $\det(-I_2) = 1 = \det(I_2) \neq -1$.

\textbf{3. Faux.} C'est exactement le contraire du théorème du cours : une matrice carrée d'ordre $2$ est inversible \textbf{si et seulement si} son déterminant est \emph{non nul}. Si $\det(A) = 0$, alors $A$ n'admet pas d'inverse. Exemple : $\begin{pmatrix} 1 & 2 \\ 2 & 4 \end{pmatrix}$ a un déterminant nul et n'est pas inversible.

\textbf{4. Vrai.} La matrice identité $I$ est l'élément neutre du produit matriciel : pour toute matrice carrée $A$ d'ordre $2$, $AI = IA = A$. Vérification directe :
\[ AI = \begin{pmatrix} 1 & 1 \\ 0 & 1 \end{pmatrix}\begin{pmatrix} 1 & 0 \\ 0 & 1 \end{pmatrix} = \begin{pmatrix} 1\times 1 + 1\times 0 & 1\times 0 + 1\times 1 \\ 0\times 1 + 1\times 0 & 0\times 0 + 1\times 1 \end{pmatrix} = \begin{pmatrix} 1 & 1 \\ 0 & 1 \end{pmatrix} = A, \]
\[ IA = \begin{pmatrix} 1 & 0 \\ 0 & 1 \end{pmatrix}\begin{pmatrix} 1 & 1 \\ 0 & 1 \end{pmatrix} = \begin{pmatrix} 1\times 1 + 0\times 0 & 1\times 1 + 0\times 1 \\ 0\times 1 + 1\times 0 & 0\times 1 + 1\times 1 \end{pmatrix} = \begin{pmatrix} 1 & 1 \\ 0 & 1 \end{pmatrix} = A. \]
\end{corrige}

\begin{corrige}[Exercice 3 — Définitions et complétion]
\textbf{1.} Le \cle{déterminant} de la matrice $A = \begin{pmatrix} a & b \\ c & d \end{pmatrix}$ est le nombre réel
\[ \det(A) = ad - bc, \qquad \text{noté aussi } \begin{vmatrix} a & b \\ c & d \end{vmatrix}. \]
C'est la différence entre le produit des termes de la diagonale principale et le produit des termes de l'autre diagonale.

\textbf{2.} Une matrice carrée $A$ d'ordre $2$ est inversible si et seulement si $\det(A) \neq 0$ ; dans ce cas :
\[ A^{-1} = \frac{1}{\det(A)}\begin{pmatrix} d & -b \\ -c & a \end{pmatrix}. \]
(Moyen mnémotechnique : on échange $a$ et $d$, on change les signes de $b$ et $c$, et on divise par le déterminant.)

\textbf{3.} La matrice d'une application linéaire $f$ dans une base $(\vec{\imath}, \vec{\jmath})$ a pour \textbf{colonnes} les coordonnées des vecteurs $f(\vec{\imath})$ et $f(\vec{\jmath})$ exprimés dans cette base : la première colonne est formée des coordonnées de $f(\vec{\imath})$, la deuxième de celles de $f(\vec{\jmath})$.

\textbf{4.} Un \cle{automorphisme} d'un plan vectoriel est une application linéaire \textbf{bijective} du plan dans lui-même (un endomorphisme bijectif). De façon équivalente, c'est une application linéaire dont la matrice dans une base quelconque est inversible, c'est-à-dire de déterminant non nul.
\end{corrige}

\begin{corrige}[Exercice 4 — Calculs directs]
\textbf{1.} La somme se calcule terme à terme :
\[ A + B = \begin{pmatrix} 2+1 & -1+0 \\ 3+(-2) & 4+5 \end{pmatrix} = \begin{pmatrix} 3 & -1 \\ 1 & 9 \end{pmatrix}. \]
Le produit par le réel $3$ multiplie chaque coefficient :
\[ 3A = \begin{pmatrix} 3\times 2 & 3\times(-1) \\ 3\times 3 & 3\times 4 \end{pmatrix} = \begin{pmatrix} 6 & -3 \\ 9 & 12 \end{pmatrix}. \]

\textbf{2.} $\det(A) = 2\times 4 - (-1)\times 3 = 8 + 3 = 11$ et $\det(B) = 1\times 5 - 0\times(-2) = 5$.

\textbf{3.} Comme $\det(A) = 11 \neq 0$, la matrice $A$ est inversible, et :
\[ A^{-1} = \frac{1}{11}\begin{pmatrix} 4 & -(-1) \\ -3 & 2 \end{pmatrix} = \frac{1}{11}\begin{pmatrix} 4 & 1 \\ -3 & 2 \end{pmatrix} = \begin{pmatrix} \dfrac{4}{11} & \dfrac{1}{11} \\[6pt] -\dfrac{3}{11} & \dfrac{2}{11} \end{pmatrix}. \]
\end{corrige}

\begin{corrige}[Exercice 5 — Produits de matrices]
\textbf{1.} Calcul de $AB$ (ligne par colonne) :
\begin{align*}
AB &= \begin{pmatrix} 2 & -1 \\ 3 & 4 \end{pmatrix}\begin{pmatrix} 1 & 0 \\ -2 & 5 \end{pmatrix} \\
&= \begin{pmatrix} 2\times 1 + (-1)\times(-2) & 2\times 0 + (-1)\times 5 \\ 3\times 1 + 4\times(-2) & 3\times 0 + 4\times 5 \end{pmatrix}
= \begin{pmatrix} 4 & -5 \\ -5 & 20 \end{pmatrix}.
\end{align*}
Calcul de $BA$ :
\begin{align*}
BA &= \begin{pmatrix} 1 & 0 \\ -2 & 5 \end{pmatrix}\begin{pmatrix} 2 & -1 \\ 3 & 4 \end{pmatrix} \\
&= \begin{pmatrix} 1\times 2 + 0\times 3 & 1\times(-1) + 0\times 4 \\ (-2)\times 2 + 5\times 3 & (-2)\times(-1) + 5\times 4 \end{pmatrix}
= \begin{pmatrix} 2 & -1 \\ 11 & 22 \end{pmatrix}.
\end{align*}
\textbf{Constat :} $AB \neq BA$ : le produit de deux matrices n'est pas commutatif en général.

\textbf{2.} Calcul de $A^2$ :
\begin{align*}
A^2 &= \begin{pmatrix} 2 & -1 \\ 3 & 4 \end{pmatrix}\begin{pmatrix} 2 & -1 \\ 3 & 4 \end{pmatrix} \\
&= \begin{pmatrix} 4 - 3 & -2 - 4 \\ 6 + 12 & -3 + 16 \end{pmatrix}
= \begin{pmatrix} 1 & -6 \\ 18 & 13 \end{pmatrix}.
\end{align*}

\textbf{3.} D'après l'exercice 4, $A^{-1} = \dfrac{1}{11}\begin{pmatrix} 4 & 1 \\ -3 & 2 \end{pmatrix}$. Alors :
\begin{align*}
A \times A^{-1} &= \frac{1}{11}\begin{pmatrix} 2 & -1 \\ 3 & 4 \end{pmatrix}\begin{pmatrix} 4 & 1 \\ -3 & 2 \end{pmatrix} \\
&= \frac{1}{11}\begin{pmatrix} 8 + 3 & 2 - 2 \\ 12 - 12 & 3 + 8 \end{pmatrix}
= \frac{1}{11}\begin{pmatrix} 11 & 0 \\ 0 & 11 \end{pmatrix}
= \begin{pmatrix} 1 & 0 \\ 0 & 1 \end{pmatrix} = I_2.
\end{align*}
La vérification est conforme à la définition de l'inverse : $A A^{-1} = I_2$.
\end{corrige}

\begin{corrige}[Exercice 6 — Condition d'inversibilité avec paramètre]
\textbf{1.} $\det(M) = m(m-1) - 2\times 3 = m^2 - m - 6$.

\textbf{2.} $M$ n'est pas inversible si et seulement si $\det(M) = 0$, c'est-à-dire :
\[ m^2 - m - 6 = 0. \]
Le discriminant est $\Delta = (-1)^2 - 4\times 1\times(-6) = 1 + 24 = 25 = 5^2$. Les racines sont :
\[ m_1 = \frac{1 - 5}{2} = -2 \qquad \text{et} \qquad m_2 = \frac{1 + 5}{2} = 3. \]
(On peut vérifier la factorisation : $m^2 - m - 6 = (m-3)(m+2)$.)
\textbf{Conclusion :} $M$ n'est pas inversible pour $m = -2$ et $m = 3$ ; elle est inversible pour toute autre valeur réelle de $m$.

\textbf{3.} Pour $m = 4$ : $M = \begin{pmatrix} 4 & 2 \\ 3 & 3 \end{pmatrix}$ et $\det(M) = 4^2 - 4 - 6 = 16 - 10 = 6 \neq 0$ (on retrouve bien que $4 \notin \{-2\,; 3\}$). Donc $M$ est inversible et :
\[ M^{-1} = \frac{1}{6}\begin{pmatrix} 3 & -2 \\ -3 & 4 \end{pmatrix} = \begin{pmatrix} \dfrac{1}{2} & -\dfrac{1}{3} \\[6pt] -\dfrac{1}{2} & \dfrac{2}{3} \end{pmatrix}. \]
Vérification rapide : $M M^{-1} = \dfrac{1}{6}\begin{pmatrix} 4 & 2 \\ 3 & 3 \end{pmatrix}\begin{pmatrix} 3 & -2 \\ -3 & 4 \end{pmatrix} = \dfrac{1}{6}\begin{pmatrix} 6 & 0 \\ 0 & 6 \end{pmatrix} = I_2$. \checkmark
\end{corrige}

\begin{corrige}[Exercice 7 — Matrice d'une application linéaire]
\textbf{1.} Les colonnes de $M$ sont les coordonnées de $f(\vec{\imath})$ et $f(\vec{\jmath})$ dans la base $(\vec{\imath}, \vec{\jmath})$ :
\[ f(\vec{\imath}) = 2\vec{\imath} + \vec{\jmath} \;\longrightarrow\; \begin{pmatrix} 2 \\ 1 \end{pmatrix}, \qquad f(\vec{\jmath}) = -\vec{\imath} + 3\vec{\jmath} \;\longrightarrow\; \begin{pmatrix} -1 \\ 3 \end{pmatrix}. \]
D'où $M = \begin{pmatrix} 2 & -1 \\ 1 & 3 \end{pmatrix}$.

\textbf{2.} Les coordonnées de $f(\vec{u})$ s'obtiennent par le produit $M \begin{pmatrix} 3 \\ -1 \end{pmatrix}$ :
\[ \begin{pmatrix} 2 & -1 \\ 1 & 3 \end{pmatrix}\begin{pmatrix} 3 \\ -1 \end{pmatrix} = \begin{pmatrix} 2\times 3 + (-1)\times(-1) \\ 1\times 3 + 3\times(-1) \end{pmatrix} = \begin{pmatrix} 7 \\ 0 \end{pmatrix}. \]
Donc $f(\vec{u})$ a pour coordonnées $(7\,; 0)$.

\textbf{3.} Pour tout vecteur de coordonnées $(x\,; y)$ :
\[ \begin{pmatrix} 2 & -1 \\ 1 & 3 \end{pmatrix}\begin{pmatrix} x \\ y \end{pmatrix} = \begin{pmatrix} 2x - y \\ x + 3y \end{pmatrix}, \]
donc $f(x\,; y) = (2x - y\,; \; x + 3y)$.

\textbf{4.} $\det(M) = 2\times 3 - (-1)\times 1 = 6 + 1 = 7 \neq 0$. La matrice $M$ est inversible, donc $f$ est bijective : $f$ est bien un \cle{automorphisme} du plan.
\end{corrige}

\begin{corrige}[Exercice 8 — Opérations sur les applications linéaires]
\textbf{1.} La matrice de $f + g$ est $A + B$, et celle de $2f$ est $2A$ :
\[ A + B = \begin{pmatrix} 1+2 & 1+0 \\ 0+1 & 2+1 \end{pmatrix} = \begin{pmatrix} 3 & 1 \\ 1 & 3 \end{pmatrix}, \qquad 2A = \begin{pmatrix} 2 & 2 \\ 0 & 4 \end{pmatrix}. \]

\textbf{2.} La matrice de la composée est le produit des matrices, \emph{dans le même ordre} : la matrice de $g \circ f$ est $BA$, celle de $f \circ g$ est $AB$.
\[ BA = \begin{pmatrix} 2 & 0 \\ 1 & 1 \end{pmatrix}\begin{pmatrix} 1 & 1 \\ 0 & 2 \end{pmatrix} = \begin{pmatrix} 2+0 & 2+0 \\ 1+0 & 1+2 \end{pmatrix} = \begin{pmatrix} 2 & 2 \\ 1 & 3 \end{pmatrix}. \]
\[ AB = \begin{pmatrix} 1 & 1 \\ 0 & 2 \end{pmatrix}\begin{pmatrix} 2 & 0 \\ 1 & 1 \end{pmatrix} = \begin{pmatrix} 2+1 & 0+1 \\ 0+2 & 0+2 \end{pmatrix} = \begin{pmatrix} 3 & 1 \\ 2 & 2 \end{pmatrix}. \]
\textbf{Comparaison :} $BA \neq AB$, donc $g \circ f \neq f \circ g$ : la composition des applications linéaires n'est pas commutative en général.

\textbf{3.} $\det(A) = 1\times 2 - 1\times 0 = 2 \neq 0$, donc $A$ est inversible et $f$ est un automorphisme du plan. La matrice de $f^{-1}$ est :
\[ A^{-1} = \frac{1}{2}\begin{pmatrix} 2 & -1 \\ 0 & 1 \end{pmatrix} = \begin{pmatrix} 1 & -\dfrac{1}{2} \\[6pt] 0 & \dfrac{1}{2} \end{pmatrix}. \]
Vérification : $A A^{-1} = \dfrac{1}{2}\begin{pmatrix} 1 & 1 \\ 0 & 2 \end{pmatrix}\begin{pmatrix} 2 & -1 \\ 0 & 1 \end{pmatrix} = \dfrac{1}{2}\begin{pmatrix} 2 & 0 \\ 0 & 2 \end{pmatrix} = I_2$. \checkmark
\end{corrige}

\begin{corrige}[Exercice 9 — Résolution d'un système par la méthode matricielle]
\textbf{1.} Le système $(S) : \begin{cases} 2x + 3y = 12 \\ x - y = 1 \end{cases}$ s'écrit $AX = B$ avec :
\[ A = \begin{pmatrix} 2 & 3 \\ 1 & -1 \end{pmatrix}, \qquad X = \begin{pmatrix} x \\ y \end{pmatrix}, \qquad B = \begin{pmatrix} 12 \\ 1 \end{pmatrix}. \]

\textbf{2.} $\det(A) = 2\times(-1) - 3\times 1 = -2 - 3 = -5 \neq 0$. La matrice $A$ est donc inversible, ce qui prouve que le système $(S)$ admet une \textbf{unique solution}, donnée par $X = A^{-1}B$.

\textbf{3.} $A^{-1} = \dfrac{1}{-5}\begin{pmatrix} -1 & -3 \\ -1 & 2 \end{pmatrix} = \dfrac{1}{5}\begin{pmatrix} 1 & 3 \\ 1 & -2 \end{pmatrix}$.

\textbf{4.} On calcule :
\[ X = A^{-1}B = \frac{1}{5}\begin{pmatrix} 1 & 3 \\ 1 & -2 \end{pmatrix}\begin{pmatrix} 12 \\ 1 \end{pmatrix} = \frac{1}{5}\begin{pmatrix} 12 + 3 \\ 12 - 2 \end{pmatrix} = \frac{1}{5}\begin{pmatrix} 15 \\ 10 \end{pmatrix} = \begin{pmatrix} 3 \\ 2 \end{pmatrix}. \]
L'unique solution de $(S)$ est donc $(x\,; y) = (3\,; 2)$.

\textbf{5.} Vérification : $2\times 3 + 3\times 2 = 6 + 6 = 12$ \checkmark\; et $3 - 2 = 1$ \checkmark. La solution convient.

\medskip
\textbf{Barème indicatif (sur 5 pts) :} écriture matricielle correcte : 1 pt ; calcul du déterminant et conclusion d'unicité : 1 pt ; calcul de $A^{-1}$ : 1 pt ; calcul de $X = A^{-1}B$ : 1 pt ; vérification : 1 pt.

\textbf{Points de vigilance :} citer explicitement la condition $\det(A) \neq 0$ avant de conclure à l'unicité ; attention à l'ordre dans le produit $A^{-1}B$ (et non $BA^{-1}$, qui n'aurait d'ailleurs pas de sens ici) ; ne pas oublier le facteur $\dfrac{1}{\det(A)}$ dans le calcul de l'inverse.
\end{corrige}

\begin{corrige}[Exercice 10 — Les mélanges d'arachides du commerçant de Douala]
\textbf{1.} La recette de lundi donne $4x + 2y = \num{5200}$ et celle de mardi $3x + 5y = \num{7400}$. D'où le système :
\[ (S) : \begin{cases} 4x + 2y = \num{5200} \\ 3x + 5y = \num{7400} \end{cases} \]

\textbf{2.} Forme matricielle $AX = B$ :
\[ A = \begin{pmatrix} 4 & 2 \\ 3 & 5 \end{pmatrix}, \qquad X = \begin{pmatrix} x \\ y \end{pmatrix}, \qquad B = \begin{pmatrix} \num{5200} \\ \num{7400} \end{pmatrix}. \]

\textbf{3.} $\det(A) = 4\times 5 - 2\times 3 = 20 - 6 = 14 \neq 0$. La matrice $A$ est inversible : le système admet une \textbf{unique solution}. Il existe donc un unique couple de prix $(x\,; y)$ compatible avec les recettes des deux jours.

\textbf{4.} Calcul de l'inverse :
\[ A^{-1} = \frac{1}{14}\begin{pmatrix} 5 & -2 \\ -3 & 4 \end{pmatrix}. \]
Puis :
\begin{align*}
X = A^{-1}B &= \frac{1}{14}\begin{pmatrix} 5 & -2 \\ -3 & 4 \end{pmatrix}\begin{pmatrix} \num{5200} \\ \num{7400} \end{pmatrix} \\
&= \frac{1}{14}\begin{pmatrix} 5\times \num{5200} - 2\times \num{7400} \\ -3\times \num{5200} + 4\times \num{7400} \end{pmatrix}
= \frac{1}{14}\begin{pmatrix} \num{26000} - \num{14800} \\ -\num{15600} + \num{29600} \end{pmatrix} \\
&= \frac{1}{14}\begin{pmatrix} \num{11200} \\ \num{14000} \end{pmatrix}
= \begin{pmatrix} 800 \\ \num{1000} \end{pmatrix}.
\end{align*}
Le kilogramme d'arachides grillées coûte $x = 800$ FCFA et le kilogramme d'arachides bouillies coûte $y = \num{1000}$ FCFA.

Vérification : $4\times 800 + 2\times \num{1000} = \num{3200} + \num{2000} = \num{5200}$ \checkmark\; et $3\times 800 + 5\times \num{1000} = \num{2400} + \num{5000} = \num{7400}$ \checkmark.

\textbf{5.} Le prix du panier est :
\[ 2x + 3y = 2\times 800 + 3\times \num{1000} = \num{1600} + \num{3000} = \num{4600} \text{ FCFA}. \]

\medskip
\textbf{Barème indicatif (sur 6 pts) :} mise en équation : 1 pt ; écriture matricielle : 0,5 pt ; déterminant et conclusion : 1 pt ; inverse : 1 pt ; résolution et interprétation des prix : 1,5 pt ; prix du panier : 1 pt.

\textbf{Points de vigilance :} bien distinguer les inconnues (prix unitaires) des coefficients (quantités vendues) ; conclure en langage courant (prix en FCFA) et pas seulement par un couple de nombres ; la vérification par substitution est un réflexe valorisé au Probatoire.
\end{corrige}

\begin{corrige}[Exercice 11 — Étude complète d'applications linéaires du plan]
\textbf{Partie A : étude de $f$}

\textbf{1.} On a $f(x\,; y) = (x + 2y\,; \; 3x + 6y)$, donc $f(\vec{\imath}) = f(1\,;0) = (1\,; 3)$ et $f(\vec{\jmath}) = f(0\,;1) = (2\,; 6)$. Les colonnes de $A$ sont ces coordonnées :
\[ A = \begin{pmatrix} 1 & 2 \\ 3 & 6 \end{pmatrix}. \]

\textbf{2.} $\det(A) = 1\times 6 - 2\times 3 = 6 - 6 = 0$. Le déterminant est nul, donc la matrice $A$ n'est pas inversible : $f$ \textbf{n'est pas un automorphisme} du plan (elle n'est pas bijective).

\textbf{3.} On cherche les vecteurs $(x\,; y)$ tels que $f(x\,; y) = (0\,; 0)$, c'est-à-dire :
\[ \begin{cases} x + 2y = 0 \\ 3x + 6y = 0 \end{cases} \]
La deuxième équation est le triple de la première : le système se réduit à $x + 2y = 0$, soit $x = -2y$. En posant $y = t$ ($t \in \mathbb{R}$), les solutions sont les vecteurs $(-2t\,; t)$, $t \in \mathbb{R}$.

\textbf{Interprétation géométrique :} l'ensemble des vecteurs d'image nulle (le \emph{noyau} de $f$) est la \textbf{droite vectorielle} dirigée par le vecteur $(-2\,; 1)$. Il y a une infinité de vecteurs ayant la même image nulle, ce qui confirme que $f$ n'est pas injective.

\textbf{4.} Soit $(x\,; y)$ un vecteur quelconque. Posons $t = x + 2y$. Alors :
\[ f(x\,; y) = (x + 2y\,; \; 3x + 6y) = (t\,; \; 3(x + 2y)) = (t\,; 3t). \]
Réciproquement, tout vecteur $(t\,; 3t)$ est atteint (par exemple $f(t\,; 0) = (t\,; 3t)$). L'ensemble des images de $f$ est donc la \textbf{droite vectorielle} dirigée par le vecteur $(1\,; 3)$ : $f$ « écrase » tout le plan sur une droite, ce qui confirme qu'elle n'est pas surjective.

\textbf{Partie B : étude de $g$}

\textbf{1.} $g(x\,; y) = (2x - y\,; \; x + y)$ donne $g(\vec{\imath}) = (2\,; 1)$ et $g(\vec{\jmath}) = (-1\,; 1)$, d'où :
\[ B = \begin{pmatrix} 2 & -1 \\ 1 & 1 \end{pmatrix}. \]

\textbf{2.} $\det(B) = 2\times 1 - (-1)\times 1 = 2 + 1 = 3 \neq 0$, donc $B$ est inversible et $g$ est un \textbf{automorphisme} du plan. La matrice de $g^{-1}$ est :
\[ B^{-1} = \frac{1}{3}\begin{pmatrix} 1 & 1 \\ -1 & 2 \end{pmatrix}. \]

\textbf{3.} Le vecteur $\vec{u}$ tel que $g(\vec{u}) = (5\,; 4)$ est $\vec{u} = g^{-1}(5\,; 4)$, dont les coordonnées sont :
\[ B^{-1}\begin{pmatrix} 5 \\ 4 \end{pmatrix} = \frac{1}{3}\begin{pmatrix} 1 & 1 \\ -1 & 2 \end{pmatrix}\begin{pmatrix} 5 \\ 4 \end{pmatrix} = \frac{1}{3}\begin{pmatrix} 5 + 4 \\ -5 + 8 \end{pmatrix} = \frac{1}{3}\begin{pmatrix} 9 \\ 3 \end{pmatrix} = \begin{pmatrix} 3 \\ 1 \end{pmatrix}. \]
Donc $\vec{u}$ a pour coordonnées $(3\,; 1)$. Vérification : $g(3\,; 1) = (2\times 3 - 1\,; \; 3 + 1) = (5\,; 4)$ \checkmark.

\textbf{Partie C : composition}

\textbf{1.} La matrice de $g \circ f$ est le produit $BA$ :
\[ BA = \begin{pmatrix} 2 & -1 \\ 1 & 1 \end{pmatrix}\begin{pmatrix} 1 & 2 \\ 3 & 6 \end{pmatrix} = \begin{pmatrix} 2 - 3 & 4 - 6 \\ 1 + 3 & 2 + 6 \end{pmatrix} = \begin{pmatrix} -1 & -2 \\ 4 & 8 \end{pmatrix}. \]

\textbf{2.} $g \circ f$ \textbf{n'est pas un automorphisme}. Justification sans calcul de déterminant : pour tout vecteur $\vec{v}$ du plan, $(g \circ f)(\vec{v}) = g(f(\vec{v}))$ appartient à l'image de $g$ restreinte à l'image de $f$ ; or d'après la partie A, l'image de $f$ est une droite vectorielle (et non le plan tout entier). L'image de $g \circ f$ est donc contenue dans l'image par $g$ de cette droite, qui est au plus une droite : $g \circ f$ ne peut pas être surjective, donc n'est pas bijective. (De façon cohérente, on observerait $\det(BA) = -8 + 8 = 0$ : composer une application non bijective avec une bijection ne peut pas « réparer » la non-bijectivité.)

\medskip
\textbf{Barème indicatif (sur 10 pts) :} Partie A : matrice 0,5 pt ; déterminant et conclusion 1 pt ; noyau et interprétation 1,5 pt ; image et interprétation 1 pt. Partie B : matrice 0,5 pt ; automorphisme et inverse 1,5 pt ; vecteur $\vec{u}$ 1 pt. Partie C : matrice de $g \circ f$ 1,5 pt ; justification 1,5 pt.

\textbf{Points de vigilance :} la bijectivité d'une application linéaire du plan se justifie \emph{uniquement} par $\det \neq 0$ (citer le théorème) ; pour la question C.2, le barème récompense le raisonnement sur les images (ou sur la non-injectivité héritée de $f$) et non un simple calcul de déterminant, puisque l'énoncé l'interdit ; dans la question A.4, il faut montrer l'inclusion \emph{et} préciser que tout $t$ est atteint pour conclure que l'image est exactement la droite.
\end{corrige}

\newpage

\coursec{Fiche bilan}

\coursec{Matrices et applications linéaires d'un plan vectoriel}

\begin{popfigure}
\begin{tikzpicture}[mindmap, grow cyclic, every node/.style={concept, font=\small},
  concept color=popblue!80,
  level 1/.append style={level distance=3.4cm, sibling angle=60, font=\footnotesize},
  level 2/.append style={level distance=2.1cm, sibling angle=45, font=\scriptsize},
  scale=0.82, transform shape]
\node[concept, concept color=popblue!85, text=white, font=\bfseries\small] {Matrices et applications lin\'eaires}
  child[concept color=poporange!85] { node {Pr\'esentation}
    child { node {Type $(n,p)$ : $n$ lignes, $p$ colonnes} }
    child { node {Coefficients $a_{ij}$} }
  }
  child[concept color=popgreen!75] { node {Op\'erations}
    child { node {Somme $A+B$ (m\^eme type)} }
    child { node {Produit $\lambda A$ (scalaire)} }
    child { node {Produit $A\times B$ : $A_{n\times p}B_{p\times q}$} }
  }
  child[concept color=popgold!90] { node {D\'eterminant}
    child { node {$\det A = ad-bc$} }
    child { node {$\det A \neq 0 \Leftrightarrow A$ inversible} }
  }
  child[concept color=poppink!80] { node {Inverse}
    child { node[align=center] {$A^{-1}=\frac{1}{\det A}\begin{pmatrix} d & -b \\ -c & a \end{pmatrix}$} }
    child { node {Syst\`emes $AX=B \Rightarrow X=A^{-1}B$} }
  }
  child[concept color=poppurple!75] { node {Application lin\'eaire}
    child { node {$f(u+v)=f(u)+f(v)$} }
    child { node {$f(\lambda u)=\lambda f(u)$} }
    child { node {Matrice : images de la base en colonnes} }
  }
  child[concept color=popblue!60] { node {Liens matrice--application}
    child { node {$\mathrm{Mat}(f+g)$, $\mathrm{Mat}(\lambda f)$} }
    child { node {$\mathrm{Mat}(g\circ f)=\mathrm{Mat}(g)\times\mathrm{Mat}(f)$} }
    child { node {$\mathrm{Mat}(f^{-1})=M^{-1}$} }
  };
\end{tikzpicture}
\legende{Carte mentale de la le\c{c}on : les six id\'ees ma\^itresses des matrices et des applications lin\'eaires du plan.}
\end{popfigure}

\begin{aretenir}[L'essentiel de la le\c{c}on]
\textbf{D\'efinitions cl\'es.}
\begin{itemize}
  \item Une \cle{matrice} de type $(n,p)$ est un tableau de nombres r\'eels \`a $n$ lignes et $p$ colonnes ; on note $A=(a_{ij})$ o\`u $a_{ij}$ est le coefficient de la ligne $i$ et de la colonne $j$.
  \item Une matrice est \cle{carr\'ee d'ordre 2} si $n=p=2$ : $A=\begin{pmatrix} a & b \\ c & d \end{pmatrix}$.
  \item Une application $f$ du plan vectoriel dans lui-m\^eme est \cle{lin\'eaire} si pour tous vecteurs $u,v$ et tout r\'eel $\lambda$ :
        $f(u+v)=f(u)+f(v)$ et $f(\lambda u)=\lambda f(u)$.
  \item La \cle{matrice de $f$} dans une base $\mathcal{B}=(\vec{u},\vec{v})$ a pour colonnes les coordonn\'ees de $f(\vec{u})$ et $f(\vec{v})$ dans $\mathcal{B}$.
\end{itemize}

\textbf{Formules \`a conna\^itre par c\oe ur.}
\begin{center}
\begin{tabularx}{\linewidth}{|l|X|}\hline
\rowcolor{popblueL} \textbf{Notion} & \textbf{Formule / R\`egle} \\ \hline
Somme & $A+B=(a_{ij}+b_{ij})$ \textbf{(m\^eme type exig\'e)} \\ \hline
Produit par un r\'eel & $\lambda A=(\lambda a_{ij})$ \\ \hline
Produit de matrices & $(AB)_{ij}=$ ligne $i$ de $A$ $\times$ colonne $j$ de $B$ ; \textbf{d\'efini si nb col($A$)=nb lig($B$)} ; en g\'en\'eral $AB\neq BA$ \\ \hline
D\'eterminant & $\det A = ad-bc$ pour $A=\begin{pmatrix} a & b \\ c & d \end{pmatrix}$ \\ \hline
Inversibilit\'e & $A$ inversible $\iff \det A \neq 0$ \\ \hline
Inverse & $A^{-1}=\dfrac{1}{ad-bc}\begin{pmatrix} d & -b \\ -c & a \end{pmatrix}$ et $AA^{-1}=A^{-1}A=I_2$ \\ \hline
Image d'un vecteur & Si $X=\begin{pmatrix} x \\ y \end{pmatrix}$ sont les coord. de $u$ dans $\mathcal{B}$, alors $f(u)$ a pour coord. $X'=MX$ \\ \hline
Compos\'ee & $\mathrm{Mat}(g\circ f)=\mathrm{Mat}(g)\times\mathrm{Mat}(f)$ \textbf{(ordre : on applique $f$ puis $g$)} \\ \hline
R\'eciproque & $f$ automorphisme $\iff \det M \neq 0$, et $\mathrm{Mat}(f^{-1})=M^{-1}$ \\ \hline
\end{tabularx}
\end{center}

\textbf{M\'ethodes express.}
\begin{itemize}
  \item \emph{Produit $A\times B$} : disposer $B$ au-dessus et \`a droite de $A$ ; chaque coefficient est une somme de produits ligne-colonne.
  \item \emph{Syst\`eme $\begin{cases} ax+by=e \\ cx+dy=f \end{cases}$} : \'ecrire $AX=B$ avec $X=\begin{pmatrix} x \\ y \end{pmatrix}$, $B=\begin{pmatrix} e \\ f \end{pmatrix}$ ; si $\det A\neq 0$, solution unique $X=A^{-1}B$.
  \item \emph{Matrice de $f$ dans $\mathcal{B}=(\vec{u},\vec{v})$} : calculer $f(\vec{u})$ et $f(\vec{v})$, \'ecrire leurs coordonn\'ees dans $\mathcal{B}$ \textbf{en colonnes}.
  \item \emph{Montrer qu'une application est lin\'eaire} : v\'erifier les deux identit\'es d\'efinitoires, ou exhiber sa matrice dans une base.
\end{itemize}
\end{aretenir}

\begin{aretenir}[Checklist avant le Probatoire]
\begin{itemize}
  \item[$\square$] Je sais lire le type $(n,p)$ d'une matrice et rep\'erer le coefficient $a_{ij}$.
  \item[$\square$] Je sais additionner deux matrices de m\^eme type et multiplier une matrice par un r\'eel.
  \item[$\square$] Je sais calculer un produit $A\times B$ par la r\`egle ligne $\times$ colonne (en v\'erifiant la condition de compatibilit\'e).
  \item[$\square$] Je me souviens que le produit matriciel n'est pas commutatif : $AB\neq BA$ en g\'en\'eral.
  \item[$\square$] Je calcule sans h\'esiter $\det A=ad-bc$ et je conclus sur l'inversibilit\'e.
  \item[$\square$] Je connais la formule de $A^{-1}$ et je v\'erifie mon r\'esultat par $AA^{-1}=I_2$.
  \item[$\square$] Je sais traduire un syst\`eme de deux \'equations en \'equation matricielle $AX=B$ et le r\'esoudre par $X=A^{-1}B$.
  \item[$\square$] Je sais \'ecrire la matrice d'une application lin\'eaire dans une base quelconque (images des vecteurs de base en colonnes).
  \item[$\square$] Je calcule les coordonn\'ees de $f(u)$ par le produit $X'=MX$.
  \item[$\square$] Je respecte l'ordre dans $\mathrm{Mat}(g\circ f)=\mathrm{Mat}(g)\times\mathrm{Mat}(f)$.
  \item[$\square$] Je sais qu'un automorphisme a une matrice inversible et que $\mathrm{Mat}(f^{-1})=M^{-1}$.
  \item[$\square$] Je r\'edige proprement : conditions d'existence des produits, conclusion $\det A\neq 0$ avant d'inverser.
\end{itemize}
\end{aretenir}

\findecours

\end{document}
